% ===========================================================================
%  Capítulo 3 — As Árvores de Progressão
%  Prosa à mão; o catálogo das 18 árvores vem de gen/arvores.tex.
% ===========================================================================
\chapter{As Árvores de Progressão}

Não existem classes engessadas. O sistema funciona através de Árvores de Progressão,
divididas em três grandes pilares:

\begin{description}
    \item[Árvore da Magia] feitiços de ataque, suporte e invocação. Recurso: \rec{PM}.
    \item[Árvore do Corpo] os três Estilos Divinos de esgrima, mais armas pesadas, escudos,
          arquearia e a híbrida do vento. Recurso: \rec{PT}.
    \item[Árvore de Utilidade] os especialistas em mundo. Recurso: \rec{PP}.
\end{description}

Este capítulo cobre primeiro as regras \emph{compartilhadas} entre árvores do mesmo pilar,
e termina com o catálogo completo --- magias, técnicas, talentos e Maestrias das 18
sub-árvores, o mesmo dado que alimenta a ficha do site, então nunca diverge dela.

\begin{nota}{Escolas Formais e Ofícios}
    As oito escolas de magia e os três Estilos Divinos são \termo{Escolas Formais}: têm
    mestres vivos, sedes, hierarquia e títulos reconhecidos no mundo inteiro --- usam os
    nomes canônicos de rank (Principiante $\to$ Imperador) e conferem status social. As
    demais são \termo{Ofícios} aprendidos na estrada, sem diploma: mecanicamente idênticos
    (mesmo Bônus de Rank, mesmos custos de PA, mesma contagem de conhecimentos), mas os
    nomes dos patamares mudam.

    \medskip
    O sétimo degrau (Deus) existe só nas Escolas Formais, e mesmo lá é narrativo. Um Ofício
    termina no sexto patamar.
\end{nota}

\begin{nota}{Quem veste Touki}
    Toda sub-árvore da Árvore do Corpo desbloqueia o Touki no \textbf{terceiro} patamar ---
    inclusive Arquearia, inclusive Escudos. A única exceção é o Estilo Deus da Espada, que
    desperta a aura no \textbf{segundo} patamar (a doutrina inteira dele é velocidade). As
    árvores de Magia e de Utilidade nunca recebem Touki nem Pontos de Touki, por mais alto
    que seja o rank.
\end{nota}

% ---------------------------------------------------------------------------
\section{A Árvore do Corpo --- Sistemas Compartilhados}

Antes de qualquer estilo específico, três sistemas governam todos os guerreiros: o Dado de
Arma, o Touki e o Triângulo dos Estilos.

\begin{nota}{Espadachim ou Guerreiro}
    Apenas quem estudou uma das Três Grandes Escolas --- Deus da Espada, Deus da Água e Deus
    do Norte --- é chamado de \termo{Espadachim}. Todos os outros, mesmo empunhando espada,
    são apenas \termo{Guerreiros}.
\end{nota}

\subsection{O Dado de Arma e a Escalada de Maestria}

O dano de um guerreiro vem da arma, não do corpo. Conforme você sobe de Rank num estilo, o
Dado Base sobe degraus na Escada de Dados --- a tabela completa está no Capítulo 1, junto
das outras tabelas de referência.

\begin{formula}
    \textbf{Dano marcial $=$ Dado de Arma (escalado) $+$ Atributo $+$ Bônus do Rank}
\end{formula}

Qual Rank conta depende do tipo de ataque, e a regra tem exatamente dois casos --- nunca um
terceiro:

\begin{itemize}
    \item \textbf{Ataque comum} (golpe simples, sem nome, sem técnica): use o \textbf{maior}
          Rank que você tiver entre todas as suas árvores do Corpo. Um personagem com Norte
          Santo e Espada Principiante rola os degraus do Norte Santo em qualquer golpe
          comum, não importa com qual arma.
    \item \textbf{Técnica nomeada} (qualquer habilidade comprada de uma árvore específica):
          use sempre o Rank \textbf{daquela árvore que concedeu a técnica}, mesmo que seja
          menor que o seu maior Rank geral. O mesmo personagem usando a Espada de Luz rola
          só os degraus do Rank Principiante --- é por isso que ela sai fraca na mão dele,
          apesar do Norte Santo.
\end{itemize}

A mesma lógica dos dois casos vale pra \emph{qualquer} talento ou regra do livro que
mencione \enquote{seu Bônus de Rank} sem dizer de qual árvore.

\begin{nota}{Acima do topo da escada}
    O 5d12 é o último degrau. Se um talento ou Maestria te der um degrau além dele --- Espada
    Emprestada e Punho Duplo são os dois casos do livro --- cada degrau excedente vira
    \textbf{$+2$ de dano fixo} em vez de sumir. Nenhum PA gasto em degrau é jogado fora.
\end{nota}

\begin{aviso}{A progressão de degraus é padrão, não universal}
    A maioria das árvores do Corpo sobe exatamente $+1$ degrau por Rank (total $+6$ no
    Imperador), mas não é regra fixa --- cada árvore define os próprios degraus por patamar
    no catálogo dela. O Deus da Espada sobe mais rápido (nove degraus no total, de propósito
    --- é a identidade da árvore); Cavalaria e Escudos sobe mais devagar, porque o valor
    dela está em proteger o grupo, não em dano. Confira o catálogo da árvore específica
    antes de calcular o dado final de alguém.
\end{aviso}

\subsection{Touki (Aura de Batalha)}

O Touki é uma camada de mana que o guerreiro veste sobre o próprio corpo --- endurece a pele
como aço, reforça o fio da lâmina e amplifica força, velocidade e reflexos. Rank
Intermediário é o mais alto que alguém alcança sem Touki.

\begin{nota}{Pontos de Touki (PT) --- as duas reservas}
    \textbf{PT Menor} --- do 1º patamar em diante: PT iguais ao seu \textbf{Vigor} (mínimo
    1). Não é aura, é fôlego --- paga técnicas, e nada mais.

    \medskip
    \textbf{PT Pleno} --- a partir do 3º patamar (2º no Deus da Espada): a reserva passa a
    ser \textbf{Espírito $+$ Vigor}, e você desbloqueia o Manto de Touki e as manobras de
    gasto abaixo. \textbf{Crescimento:} $+1$ PT por patamar que já tenha PT Pleno em
    qualquer árvore do Corpo. Cavalaria e Escudos concede $+2$ por patamar em vez de $+1$,
    por gastar PT mais rápido que qualquer outra árvore.

    \begin{itemize}
        \item PT são recuperados \textbf{integralmente} em um Descanso Curto --- e são o
              único recurso que volta inteiro nele (PM voltam pela metade). O Cap. 4, §6
              limita a dois Descansos Curtos entre dois Longos.
        \item PT não podem ser convertidos em PM, nem PM em PT.
        \item Um personagem com Ranks em mais de um estilo marcial usa uma reserva única.
    \end{itemize}
\end{nota}

\textbf{O Manto de Touki} (passivo, gratuito, independente de PT) --- a partir do Rank
Avançado, enquanto consciente e não Exausto: $+$CA igual à metade do Bônus de Rank
(arredondado pra cima); Redução contra projéteis mundanos igual ao dobro do Bônus de Rank (e
nunca sofre crítico deles); ataques desarmados e com objetos improvisados contam como
mágicos. \textbf{O Manto não consome PT e continua ativo mesmo com a reserva em 0} --- ele é
vestido pelo Rank, não comprado com o recurso.

\begin{tabelalivro}{C{1.6cm} L{3.6cm} L{8.6cm}}{Manobras de Touki (custam PT)}
\cabtabela{\textbf{Custo} & \textbf{Manobra} & \textbf{Efeito}}
1 PT & Touki Concentrado & Sem Ação. Até o fim do turno, some seu Bônus de Rank ao dano de novo e reduza todo dano físico recebido pelo mesmo valor. \\
1 PT & Lâmina de Touki & Sem Ação. Por 1 minuto, sua arma corta pedra e aço, conta como mágica e ignora Resistência a cortante/perfurante. \\
2 PT & Golpe Estendido & 1 Ação. Clarão da lâmina que atinge um alvo a até 9m. Dano de arma normal. \\
2 PT & Aguentar & 1 Reação. Ao sofrer dano que te levaria a 0 PV, fica com 1 PV em vez disso. Uma vez por combate. \\
3 PT & Explosão de Aura & 1 Ação. Criaturas a 3m fazem teste de Força (CD 8 $+$ Força $+$ Rank) ou são arremessadas 4,5m e ficam Caídas. \\
\end{tabelalivro}

\begin{aviso}{\enquote{Aguentar} é esta manobra, e só ela}
    Toda árvore do Corpo tem acesso à manobra acima. A técnica de Cavalaria e Escudos que
    intercepta dano por 1 PT chama-se \textbf{Aguentar o Baque} --- nome diferente, de
    propósito, porque um personagem de Escudos tem as duas na ficha e \enquote{gastei
    Aguentar} precisava ser uma frase sem ambiguidade.
\end{aviso}

\subsection{O Triângulo dos Estilos}

\begin{tabelalivro}{L{3cm} L{6.6cm} C{2.2cm} C{2.2cm}}{As Três Grandes Escolas}
\cabtabela{\textbf{Estilo} & \textbf{Filosofia} & \textbf{Vence} & \textbf{Perde}}
Deus da Espada & Velocidade e agressão; matar em um golpe. Sem defesa, sem contra-ataque. & Deus do Norte & Deus da Água \\
Deus da Água & Defesa e contragolpe. Deixa o inimigo atacar e devolve. & Deus da Espada & Deus do Norte \\
Deus do Norte & Sobreviver e vencer por qualquer meio. Truques, terreno, improviso. & Deus da Água & Deus da Espada \\
\end{tabelalivro}

\begin{nota}{Regra da Vantagem de Estilo}
    Quando você luta contra um praticante do estilo que o seu contra-ataca, e ambos possuem
    Rank naqueles estilos: você rola com Vantagem em todas as Disputas contra ele, e as
    Reações defensivas dele falham automaticamente contra sua primeira Ação de cada turno.
    Se o Rank dele for dois ou mais acima do seu, a vantagem se anula --- treino bruto supera
    a tabela de tipos.
\end{nota}

% ---------------------------------------------------------------------------
\section{A Árvore de Utilidade --- Sistemas Compartilhados}

O terceiro pilar não compete em dano. O que a Utilidade faz é decidir as \textbf{condições}
em que a luta, a negociação ou o roubo acontecem.

\begin{itemize}
    \item \textbf{Sem Escada de Dados:} árvores de Utilidade não recebem degraus no Dado de
          Arma.
    \item \textbf{Sem Touki, nunca:} nenhum patamar de Utilidade concede Manto de Touki nem
          PT.
    \item \textbf{O Rank soma só nas Perícias que a sua árvore cobre} --- nunca em todas as
          Perícias do jogo.
\end{itemize}

\begin{tabelalivro}{L{5cm} L{9cm}}{Perícias cobertas pelo Bônus de Rank}
\cabtabela{\textbf{Árvore} & \textbf{Perícias cobertas}}
Furtividade e Armadilhas (Ladino) & Furtividade e Percepção --- só essas duas. \\
Bardo e Interação & Atuação, Persuasão, Intuição, História. \\
Navegação e Liderança (Tático) & Sobrevivência, Natureza, Investigação, Percepção (rastreio). \\
\end{tabelalivro}

Duas árvores cobrem \textbf{Percepção}. Se você tiver Rank em ambas, use o \textbf{maior}
Bônus de Rank entre as duas --- nunca some os dois, e nunca use os dois pra dobrar a vantagem
no mesmo teste.

\subsection{Pontos de Preparação (PP)}

Magos gastam PM pra fazer algo acontecer agora. Guerreiros gastam PT pra aguentar o que está
acontecendo agora. A Utilidade gasta um recurso que nenhum dos dois tem: \rec{PP} serve pra
declarar que algo \emph{já aconteceu antes}.

\begin{formula}
    \textbf{PP Máximos $=$ Intelecto $+$ atributo-chave da árvore} (mínimo 1),\\
    $+1$ por patamar a partir do terceiro
\end{formula}

Ladino usa Agilidade, Bardo usa Espírito. Com mais de uma árvore de Utilidade, a reserva é
única --- use o maior atributo-chave. Recupera-se tudo em Descanso Longo.

\begin{aviso}{Navegação e Liderança é a exceção}
    O atributo-chave do Tático \emph{é} Intelecto: nesse caso ele não conta duas vezes ---
    some, no lugar, o seu \textbf{Bônus de Rank} naquela árvore. Sem essa cláusula o Tático
    tinha a maior reserva de PP do livro (20 no Imperador, contra 15 do Bardo) investindo
    \emph{um} atributo onde as outras duas árvores de Utilidade investem dois.
\end{aviso}

Gastando 1 PP, você declara em voz alta um fato sobre o passado que passa a ser verdade no
jogo. Cinco condições: (1) precisa caber no seu Escopo; (2) precisa caber no seu Domínio;
(3) é sempre pretérito; (4) custa 2 PP se resolver o obstáculo central da cena; (5) o Mestre
não pode negar, mas pode anexar uma complicação.

\subsection{As Três Faixas}

A solução pra três árvores não virarem \enquote{a mesma pessoa com roupa diferente}: dividir
o passado em três domínios que não se tocam, cada um travado numa faixa exclusiva de combate.

\begin{tabelalivro}{L{3.2cm} L{3.6cm} L{3.6cm} L{3.6cm}}{As Três Faixas}
\cabtabela{ & \textbf{Ladino} & \textbf{Bardo} & \textbf{Tático}}
Atributo-chave & Agilidade & Espírito & Intelecto \\
Domínio & Coisas e lugares. & Pessoas e reputação. & Tempo e logística. \\
Exemplo de fato & \enquote{Essa fechadura eu já limei.} & \enquote{O capitão da guarda me deve um favor.} & \enquote{O suprimento deles acabou anteontem.} \\
Faixa exclusiva & Dano Furtivo. & Estado emocional. & Economia de ação. \\
A pergunta dele & Como eu entro? & Quem eu convenço? & Onde e quando isso acontece? \\
\end{tabelalivro}

\begin{aviso}{A Regra da Faixa}
    Nenhuma habilidade pode invadir a faixa de outra árvore de Utilidade --- um talento que dê
    Dano Furtivo a um Bardo, ou que deixe um Ladino conceder uma Ação, está errado. A Faixa
    vale só entre Ladino, Bardo e Tático; árvores do Corpo e de Magia cruzam essas linhas
    livremente.
\end{aviso}

\begin{nota}{A régua honesta}
    Contra um Norte Imperador batendo 81 por turno, os três juntos talvez somem 30 de dano
    direto na luta inteira. E ainda assim eles são a razão de o combate ter começado com o
    grupo em cima do telhado, os reforços trancados do lado de fora, metade dos inimigos
    apavorados, e o chefe já sabendo que perdeu.
\end{nota}

% ---------------------------------------------------------------------------
\chapter*{Catálogo: Todas as Sub-árvores}
\addcontentsline{toc}{chapter}{Catálogo: Todas as Sub-árvores}
\markboth{Catálogo das Sub-árvores}{Catálogo das Sub-árvores}

As páginas a seguir são geradas diretamente dos dados do sistema. Cada patamar lista, nesta
ordem: a \simbmaestria{} \textbf{Maestria} (gratuita, automática, não conta como
conhecimento), os \textbf{Talentos}, e as \textbf{Magias} ou \textbf{Técnicas} --- com
\assinatura{} marcando a Assinatura do patamar.

% GERADO — fonte: src/data/trees/*.ts

\section{Árvore de Magia}

\subsection{Magia de Fogo}
\begin{arvorecard}{Magia de Fogo}{Magia Ofensiva}{Intelecto \textbullet\ Recurso: PM}
Dano bruto e consequência --- a única escola que destrói o que estava em volta.
\end{arvorecard}
\patamar{Principiante}{Dado de PV \dado{1d4+1}}
\maestria{Chama Viva}{Acende, apaga, aquece e controla qualquer chama já existente a até 9m, sem PM e sem Ação. Você é imune a dano ígneo não-mágico e não sofre penalidade de calor extremo.}
\talento{Pavio Curto}{1}{Magias suas que aplicam Em Chamas causam +2 de dano contra alvos que já estejam Em Chamas.}
\talento{Fôlego de Forja}{1}{+2 PM por patamar seu em Fogo. Aplicado sozinho na ficha, e cresce a cada patamar novo que você abrir nela.}
\talento{Mãos de Ferreiro}{1}{Você trabalha metal, vidro e cerâmica sem forja, usando as próprias mãos.}
\magia{Bola de Fogo}{s}{2}{1}{2 Ações / encurtada 1 / silenciosa 1}{18 metros}{\textbf{Dano:} 1d8 + BC (ígneo)\par
Ataque mágico à distância. Se acertar, o alvo fica Em Chamas.}{Chama que dorme na pedra e na madeira, acorde na minha mão e vá até ele. Bola de Fogo!}
\magia{Fagulha}{\relax}{1}{1}{2 Ações / encurtada 1 / silenciosa 1}{9 metros}{Sem dano direto. Um ponto de calor acende qualquer coisa inflamável que você consiga ver.}{\relax}
\magia{Toque Escaldante}{\relax}{1}{2}{2 Ações / encurtada 1 / silenciosa 1}{Toque}{\textbf{Dano:} 2d6 + BC (ígneo)\par
Ataque corpo a corpo mágico. Aplica Em Chamas. Se você estiver com metade ou menos dos PV, causa +1d6.}{\relax}
\magia{Muro de Chamas}{\relax}{1}{3}{2 Ações / encurtada 1 / silenciosa 1}{Linha de 9 metros}{\textbf{Dano:} 3d6 de dano ígneo a quem atravessa\par
Parede de fogo de 3m de altura por 1 minuto. Atravessá-la aplica Em Chamas. Bloqueia visão e assusta animais.}{\relax}
\magia{Clarão}{\relax}{1}{1}{2 Ações / encurtada 1 / silenciosa 1}{Esfera de 6 metros}{Teste de Resistência de Vigor (CD 8 + BC) ou o alvo fica Cego até o fim do próximo turno. Não causa dano nem incendeia nada.}{\relax}
\patamar{Intermediário}{Dado de PV \dado{1d4+1}}
\maestria{Propagação}{A condição Em Chamas aplicada por você causa 1d8 em vez de 1d6. Uma vez por turno, sem gastar Ação, você faz o fogo de um alvo Em Chamas saltar para outra criatura a até 3m (teste de Agilidade CD 8 + BC ou também pega fogo).}
\talento{Calor Dirigido}{1}{Você escolhe até Intelecto criaturas na área das suas magias de Fogo. Elas não são afetadas.}
\talento{Combustão Lenta}{1}{Em Chamas aplicada por você não pode ser apagada gastando Ação --- só com água, frio ou submersão.}
\talento{Nada Sobra}{1}{Suas magias causam dano dobrado contra objetos, estruturas, cordas, portas e barreiras não-mágicas.}
\magia{Lança de Fogo}{s}{2}{3}{2 Ações / encurtada 1 / silenciosa 1}{27 metros}{\textbf{Dano:} 3d8 + BC (ígneo, ignora Resistência); +2d8 contra alvo Em Chamas\par
Ataque mágico à distância. Fogo comprimido em haste: não incendeia, perfura e sela a ferida no caminho.}{Calor branco que não conhece fumaça, tome a forma da lança e atravesse. Lança de Fogo!}
\magia{Explosão}{\relax}{1}{4}{2 Ações / encurtada 1 / silenciosa 1}{Esfera de 6 metros}{\textbf{Dano:} 3d6 + BC de explosão (+3d6 contra alvos Em Chamas)\par
Teste de Resistência de Agilidade (CD 8 + BC), metade se passar. Falha: arremessadas 3m.}{\relax}
\magia{Chuva de Brasas}{\relax}{1}{3}{2 Ações / encurtada 1 / silenciosa 1}{Esfera de 9m de raio}{Sem dano imediato. Teste de Agilidade (CD 8 + BC) ou Em Chamas. Inflamáveis na área acendem; terreno difícil por brasa acesa 1 minuto.}{\relax}
\magia{Sopro}{\relax}{1}{3}{2 Ações / encurtada 1 / silenciosa 1}{Cone de 9 metros}{\textbf{Dano:} 2d10 + BC (ígneo)\par
Teste de Agilidade para metade do dano. Em Chamas a quem falhar.}{\relax}
\patamar{Avançado}{Dado de PV \dado{1d6+2}}
\maestria{Termodinâmica Inversa}{Você controla a temperatura dentro das suas áreas: pode poupar aliados/objetos, ou fazer o fogo queimar só um material escolhido. Desbloqueia o direito de combinar escolas (Magia Combinada, Cap. 2).}
\talento{Coração de Brasa}{2}{Você é imune a todo dano ígneo, mágico ou não, e à condição Em Chamas.}
\talento{Detonação}{2}{Uma vez por turno, sem gastar Ação, apague a condição Em Chamas de um alvo para causar imediatamente 3d8 de dano de explosão nele.}
\talento{Cântico de Cinzas}{2}{Magias de Fogo de rank Intermediário ou inferior custam 1 PM a menos (mínimo 1).}
\magia{Tempestade de Fogo}{s}{3}{6}{3 Ações / encurtada 2 / silenciosa 1}{Esfera de 12m de raio}{\textbf{Dano:} 6d8 + BC (ígneo)\par
Teste de Agilidade (CD 8 + BC), metade se passar. Falha: Em Chamas. A área continua queimando 1 minuto: quem começar o turno dentro sofre +2d6.}{Vento que alimenta e chama que devora, girem juntos até que não reste ar nem nome neste lugar. Tempestade de Fogo!}
\magia{Coluna Solar}{\relax}{2}{5}{3 Ações / encurtada 2 / silenciosa 1}{45 metros}{\textbf{Dano:} 5d8 + BC (ígneo)\par
Pilar de fogo de 3m de raio. Teste de Agilidade para metade, Em Chamas automático em falha. Ignora Cobertura que não seja um teto sólido.}{\relax}
\magia{Vapor Seco}{\relax}{2}{4}{3 Ações / encurtada 2 / silenciosa 1}{Esfera de 9m}{\textbf{Dano:} 4d8 de dano ígneo\par
Requer 1 patamar em Vento (ou aliado mago de Vento conjurando junto). Teste de Vigor (CD 8 + BC): falha não fala nem recita cânticos por 2 turnos e fica Em Chamas por dentro (não abafável com 1 Ação).}{\relax}
\magia{Escudo de Cinzas}{\relax}{2}{4}{3 Ações / encurtada 2 / silenciosa 1}{Pessoal}{1 minuto: quem te atingir corpo a corpo sofre 2d6 ígneo e fica Em Chamas. Você recebe Resistência a dano físico de armas mundanas.}{\relax}
\patamar{Santo}{Dado de PV \dado{1d6+2}}
\maestria{Domínio da Combustão}{Suas magias de Fogo ignoram Resistência a dano ígneo e de explosão. Você pode conjurar qualquer magia de Fogo de rank Avançado ou inferior com metade da área pelo dobro do dano, ou o inverso. Fogo aceso por você não se apaga enquanto você quiser.}
\talento{Segunda Ignição}{3}{Uma vez por combate, ao conjurar magia de Fogo de rank Avançado ou inferior, conjure-a duas vezes pagando o PM uma vez; a segunda pode ter alvo diferente.}
\magia{Mar de Chamas}{s}{4}{12}{4 Ações / encurtada 3 / silenciosa 2}{Raio de 60m}{\textbf{Dano:} 10d8 + BC (ígneo)\par
\textbf{Ritual:} não pode ser encurtado.\par
Teste de Agilidade com Desvantagem. Falha: dano cheio e Em Chamas. Sucesso: metade. Aliados só são poupados pelo talento Calor Dirigido. Área continua em chamas por 10 minutos.}{Que o chão lembre do dia em que foi lava. Que nada aqui volte a ter nome. Mar de Chamas!}
\magia{Corpo de Fogo}{\relax}{3}{8}{4 Ações / encurtada 3 / silenciosa 2}{Pessoal}{\textbf{Dano:} 3d6 de dano ígneo a quem tocar\par
1 minuto: Resistência a todo dano físico, atravessa frestas, ignora terreno difícil. Não pode usar itens, empunhar armas nem ser curado.}{\relax}
\patamar{Rei}{Dado de PV \dado{1d8+3}}
\maestria{Plasma}{Você conjura dano de plasma, que ignora Resistência e Imunidade a dano ígneo e derrete metal, pedra e Touki com igual indiferença. Dano de plasma reduz permanentemente em 2 a CA de armadura não-mágica atingida, cumulativamente.}
\talento{Ponto de Fusão}{4}{Seu dano de plasma passa a ignorar também barreiras mágicas com PV --- ele fura, em vez de gastar.}
\magia{Flashover}{s}{5}{14}{5 Ações / encurtada 4 / silenciosa 3}{Raio de 90m}{\textbf{Dano:} 12d10 + BC (ígneo, d12 contra alvo já Em Chamas)\par
Teste de Vigor com Desvantagem (exceto submersos/barreira/clima chuvoso). Falha: dano cheio e Em Chamas incombatível. Sucesso: metade.}{Não peço chama. Peço o instante em que tudo o que respira descobre que já estava queimando. Flashover!}
\magia{Lança de Plasma}{\relax}{4}{10}{5 Ações / encurtada 4 / silenciosa 3}{45 metros}{\textbf{Dano:} 8d8 + BC (plasma)\par
Ataque mágico à distância. Atravessa o alvo em linha reta e atinge tudo atrás por mais 15m com metade do dano.}{\relax}
\patamar{Imperador}{Dado de PV \dado{1d8+3}}
\maestria{A Segunda Estrela}{Suas magias de Fogo e plasma não podem ser aparadas, refletidas, absorvidas nem redirecionadas por rank inferior ao seu. Criaturas reduzidas a 0 PV são reduzidas a cinza --- sem ressurreição abaixo de rank Deus. Uma vez por turno, conjure magia de Fogo Avançado ou inferior em Silenciosa sem gastar Ação.}
\talento{A Chama Que Escolhe}{4}{Você poupa automaticamente um número de criaturas igual ao seu Espírito em qualquer magia sua, incluindo o Sol Menor.}
\magia{Sol Menor}{s}{6}{22}{6 Ações / silenciosa 4}{Esfera de 30m}{\textbf{Dano:} 14d12 de plasma (20d12 contra alvos Em Chamas)\par
\textbf{Ritual:} não pode ser encurtado.\par
Teste de Vigor com Desvantagem Absoluta. Aliados não são poupados automaticamente. Construções, terreno e cadáveres na área deixam de existir; a cratera é permanente.}{Eu junto tudo o que o fogo seria em mil anos e devolvo em um segundo. Que o céu abaixe até o chão. Sol Menor!}
\magia{Nunca Apaga}{\relax}{5}{16}{6 Ações / silenciosa 4}{Raio de 1 km}{A região pega fogo e continua pegando fogo por três dias, sem combustível, ignorando chuva de rank Santo ou inferior. Rotas fecham, cidades evacuam.}{\relax}

\subsection{Magia de Água}
\begin{arvorecard}{Magia de Água}{Magia Ofensiva}{Intelecto \textbullet\ Recurso: PM}
Atrição e controle de terreno --- vence decidindo onde a luta acontece, não trocando golpes.
\end{arvorecard}
\patamar{Principiante}{Dado de PV \dado{1d4+1}}
\maestria{Afinidade Aquática}{Cria, move, aquece levemente ou evapora até 20 litros de água limpa por minuto, sem gastar PM nem Ação. Não causa dano, mas serve pra tudo o mais: encher cantis, apagar fogueiras, limpar ferimentos, dar água a um cavalo.}
\talento{Condutor de Gelo}{1}{Magias de gelo suas contra alvos Molhados impõem Desvantagem no teste de resistência.}
\talento{Nascente de Mana}{1}{+2 PM por patamar seu em Água. Aplicado sozinho na ficha, e cresce a cada patamar novo que você abrir nela.}
\talento{Mão Firme}{1}{Você não sofre Desvantagem ao conjurar com um inimigo adjacente a você.}
\magia{Bola de Água}{\relax}{1}{1}{2 Ações / encurtada 1 / silenciosa 1}{18 metros}{\textbf{Dano:} 1d6 + BC (contundente)\par
Ataque mágico à distância. O alvo é empurrado 1,5m e fica Molhado.}{Que a grande proteção da água esteja no lugar que buscas. Eu clamo por um riacho refrescante e borbulhante aqui e agora. Bola de Água!}
\magia{Flecha de Água}{s}{2}{1}{2 Ações / encurtada 1 / silenciosa 1}{27 metros}{\textbf{Dano:} 1d8 + BC (perfurante)\par
Ataque mágico à distância. Se acertar, o alvo fica Molhado. A magia comum que define um mago de Água.}{Água que flui, tome a forma perfurante da caçada e atinja meu inimigo. Flecha de Água!}
\magia{Impacto de Gelo}{\relax}{1}{2}{2 Ações / encurtada 1 / silenciosa 1}{18 metros}{\textbf{Dano:} 1d4 contundente + 1d6 de frio (frio dobra contra Molhado)\par
Teste de Resistência de Agilidade (CD 8 + BC). Falha: deslocamento reduzido em 3m até o fim do próximo turno.}{Coloco diante de ti um berço de gelo como desejas, agora libere tuas correntes glaciais. Impacto de Gelo!}
\magia{Lâmina de Gelo}{\relax}{1}{2}{2 Ações / encurtada 1 / silenciosa 1}{Toque}{\textbf{Dano:} 1d8 + BC (cortante) + 1d4 de frio\par
Cria uma arma de gelo por 1 minuto. Ataques com ela usam Força ou Intelecto, à sua escolha.}{Frio cortante do inverno, esculpa a arma que ceifará os meus inimigos. Lâmina de Gelo!}
\magia{Escudo de Água}{\relax}{1}{2}{2 Ações / encurtada 1 / silenciosa 1}{6 metros}{Barreira de água de 3m de largura por 2 turnos. Aliados atrás recebem Cobertura Superior (+5 CA) contra ataques físicos e projéteis. Dano ígneo que a atravessa é reduzido à metade.}{Espírito das correntes, erga-se da terra e forme a muralha que me protege do calor. Escudo de Água!}
\magia{Névoa Densa}{\relax}{1}{1}{2 Ações / encurtada 1 / silenciosa 1}{Esfera de 6 metros}{Área fortemente obscurecida por 5 minutos (efetivamente cega quem está dentro). Vento forte dissipa em 1 turno.}{Respiração fria da manhã, roube deles a visão do mundo. Névoa Densa!}
\patamar{Intermediário}{Dado de PV \dado{1d4+2}}
\maestria{Cântico Fluido}{Suas magias de Rank Principiante não sofrem mais penalidade na versão Encurtada --- dano cheio, área cheia. Falhas críticas em magias de Principiante apenas falham em silêncio.}
\talento{Fluidez Defensiva}{1}{1 Ação e 2 PM: você desliza 4,5m em qualquer direção sobre água conjurada, sem provocar ataques de oportunidade.}
\talento{Pressão Profunda}{1}{Magias que empurram agora empurram o dobro da distância. Colisão com obstáculo causa +1d6 contundente.}
\talento{Cristalização Rápida}{1}{Uma vez por combate, aplique Congelado a um alvo já Molhado sem conjurar magia --- apenas 1 Ação e 2 PM.}
\magia{Canhão de Água}{\relax}{1}{3}{2 Ações / encurtada 1 / silenciosa 1}{Linha de 18m × 1,5m}{\textbf{Dano:} 3d6 + BC (contundente)\par
Teste de Resistência de Força (CD 8 + BC). Falha: empurradas 4,5m e ficam Molhadas. Sucesso: metade do dano, mas ficam Molhadas mesmo assim.}{Flexível espírito da água... varra todas as coisas com seu poder oculto. Canhão de Água!}
\magia{Lança de Gelo}{s}{2}{3}{2 Ações / encurtada 1 / silenciosa 1}{27 metros}{\textbf{Dano:} 2d8 + BC (perfurante) + 1d8 de frio (dobra contra Molhado)\par
Ataque mágico à distância. Se acertar, o alvo fica Molhado pelo degelo do impacto. A magia que consagra um Intermediário.}{Águas eternas, condensem-se na haste que perfura a armadura e o osso. Lança de Gelo!}
\magia{Pilar de Gelo}{\relax}{1}{3}{2 Ações / encurtada 1 / silenciosa 1}{18 metros}{\textbf{Dano:} 2d6 + BC (contundente)\par
Coluna de 1,5m de raio e 3m de altura. Teste de Agilidade (CD 8 + BC): falha sofre o dano e é arremessado ao topo. Dura 10 minutos e concede Cobertura Superior.}{Águas adormecidas sob a terra, congelem e ergam-se subitamente para os céus. Pilar de Gelo!}
\magia{Respingos de Água}{\relax}{1}{2}{2 Ações / encurtada 1 / silenciosa 1}{Esfera de 9m de raio}{Sem dano. Toda criatura na área fica Molhada, sem teste de resistência. O chão vira terreno difícil por 1 minuto.}{Espalhe as gotas que caem, cubra o mundo em água. Respingos de Água!}
\magia{Enxurrada}{\relax}{1}{3}{2 Ações / encurtada 1 / silenciosa 1}{Cone de 9 metros}{\textbf{Dano:} 2d6 + BC (contundente)\par
Teste de Resistência de Força (CD 8 + BC). Falha: empurrão de 4,5m, Caído e Molhado. Sucesso: metade do dano, Molhado.}{Correnteza que arranca a montanha, desça sobre eles e leve tudo consigo. Enxurrada!}
\patamar{Avançado}{Dado de PV \dado{1d6+2}}
\maestria{Termodinâmica Aplicada}{Você troca livremente o dano de frio das suas magias por contundente (água pressurizada) ou ígneo (vapor), sem alterar os dados. Desbloqueia o direito de aprender e executar Magias Combinadas (Cap. 2).}
\talento{Zero Perfurante}{2}{Seu gelo ignora Resistência a dano de frio.}
\talento{Mestre da Adaptação}{2}{Suas magias Encurtadas de Água não perdem mais dados de dano, apenas a redução de área.}
\talento{Núcleo Gélido}{2}{Você é imune a dano de frio não-mágico e não sofre penalidades de clima gelado.}
\magia{Quebra de Gelo}{\relax}{2}{4}{3 Ações / encurtada 2 / silenciosa 1}{27 metros}{\textbf{Dano:} 3d8 + BC (perfurante) + 2d6 de frio (dobra contra Molhado)\par
Ataque mágico à distância. Contra objetos e estruturas, o dano é dobrado.}{Magníficos espíritos da água, atinjam o insolente com sua majestosa espada de gelo e estilhacem sua alma! Quebra de Gelo!}
\magia{Corte de Gelo}{\relax}{2}{4}{3 Ações / encurtada 2 / silenciosa 1}{Linha de 18 metros}{\textbf{Dano:} 3d8 + BC (cortante) + 1d6 de frio\par
Teste de Resistência de Agilidade (CD 8 + BC) para cada criatura na linha. Falha: dano cheio. Sucesso: metade.}{Frio implacável, tome a forma da execução perfeita! Eu a convoco para abater e fatiar o meu inimigo! Corte de Gelo!}
\magia{Campo de Gelo}{\relax}{2}{4}{3 Ações / encurtada 2 / silenciosa 1}{Esfera de 9m de raio}{\textbf{Dano:} 2d8 de frio (dobra contra Molhado)\par
Teste de Resistência de Vigor (CD 8 + BC). Falha: deslocamento reduzido a 0 até o fim do próximo turno. Alvos Molhados que falharem ficam Congelados.}{Deusa Azul que desce dos céus, empunhe seu cajado e cubra este mundo maldito em geada! Campo de Gelo!}
\magia{Nevasca}{s}{3}{5}{3 Ações / encurtada 2 / silenciosa 1}{Explosão de 9m centrada em você}{\textbf{Dano:} 2d6 perfurante + 3d6 de frio + BC (dobra contra Molhado)\par
Teste de Resistência de Agilidade (CD 8 + BC). Falha: dano cheio e arremessadas 3m. Sucesso: metade e não empurradas. Você não é afetado.}{Soberano envolto em branco absoluto, cujo frio glacial rouba todo o calor. Congele aqueles que ousam se aproximar! Nevasca!}
\magia{Tempestade}{\relax}{2}{5}{3 Ações / encurtada 2 / silenciosa 1}{Raio de 1 km}{Chuva pesada por 1 hora. Todos sob ela mantêm Molhado permanentemente. Visibilidade reduzida à metade. Fogueiras e magia de fogo rank Intermediário ou inferior não funcionam ao ar livre.}{Nuvens carregadas que viajam pelos ventos uivantes, derramem suas lágrimas sobre esta terra árida e lavem o mundo! Tempestade!}
\magia{Fortaleza de Gelo}{\relax}{2}{5}{3 Ações / encurtada 2 / silenciosa 1}{9 metros}{Muralha ou domo com 80 PV (Cobertura Total). Se conjurada como Reação em Silenciosa, surge com apenas 30 PV. Dura 10 minutos.}{Guardião das geleiras eternas, erga-se das profundezas e torne-se o escudo intransponível que protegerá minha vida. Fortaleza de Gelo!}
\patamar{Santo}{Dado de PV \dado{1d6+3}}
\maestria{Domínio Climático}{Você é imune aos danos e efeitos colaterais das suas próprias magias de área e de clima, e pode poupar um número de aliados igual ao seu Espírito. Enxerga perfeitamente através de chuva, névoa e nevasca, e mantém uma magia de clima ativa sem gastar concentração nem Ações.}
\talento{Olho da Tempestade}{3}{Você mantém duas magias de clima simultaneamente, e o custo em PM de qualquer magia de clima da escola é reduzido à metade.}
\magia{Cumulonimbus}{s}{4}{10}{4 Ações / encurtada 3 / silenciosa 2}{Raio de 1,5 km}{\textbf{Ritual:} não pode ser encurtado.\par
A nuvem paira por 1 minuto. Todos na área ficam Molhados incondicionalmente. Enquanto ativa, gastando 1 Ação e 2 PM você castiga um alvo visível com um relâmpago: teste de Agilidade (CD 8 + BC) ou 4d10 de dano elétrico (metade se passar).}{Grande espírito da água e príncipe imperial do relâmpago que ascende aos céus! Conceda meu desejo e traga uma bênção selvagem. Cumulonimbus!}
\magia{Prisão de Gelo Eterno}{\relax}{3}{8}{4 Ações / encurtada 3 / silenciosa 2}{18 metros}{\textbf{Dano:} 4d8 de frio ao falhar (dobra contra Molhado)\par
Teste de Resistência de Vigor (CD 8 + BC), com Desvantagem se Molhado. Falha: Paralisado em animação suspensa por até 1 hora, imune a dano. Dano ígneo aplicado ao bloco concede novo teste.}{Coração parado do inverno profundo, tome esta alma e guarde-a onde o tempo não alcança. Prisão de Gelo Eterno!}
\magia{Maremoto}{\relax}{3}{9}{4 Ações / encurtada 3 / silenciosa 2}{Linha de 45m × 9m}{\textbf{Dano:} 8d8 + BC (contundente)\par
Teste de Resistência de Força (CD 8 + BC). Falha: dano cheio, arrastadas 9m, Caídas e Molhadas. Sucesso: metade, sem arrasto. Estruturas de madeira sofrem dano dobrado.}{Mar que engoliu continentes antes que houvesse nomes, lembre-se do que você é e venha reivindicar esta terra! Maremoto!}
\patamar{Rei}{Dado de PV \dado{1d8+3}}
\maestria{Condutividade}{A escola de Água passa a incluir o elemento Eletricidade. Todo dano elétrico que você causar a um alvo Molhado impõe Desvantagem no teste de resistência dele; se tirar 5 ou menos, fica Atordoado por 1 turno.}
\talento{Soberania Elétrica}{4}{O Atordoamento causado pela Maestria Condutividade dura 2 turnos e se estende a todas as criaturas adjacentes ao alvo original.}
\magia{Relâmpago}{s}{5}{12}{5 Ações / encurtada 4 / silenciosa 3}{Alcance ilimitado (linha de visão)}{\textbf{Dano:} 8d10 + BC (elétrico)\par
Pré-requisito: Cumulonimbus ativa acima de você. Ignora bônus de CA por Touki. Contra armadura metálica, acerto é Crítico automático. Contra alvo Molhado, dano dobrado.}{Oh espíritos das águas magníficas, eu imploro ao Príncipe do Trovão! Deixe seu poder radiante ensinar a este inimigo insolente que o Imperador ainda reina supremo! Relâmpago!}
\magia{Era Glacial}{\relax}{4}{11}{5 Ações / encurtada 4 / silenciosa 3}{Esfera de 30m}{\textbf{Dano:} 3d10 de frio (sem teste, a quem começar o turno dentro)\par
A área vira terreno congelado por 24 horas; deslocamento reduzido à metade dentro dela. Corpos d'água congelam sólidos e viram terreno transitável.}{Estação que não pede permissão, desça sobre esta terra e cubra o que os homens construíram. Era Glacial!}
\patamar{Imperador}{Dado de PV \dado{1d8+4}}
\maestria{O Silêncio Primordial}{Seu dano de frio ignora completamente Resistência e Imunidade. Criaturas reduzidas a 0 PV pelas suas magias de gelo cristalizam em pó de diamante --- não podem ser ressuscitadas por nada abaixo de rank Deus. Uma vez por turno, conjure uma magia de Água de rank Avançado ou inferior em Conjuração Silenciosa sem gastar Ação.}
\talento{Essência do Inverno}{4}{Seu dano de frio passa a ignorar também Invulnerabilidade e proteção mágica de rank inferior ao seu. Uma vez por Descanso Longo, conjure Zero Absoluto pagando metade do PM.}
\magia{Zero Absoluto}{s}{6}{20}{6 Ações / silenciosa 4}{Esfera de 45 metros}{\textbf{Dano:} 12d12 de frio (24d12 contra alvo Molhado)\par
\textbf{Ritual:} não pode ser encurtado.\par
Teste de Resistência de Vigor com Desvantagem Absoluta. Aliados são isolados magicamente e não sofrem efeito algum. Criaturas reduzidas a 0 PV são eternamente petrificadas --- sem ressurreição, sem cura, sem corpo.}{Silêncio primordial, onde até o próprio tempo congela e toda a chama cessa. Que o conceito de calor deixe de existir. Zero Absoluto!}
\magia{Dilúvio}{\relax}{5}{18}{6 Ações / silenciosa 4}{Raio de 3 km}{\textbf{Ritual:} não pode ser encurtado.\par
Efeito narrativo de escala regional. Todos na área permanecem Molhados por 3 dias, todo terreno vira difícil, magia de fogo de rank Santo ou inferior falha automaticamente ao ar livre, e qualquer exército em campo aberto perde suprimentos.}{Céus que testemunharam a queda de três mundos, lembrem-se de como se limpa uma face. Abram-se, e não parem. Dilúvio!}

\subsection{Magia de Vento}
\begin{arvorecard}{Magia de Vento}{Magia Ofensiva}{Intelecto \textbullet\ Recurso: PM}
Corte, som e mobilidade --- a escola que anula a distância, inclusive a sua própria.
\end{arvorecard}
\patamar{Principiante}{Dado de PV \dado{1d4+2}}
\maestria{Brisa}{Controle constante e gratuito de ar num raio de 18 metros, sem PM e sem Ação. Você nunca sofre dano de queda de até 15 metros, e desce devagar de qualquer altura.}
\talento{Pés Leves}{1}{+3 metros de Deslocamento, e você não deixa pegadas nem faz ruído ao andar.}
\talento{Ouvido do Vento}{1}{Você escuta qualquer conversa a até 60 metros, desde que exista ar entre vocês. Vantagem em Percepção auditiva.}
\talento{Reserva de Ar}{1}{+2 PM por patamar seu em Vento. Aplicado sozinho na ficha, e cresce a cada patamar novo que você abrir nela.}
\magia{Lâmina de Vento}{s}{2}{1}{2 Ações / encurtada 1 / silenciosa 1}{27 metros}{\textbf{Dano:} 1d10 + BC (cortante)\par
Ataque mágico à distância. Se acertar, o alvo fica Desequilibrado.}{Ar que passa por tudo sem pedir licença, tome fio e vá. Lâmina de Vento!}
\magia{Empurrão}{\relax}{1}{1}{2 Ações / encurtada 1 / silenciosa 1}{18 metros}{Teste de Força (CD 8 + BC) ou o alvo é empurrado 6m na direção escolhida e fica Desequilibrado. Sem dano.}{\relax}
\magia{Passo de Vento}{\relax}{1}{2}{2 Ações / encurtada 1 / silenciosa 1}{Pessoal}{1 Ação: desloque-se 18 metros em qualquer direção, inclusive para cima, sem provocar ataques de oportunidade e ignorando terreno difícil. Desce suavemente se terminar no ar.}{\relax}
\magia{Sopro de Poeira}{\relax}{1}{1}{2 Ações / encurtada 1 / silenciosa 1}{Cone de 6 metros}{Teste de Vigor (CD 8 + BC) ou o alvo fica Cego e Desequilibrado até o fim do próximo turno dele. Sem dano.}{\relax}
\magia{Vácuo Localizado}{\relax}{1}{2}{2 Ações / encurtada 1 / silenciosa 1}{9 metros}{\textbf{Dano:} 1d6\par
Remove o ar ao redor da cabeça do alvo. Ele não consegue recitar cânticos, falar nem gritar por 1 turno. Não funciona em quem não respira.}{\relax}
\patamar{Intermediário}{Dado de PV \dado{1d6+2}}
\maestria{Sem Peso}{Passo de Vento passa a custar 1 PM e pode ser usado uma vez por turno sem gastar Ação. Você ignora terreno difícil e não pode ser Caído nem Atolado enquanto consciente. Ataques à distância mundanos contra você sofrem Desvantagem.}
\talento{Corrente de Apoio}{1}{Quando um aliado a até 18m conjurar magia de área, gaste 1 PM como Reação para aumentar a área dela em metade.}
\talento{Corte Fino}{1}{Suas magias cortantes de Vento causam +1d6 contra criaturas sem armadura.}
\talento{Fuga Perfeita}{1}{Ao ser reduzido a 0 PV, você é automaticamente arremessado 9m para longe da fonte do dano antes de cair.}
\magia{Estrondo Sônico}{s}{2}{3}{2 Ações / encurtada 1 / silenciosa 1}{Linha de 27 metros}{\textbf{Dano:} 3d8 + BC (contundente)\par
Teste de Força (CD 8 + BC). Falha: dano cheio, empurradas 9m e Desequilibradas. Sucesso: metade.}{Ar que se dobra até doer, endireite de uma vez e leve tudo o que estiver na frente. Estrondo Sônico!}
\magia{Foice de Vácuo}{\relax}{1}{3}{2 Ações / encurtada 1 / silenciosa 1}{36 metros}{\textbf{Dano:} 2d10 + BC (cortante)\par
Ataque mágico à distância que ignora Cobertura. Contra alvo Desequilibrado, crita em 19-20 e ignora metade da CA de armadura.}{\relax}
\magia{Ciclone}{\relax}{1}{4}{2 Ações / encurtada 1 / silenciosa 1}{Esfera de 9m de raio}{\textbf{Dano:} 2d8 + BC (contundente)\par
Teste de Força (CD 8 + BC). Falha: dano, arremessadas 6m para longe do centro, Desequilibradas e Caídas. Vento forte na área por 1 minuto --- projéteis mundanos erram.}{\relax}
\magia{Asas Emprestadas}{\relax}{1}{3}{2 Ações / encurtada 1 / silenciosa 1}{Toque}{10 minutos: o alvo recebe Deslocamento de Voo igual ao dobro do deslocamento dele. Se acabar no ar, desce suavemente.}{\relax}
\patamar{Avançado}{Dado de PV \dado{1d6+2}}
\maestria{Voo}{Você recebe Deslocamento de Voo permanente igual ao seu Deslocamento normal, sem custo de PM, enquanto consciente e não Exausto. Desbloqueia o direito de combinar escolas (Magia Combinada, Cap. 2).}
\talento{Vento Constante}{2}{Você mantém duas magias de Vento sustentadas simultaneamente sem concentração.}
\talento{Redirecionar}{2}{1 Reação e 2 PM: um ataque à distância mundano ou projétil mágico dirigido a você é desviado para outra criatura a até 9m.}
\talento{Corpo de Corrente}{2}{Além de Caído e Atolado, você é imune a Preso e Agarrado, e atravessa qualquer fresta por onde caiba ar.}
\magia{Nova Congelante}{s}{3}{6}{3 Ações / encurtada 2 / silenciosa 1}{Esfera de 12m}{\textbf{Dano:} 6d8 de frio (já contando a duplicação por Molhado)\par
Requer 1 patamar em Água (ou aliado mago de Água conjurando junto). Todos na área ficam Molhados e imediatamente Congelados, sem teste.}{Umidade que viaja comigo, pare no meio do caminho e escolha ser vidro. Nova Congelante!}
\magia{Prisão de Ar}{\relax}{2}{5}{3 Ações / encurtada 2 / silenciosa 1}{27 metros}{Teste de Força (CD 8 + BC). Falha: o alvo é erguido 6m do chão, Preso e Desequilibrado por 1 minuto (repete teste no fim de cada turno).}{\relax}
\magia{Guilhotina de Vácuo}{\relax}{2}{5}{3 Ações / encurtada 2 / silenciosa 1}{Linha de 45 metros}{\textbf{Dano:} 5d8 + BC (cortante, +2d8 contra Desequilibrado)\par
Corta madeira, corda e tecido com facilidade; não corta pedra.}{\relax}
\magia{Tomar o Ar}{\relax}{2}{4}{3 Ações / encurtada 2 / silenciosa 1}{Esfera de 6m}{\textbf{Dano:} 3d6 por turno\par
Remove o ar da área por 3 turnos. Teste de Vigor por turno; quem falhar sofre dano e não pode recitar cântico. Apaga fogo na área, inclusive Em Chamas.}{\relax}
\patamar{Santo}{Dado de PV \dado{1d6+3}}
\maestria{Senhor do Céu}{Seu Voo passa ao dobro do Deslocamento, e você pode pairar imóvel indefinidamente. Você controla o vento num raio de 1 km: nega voo a hostis, encalha navios, derruba flechas, dispersa névoa. Nenhum ataque à distância mundano acerta você.}
\talento{Sem Cântico}{3}{Suas magias de Vento de rank Avançado ou inferior podem ser conjuradas em Conjuração Silenciosa sem penalidade alguma.}
\magia{Tempestade Cortante}{s}{4}{11}{4 Ações / encurtada 3 / silenciosa 2}{Esfera de 30m de raio}{\textbf{Dano:} 5d8 + BC de dano cortante por turno\par
Dura 1 minuto e se move 9m por turno para onde você quiser. Desequilibrado automático em quem falhar teste de Força.}{Que o ar se lembre de que já foi lâmina. Que ele lembre mil vezes por segundo. Tempestade Cortante!}
\magia{Céu Negado}{\relax}{3}{9}{4 Ações / encurtada 3 / silenciosa 2}{Raio de 300m}{10 minutos: nenhuma criatura hostil voa, salta acima de 3m ou dispara projéteis pra dentro da área. Aliados voam livremente.}{\relax}
\patamar{Rei}{Dado de PV \dado{1d8+3}}
\maestria{Som e Vácuo}{Você conjura dano sônico, que ignora Cobertura, armadura e o Manto de Touki (não funciona no vácuo nem debaixo d'água). Cria silêncio absoluto ou ruído insuportável em até 30m, à vontade, sem PM.}
\talento{Ouvido Absoluto}{4}{Você "vê" por som num raio de 60m, atravessando paredes, escuridão e invisibilidade.}
\magia{Grito do Mundo}{s}{5}{13}{5 Ações / encurtada 4 / silenciosa 3}{Esfera de 45m}{\textbf{Dano:} 10d10 + BC de dano sônico (ignora armadura, Cobertura e Manto de Touki)\par
Teste de Vigor (CD 8 + BC). Falha: dano cheio, Atordoado 1 turno e Surdo 10 minutos. Sucesso: metade. Estruturas de pedra racham; vidro se despedaça.}{\relax}
\magia{Vazio}{\relax}{4}{12}{5 Ações / encurtada 4 / silenciosa 3}{Esfera de 18m}{\textbf{Dano:} 4d10 por turno, sem teste\par
Remove todo o ar da área por 1 minuto. Impossível recitar, gritar ou fazer fogo ali dentro. Som não existe --- nem o seu.}{\relax}
\patamar{Imperador}{Dado de PV \dado{1d8+4}}
\maestria{A Atmosfera é Sua}{Você controla a pressão do ar num raio de 5 km. Criaturas hostis a até 18m de você estão permanentemente Desequilibradas, sem teste. Uma vez por turno, conjure magia de Vento de rank Avançado ou inferior em Silenciosa sem gastar Ação.}
\talento{Nada Toca Você}{4}{Enquanto estiver voando e não Exausto, criaturas de rank Santo ou inferior não conseguem te atingir com ataque corpo a corpo algum.}
\magia{Lâmina do Horizonte}{s}{6}{20}{6 Ações / silenciosa 4}{Linha de 3 km}{\textbf{Dano:} 16d10 + BC de dano cortante (+6d10 contra Desequilibrado)\par
\textbf{Ritual:} não pode ser encurtado.\par
Teste de Agilidade com Desvantagem Absoluta. Muralhas, torres e florestas na trajetória são cortadas ao meio, permanentemente.}{Um traço. Um só. E que tudo o que estiver do lado errado dele aprenda que estava do lado errado. Lâmina do Horizonte!}
\magia{Explosão Silenciosa}{\relax}{5}{18}{6 Ações / silenciosa 4}{Esfera de 60m}{\textbf{Dano:} 12d12 dividido igualmente entre ígneo, sônico e contundente\par
Requer 1 patamar em Fogo. Teste de Vigor com Desvantagem Absoluta para metade. Ignora Resistência aos três tipos. Não faz som no momento --- ele chega depois, a quilômetros de distância.}{\relax}

\subsection{Magia de Terra}
\begin{arvorecard}{Magia de Terra}{Magia Ofensiva}{Intelecto \textbullet\ Recurso: PM}
Construção, cerco e projétil pesado --- a única escola que constrói, e o mago com mais PV do jogo.
\end{arvorecard}
\patamar{Principiante}{Dado de PV \dado{1d6+2}}
\maestria{Moldar}{Sem PM e sem Ação, você molda terra, areia, argila e pedra macia num raio de 9 metros: abrir/fechar buraco, degrau, tigela, parede baixa. Em 10 minutos você ergue um abrigo fortificado pro grupo inteiro.}
\talento{Pele de Pedra}{1}{+4 PV por patamar seu em Terra. Aplicado sozinho na ficha, e cresce a cada patamar novo que você abrir nela.}
\talento{Sentido Sísmico}{1}{Você percebe pelo chão qualquer criatura em contato com o solo num raio de 18m, mesmo invisível ou escondida.}
\talento{Mineralogista}{1}{Você identifica minérios e pedras preciosas, e sabe, ao tocar uma parede, o que existe atrás dela.}
\magia{Bala de Pedra}{s}{2}{1}{2 Ações / encurtada 1 / silenciosa 1}{27 metros}{\textbf{Dano:} 1d10 + BC (contundente)\par
Ataque mágico à distância. Contra alvo Atolado, acerta automaticamente.}{Pó que já foi montanha, lembre do peso e vá. Bala de Pedra!}
\magia{Atoleiro}{\relax}{1}{2}{2 Ações / encurtada 1 / silenciosa 1}{Esfera de 6m de raio}{Sem dano. Teste de Agilidade (CD 8 + BC) ou fica Atolado. Área permanece terreno difícil por 10 minutos.}{Terra firme, seja honesta uma vez: você nunca foi firme. Atoleiro!}
\magia{Muro de Terra}{\relax}{1}{2}{2 Ações / encurtada 1 / silenciosa 1}{9 metros}{Parede de 6m de largura por 3m de altura, meio metro de espessura, 40 PV, dura 10 minutos. Cobertura Total.}{\relax}
\magia{Lança de Pedra}{\relax}{1}{2}{2 Ações / encurtada 1 / silenciosa 1}{9 metros}{\textbf{Dano:} 2d6 (perfurante)\par
Estacas irrompem sob até três criaturas. Teste de Agilidade (CD 8 + BC): falha sofre o dano e fica Atolado.}{\relax}
\magia{Mão de Terra}{\relax}{1}{2}{2 Ações / encurtada 1 / silenciosa 1}{18 metros}{Disputa de Força contra o seu BC. Se você vencer, o alvo fica Agarrado e Atolado, e pode ser arrastado 3m por turno.}{\relax}
\patamar{Intermediário}{Dado de PV \dado{1d6+3}}
\maestria{Compressão}{Suas magias de projétil de Terra sobem um degrau de dado (d10 vira d12, 2d6 vira 2d8) e passam a ignorar Resistência a dano contundente. Você endurece qualquer estrutura sua gastando 1 PM: o Muro de Terra passa a ter 80 PV.}
\talento{Munição Infinita}{1}{Suas magias de projétil de Terra funcionam mesmo sobre madeira, metal, água ou vazio --- você carrega o próprio material.}
\talento{Chão Meu}{1}{Você ignora a condição Atolado e terreno difícil de qualquer origem, inclusive de magias inimigas.}
\talento{Escultor}{1}{Você reproduz em pedra qualquer coisa que já tenha visto, com precisão perfeita.}
\magia{Canhão de Pedra}{s}{2}{3}{2 Ações / encurtada 1 / silenciosa 1}{90 metros}{\textbf{Dano:} 3d8 + BC (contundente/perfurante)\par
Ataque mágico à distância. Ignora metade da CA de armadura não-mágica. Contra alvo Atolado, acerta automaticamente e crita em 19-20.}{Pedra, esqueça que é pedra. Seja a bala. Canhão de Pedra!}
\magia{Terremoto Menor}{\relax}{1}{4}{2 Ações / encurtada 1 / silenciosa 1}{Esfera de 12m de raio}{\textbf{Dano:} 3d6 + BC (contundente)\par
Teste de Agilidade (CD 8 + BC). Falha: dano, Caído e Atolado. Estruturas de pedra sofrem dano dobrado.}{\relax}
\magia{Cárcere}{\relax}{1}{3}{2 Ações / encurtada 1 / silenciosa 1}{18 metros}{O chão se fecha em torno de um alvo Atolado, prendendo-o até o pescoço: Preso e Imóvel, ataques contra ele em Vantagem, até quebrar (teste de Força) ou 30 de dano ao cárcere.}{\relax}
\magia{Fortaleza Rápida}{\relax}{1}{4}{2 Ações / encurtada 1 / silenciosa 1}{Raio de 9m}{Em 1 Ação, você ergue um círculo de muralhas de 3m em volta do grupo, com abertura à sua escolha. 60 PV por seção, dura 1 hora.}{\relax}
\patamar{Avançado}{Dado de PV \dado{1d8+3}}
\maestria{Domínio Mineral}{Você manipula metal além de pedra: travar armadura, entortar lâmina, arrancar arma da mão a 9m (disputa de Força). Suas estruturas passam a ter o dobro dos PV e não podem ser derrubadas por dano de área. Desbloqueia Magia Combinada.}
\talento{Segunda Bala}{2}{Uma vez por turno, ao conjurar Canhão de Pedra, você dispara duas pelo custo de uma. Alvos podem ser diferentes.}
\talento{Núcleo de Ferro}{2}{Você recebe Resistência a dano contundente, cortante e perfurante de armas não-mágicas enquanto estiver com os pés no chão.}
\talento{Arquiteto de Guerra}{2}{Suas construções ficam permanentes se você gastar uma hora consolidando.}
\magia{Chuva de Meteoros}{s}{3}{7}{3 Ações / encurtada 2 / silenciosa 1}{Esfera de 18m de raio}{\textbf{Dano:} 8d8 + BC (contundente)\par
Teste de Agilidade (CD 8 + BC), metade se passar. A área vira terreno difícil permanente; quem falhar fica Atolado nos escombros.}{O que está embaixo já esteve em cima. Devolvo o que era seu. Chuva de Meteoros!}
\magia{Prisão de Pedra}{\relax}{2}{6}{3 Ações / encurtada 2 / silenciosa 1}{27 metros}{Um bloco maciço encapsula o alvo: não vê, não ouve, não conjura, não respira (2d10/turno se precisar de ar). O bloco tem 100 PV.}{\relax}
\magia{Lâmina de Aço}{\relax}{2}{5}{3 Ações / encurtada 2 / silenciosa 1}{27 metros}{\textbf{Dano:} 6d8 + BC (perfurante)\par
Dispara metal extraído do solo em linha, ignorando toda CA de armadura metálica.}{\relax}
\magia{Colapso}{\relax}{2}{6}{3 Ações / encurtada 2 / silenciosa 1}{45 metros}{\textbf{Dano:} 8d6\par
Uma estrutura à sua escolha desmorona. Criaturas embaixo fazem teste de Agilidade (CD 8 + BC) ou sofrem o dano e ficam Presas.}{\relax}
\patamar{Santo}{Dado de PV \dado{1d8+4}}
\maestria{O Chão Obedece}{Você molda terra, pedra e metal num raio de 1 km, sem PM, fora de combate. Em combate, gaste 1 Ação para reformular o campo de batalha (aliados escolhem se são afetados; inimigos não). Não pode ser derrubado, empurrado, agarrado nem teleportado contra a vontade enquanto tocar o solo.}
\talento{Peso Absoluto}{3}{Uma vez por combate, sem gastar Ação, triplique o peso de um alvo visível. Teste de Força ou fica Atolado e Caído; voadores caem do céu.}
\magia{Falha Geológica}{s}{4}{11}{4 Ações / encurtada 3 / silenciosa 2}{Linha de 90m × 6m}{\textbf{Dano:} 10d10 de dano de queda\par
Teste de Agilidade com Desvantagem. Falha: cai na fenda e fica Presa no fundo. Sucesso: fica na borda, Atolada e Caída. Estruturas atravessadas desabam; a fenda é permanente.}{Placa que dorme há dez mil anos, acorde por dois segundos. Só dois. Falha Geológica!}
\magia{Golem}{\relax}{3}{9}{4 Ações / encurtada 3 / silenciosa 2}{18 metros}{Anima um corpo de pedra de 3m por 10 minutos: 80 PV, CA 16, ataca com seu BC causando 3d10+BC, obedece ordens simples sem gastar suas Ações. Um por vez.}{\relax}
\patamar{Rei}{Dado de PV \dado{1d10+4}}
\maestria{Metal e Magma}{Metal: você conjura, molda e endurece aço puro --- suas armas de pedra passam a contar como mágicas e ignoram Resistência. Magma: dano de magma ignora Resistência a ígneo e contundente ao mesmo tempo, e aplica Em Chamas mesmo em imunes.}
\talento{Bala Imperial}{4}{Seu Canhão de Pedra passa a rolar 8d8 e ignora toda CA de armadura, Cobertura e Manto de Touki. Em troca, custa 6 PM e não combina com Segunda Bala.}
\magia{Rio de Magma}{s}{5}{14}{5 Ações / encurtada 4 / silenciosa 3}{Linha de 45m × 9m}{\textbf{Dano:} 12d8 + BC de dano de magma no impacto, depois 6d10 por turno\par
A área permanece coberta de magma por 10 minutos. Teste de Agilidade para metade --- criaturas Atoladas não podem testar.}{\relax}
\magia{Muralha do Fim}{\relax}{4}{12}{5 Ações / encurtada 4 / silenciosa 3}{Raio de 90m}{Você ergue uma muralha de pedra e aço de 9 metros de altura, com 400 PV por seção. Permanente até ser destruída.}{\relax}
\patamar{Imperador}{Dado de PV \dado{1d10+5}}
\maestria{O Continente é Seu}{Você molda terra, pedra, metal e magma num raio de 10 km, permanentemente. Criaturas hostis em contato com o solo a até 30m estão permanentemente Atoladas, sem teste. Uma vez por turno, conjure magia de Terra Avançado ou inferior em Silenciosa sem gastar Ação.}
\talento{Aquele que Move Montanhas}{4}{Uma vez por Descanso Longo, conjure qualquer magia de Terra pagando metade do PM, arredondado para baixo.}
\magia{Sepultamento}{s}{6}{20}{6 Ações / silenciosa 4}{Esfera de 45m}{\textbf{Dano:} 16d10 de dano contundente\par
\textbf{Ritual:} não pode ser encurtado.\par
Teste de Força com Desvantagem Absoluta. Falha: soterrada, Presa e sufocando (4d10/turno) até ser escavada. Sucesso: metade e Atolada. Aliados são poupados automaticamente.}{A terra recebe tudo de volta. Eu só estou antecipando a data. Sepultamento!}
\magia{Cordilheira}{\relax}{5}{18}{6 Ações / silenciosa 4}{Raio de 3 km}{\textbf{Ritual:} não pode ser encurtado.\par
Você levanta ou derruba a geografia da região --- montanha, vale, desfiladeiro. Efeito narrativo permanente: rotas mudam, exércitos são forçados a rodear.}{\relax}

\subsection{Magia de Cura}
\begin{arvorecard}{Magia de Cura}{Cura e Suporte}{Espírito \textbullet\ Recurso: PM}
Não cura veneno, doença nem maldição (isso é Desintoxicação) --- mas decide quem sobrevive à campanha. O Rank Deus (Ressurreição de recém-mortos) é puramente narrativo, sem custo em PA.
\end{arvorecard}
\patamar{Principiante}{Dado de PV \dado{1d4+2}}
\maestria{Diagnóstico}{Encostando ou observando por 10 segundos, você sabe PV atual/máximo, condições ativas e a categoria do problema (ferimento, veneno, doença, maldição, exaustão ou fome). Encostando num aliado a 0 PV, você o estabiliza automaticamente, sem PM, Ação nem rolagem.}
\talento{Mãos Firmes}{1}{Você conjura magias de Cura sem sofrer Desvantagem por corpo a corpo, chuva, escuro ou inimigo adjacente.}
\talento{Reserva do Curandeiro}{1}{+2 PM por patamar seu em Cura. Aplicado sozinho na ficha, e cresce a cada patamar novo que você abrir nela.}
\talento{Juramento}{1}{Enquanto jurar nunca usar magia para ferir, suas magias de Cura custam 1 PM a menos (mínimo 1). Quebrar o juramento desliga o talento por uma semana.}
\magia{Cura}{s}{2}{2}{2 Ações / encurtada 1 / silenciosa 1}{Toque}{\textbf{Dano:} 2d8 + BC de PV (4d8 + BC se Ferida Fresca)\par
Remove todas as Marcas da Morte do alvo e o acorda, se estiver a 0 PV.}{Que este poder divino seja alimento farto, e que dê a quem perdeu as forças a força de se erguer de novo. Cura!}
\magia{Estancar}{\relax}{1}{1}{2 Ações / encurtada 1 / silenciosa 1}{Toque}{Encerra imediatamente sangramento, corte aberto ou queimadura ativa. Não funciona contra veneno nem fogo mágico ainda ardendo.}{\relax}
\magia{Selar a Ferida}{\relax}{1}{1}{2 Ações / encurtada 1 / silenciosa 1}{Toque}{Custa 1 PM por Bônus de Rank seu em Cura (1 no Principiante, 6 no Imperador). O dano escolhido conta como Ferida Fresca por 1 hora, mesmo depois de o turno em que ocorreu passar --- qualquer magia de Cura conjurada nesse dano nessa 1 hora ainda cura em dobro.}{\relax}
\magia{Vigor Emprestado}{\relax}{1}{2}{2 Ações / encurtada 1 / silenciosa 1}{Toque}{\textbf{Dano:} 1d8 + BC de PV Temporários\par
Duram 10 minutos, gastos antes dos PV reais, não acumulam com outra fonte.}{\relax}
\magia{Mão que Acalma}{\relax}{1}{2}{2 Ações / encurtada 1 / silenciosa 1}{Toque}{Remove um nível de Exaustão causado por ferimento, trauma ou por acordar do Fio da Vida. Não remove Exaustão por fome, sede ou marcha.}{\relax}
\patamar{Intermediário}{Dado de PV \dado{1d6+2}}
\maestria{Alcance da Compaixão}{Suas magias de Cura passam a alcançar 9 metros. Magias de rank Principiante custam 1 PM a menos (mínimo 1), e você cura duas criaturas ao mesmo tempo com uma conjuração de Principiante, dividindo os dados.}
\talento{Diagnóstico Profundo}{1}{Seu Diagnóstico revela também a causa: quem envenenou, que criatura infectou, há quanto tempo, e se é mágico ou natural.}
\talento{Toque Duplo}{1}{Ao usar Prontidão, você pode curar dois aliados que sofreram dano do mesmo efeito, dividindo os dados.}
\talento{Curandeiro de Guerra}{1}{Você trata quatro pessoas por hora fora de combate, e nunca erra um diagnóstico sob pressão.}
\magia{Prontidão}{s}{2}{3}{1 Reação}{9 metros}{\textbf{Dano:} 2d8 + BC de PV, sempre como Ferida Fresca: 4d8 + BC\par
1 Reação, quando um aliado visível sofrer dano. A magia que define a escola --- cura enquanto o golpe ainda está acontecendo.}{\relax}
\magia{Bênção Coletiva}{\relax}{1}{4}{2 Ações / encurtada 1 / silenciosa 1}{Esfera de 6m}{\textbf{Dano:} 1d8 + BC de PV (dobrado individualmente para quem tiver Ferida Fresca)\par
Todos os aliados na área recuperam PV.}{\relax}
\magia{Escudo de Carne}{\relax}{1}{3}{2 Ações / encurtada 1 / silenciosa 1}{Toque}{1 minuto: o alvo recebe Resistência a dano físico de armas mundanas. Encerra se ele cair a 0 PV.}{\relax}
\magia{Transferência}{\relax}{1}{2}{2 Ações / encurtada 1 / silenciosa 1}{Toque}{Até BC pontos do dano de um aliado adjacente passam para você, ignorando sua Resistência. O dano transferido conta como Ferida Fresca em você.}{\relax}
\magia{Sono Reparador}{\relax}{1}{3}{2 Ações / encurtada 1 / silenciosa 1}{Toque}{O alvo dorme 1 hora e acorda como se tivesse feito um Descanso Curto completo. Não funciona duas vezes na mesma pessoa entre Descansos Longos.}{\relax}
\patamar{Avançado}{Dado de PV \dado{1d6+2}}
\maestria{A Ferida Mortal}{Você cura ferimentos potencialmente fatais: pode levar um aliado de 0 PV direto para PV positivos numa única conjuração, sem que ele receba Exaustão ao acordar. Não reimplanta membros decepados. Desbloqueia Magia Combinada.}
\talento{Mão Silenciosa}{2}{Você conjura magias de Cura de rank Principiante e Intermediário em Conjuração Silenciosa, sem penalidade alguma.}
\talento{Sangue Trocado}{2}{Transferência passa a mover o dobro do dano e alcança 9 metros.}
\talento{Mãos Repartidas}{2}{Ao conjurar Cura Suprema, você pode dividir os dados rolados entre até duas criaturas ao seu alcance, em vez de concentrar tudo em uma só. Cada uma ainda dobra individualmente se estiver com Ferida Fresca.}
\magia{Cura Suprema}{s}{3}{6}{3 Ações / encurtada 2 / silenciosa 1}{9 metros}{\textbf{Dano:} 6d8 + BC de PV (12d8 + BC se Ferida Fresca)\par
Remove sangramentos, ossos quebrados e a condição Caído. Se o alvo estiver a 0 PV, levanta com metade dos PV máximos, sem Exaustão.}{\relax}
\magia{Círculo de Recuperação}{\relax}{2}{5}{3 Ações / encurtada 2 / silenciosa 1}{Esfera de 9m}{\textbf{Dano:} 2d6 + BC de PV por turno\par
3 turnos: todo aliado que começar o turno dentro da área recupera PV. Sustentada sem concentração, um círculo por vez.}{\relax}
\magia{Rejeitar a Morte}{\relax}{2}{5}{1 Reação}{9 metros}{1 Reação, quando um aliado visível chegaria a 0 PV: ele fica com 1 PV e pode se mover 4,5m imediatamente. Uma vez por criatura por combate.}{\relax}
\magia{Anestesia}{\relax}{2}{3}{3 Ações / encurtada 2 / silenciosa 1}{Toque}{10 minutos sem dor: ignora penalidades de ferimento/Exaustão, mas o Mestre para de informar os PV do alvo ao jogador.}{\relax}
\patamar{Santo}{Dado de PV \dado{1d6+3}}
\maestria{A Luz que Reconecta}{Você reimplanta membros recém-decepados (menos de 1h, ou a qualquer momento se a ferida estiver Selada). Não recria o que não existe mais. Magias de rank Avançado ou inferior em Conjuração Silenciosa sem penalidade. Mantém duas magias sustentadas.}
\talento{Vigília}{3}{Uma vez por combate, sua Prontidão não gasta Reação.}
\magia{Cura Radiante}{s}{4}{10}{4 Ações / encurtada 3 / silenciosa 2}{18 metros}{\textbf{Dano:} 10d8 + BC de PV (20d8 + BC se Ferida Fresca)\par
Reimplanta um membro recém-decepado. Remove todas as Marcas da Morte, Exaustão física, e Caído/Atordoado/Cego de origem traumática.}{Anjo dos milagres, conceda teu sopro sagrado ao coração que pulsa diante de ti. Cura Radiante!}
\magia{Santuário Menor}{\relax}{3}{8}{4 Ações / encurtada 3 / silenciosa 2}{Esfera de 9m}{1 minuto: nenhum aliado na área recebe Marcas da Morte, e todo aliado que chegaria a 0 PV fica com 1 PV --- uma vez cada. Cada uma dessas salvações conta como a Salvação daquela criatura no combate (Cap. 4, §4, Uma Salvação por Combate).}{\relax}
\magia{Corpo de Ferro}{\relax}{3}{7}{4 Ações / encurtada 3 / silenciosa 2}{Toque}{10 minutos: +50 PV máximos e imunidade a acertos críticos.}{\relax}
\patamar{Rei}{Dado de PV \dado{1d8+3}}
\maestria{Golpe Divino}{Você canaliza cura ofensivamente contra mortos-vivos, construtos e criaturas de mana corrompida: sofrem o valor curado como dano radiante, sem teste. Contra vivos normais, não funciona. Regenera membros perdidos desde que os ossos estejam disponíveis, não importa há quanto tempo.}
\talento{Sopro do Julgamento}{4}{Seu Golpe Divino passa a funcionar também contra demônios de linhagem antiga e criaturas de rank Deus corrompidas, com metade do dano.}
\magia{Restauração}{s}{5}{15}{5 Ações / encurtada 4 / silenciosa 3}{Toque}{\textbf{Ritual:} não pode ser encurtado.\par
Ritual de 10 minutos. Regenera completamente um membro/olho/orelha/órgão perdido, desde que você tenha os ossos correspondentes. Não funciona contra perdas por magia Imperador ou superior.}{\relax}
\magia{Julgamento}{\relax}{4}{12}{5 Ações / encurtada 4 / silenciosa 3}{Esfera de 12m}{\textbf{Dano:} 10d8 + BC de dano radiante (mortos-vivos/construtos/corrompidos, sem teste)\par
Aliados vivos na área recuperam 3d8 + BC de PV na mesma conjuração.}{\relax}
\magia{Milagre Menor}{\relax}{4}{12}{5 Ações / encurtada 4 / silenciosa 3}{Toque}{Remove uma condição de qualquer origem, inclusive Paralisia, Petrificação, Congelado e cegueira mágica. Exceção absoluta: veneno, doença e maldição continuam fora de alcance.}{\relax}
\patamar{Imperador}{Dado de PV \dado{1d8+4}}
\maestria{Nada é Irreversível}{Você regenera membros perdidos independentemente de quanto tempo faz e sem precisar dos ossos. Suas magias de Cura curam o valor máximo dos dados contra Feridas Frescas, sem rolar. Uma vez por turno, conjure magia de Cura de rank Avançado ou inferior em Silenciosa sem gastar Ação.}
\talento{A Mão que Não Cansa}{4}{Uma vez por Descanso Longo, conjure qualquer magia de Cura pagando zero PM.}
\magia{Corpo Íntegro}{s}{6}{25}{6 Ações / silenciosa 4}{Toque}{\textbf{Ritual:} não pode ser encurtado.\par
Ritual de 1 hora. O alvo é restaurado à integridade física completa: membros, órgãos, sentidos e cicatrizes voltam ao estado original, não importa quando/como foram perdidos. PV restaurados, Exaustão removida. Não remove veneno, doença nem maldição; não funciona em quem foi pulverizado, petrificado permanentemente ou consumido por magia Imperador ou superior.}{\relax}
\magia{Santuário}{\relax}{5}{20}{4 Ações}{Esfera de 30m}{\textbf{Ritual:} não pode ser encurtado.\par
Ritual (4 Ações). 1 minuto: nenhuma criatura viva na área pode morrer --- ainda caem a 0 PV, mas não recebem Marcas da Morte e não podem ser mortas por efeito abaixo de rank Deus. Quando acaba, todo o dano acumulado é aplicado de uma vez. Você não pode conjurar outra magia enquanto ativo.}{\relax}
\magia{Luz Absoluta}{\relax}{5}{22}{6 Ações / silenciosa 4}{Esfera de 30m}{\textbf{Dano:} 20d8 de dano radiante\par
Mortos-vivos e construtos com menos da metade dos PV máximos são destruídos automaticamente; os demais sofrem o dano. Todo aliado vivo na área recupera 10d8 + BC de PV.}{\relax}

\subsection{Magia de Desintoxicação}
\begin{arvorecard}{Magia de Desintoxicação}{Cura e Suporte}{Espírito \textbullet\ Recurso: PM}
Trata veneno, doença, maldição e petrificação --- a única escola cuja dificuldade cresce sozinha (Profundidade). O Rank Deus (a Doença da Pedra Mágica) é puramente narrativo.
\end{arvorecard}
\patamar{Principiante}{Dado de PV \dado{1d6+2}}
\maestria{Paladar}{Tocando, cheirando ou provando qualquer substância, você sabe exatamente o que ela é, e identifica a Profundidade de qualquer aflição que veja, inclusive em cadáveres. Você é imune a veneno mundano.}
\talento{Reserva do Purificador}{1}{+2 PM por patamar seu em Desintoxicação. Aplicado sozinho na ficha, e cresce a cada patamar novo que você abrir nela.}
\talento{Herborista}{1}{Fora de combate, com uma hora e material, reduza a Profundidade de uma aflição em 1 sem gastar PM. No máximo uma vez por aflição por Descanso Longo.}
\talento{Mão que Não Contamina}{1}{Você não pode ser envenenado, infectado ou amaldiçoado por contato ao manusear aquilo que está tratando ou extraindo.}
\magia{Purgar}{s}{2}{2}{2 Ações / encurtada 1 / silenciosa 1}{Toque}{Remove uma aflição de Profundidade 1 do alvo --- veneno, doença, maldição ou petrificação incipiente, mágica ou não.}{O que entrou sem ser convidado, saia do jeito que quiser, mas saia. Purgar!}
\magia{Antídoto}{\relax}{1}{1}{2 Ações / encurtada 1 / silenciosa 1}{Toque}{O alvo recebe Vantagem contra veneno e doença por 1 hora; se já afetado, a Profundidade para de subir durante o período.}{\relax}
\magia{Água Limpa}{\relax}{1}{1}{2 Ações / encurtada 1 / silenciosa 1}{Toque}{Purifica até 20 litros de comida, água, ar ou terreno contaminado.}{\relax}
\magia{Sangria}{\relax}{1}{2}{2 Ações / encurtada 1 / silenciosa 1}{Toque}{\textbf{Dano:} 2d6 de dano ao alvo\par
Reduz em 1 a Profundidade de uma aflição sem removê-la, ao custo do dano acima. No máximo uma vez por aflição por Descanso Longo.}{\relax}
\magia{Estômago de Ferro}{\relax}{1}{1}{2 Ações / encurtada 1 / silenciosa 1}{Toque}{Por 8 horas, o alvo pode comer e beber qualquer coisa sem consequência.}{\relax}
\patamar{Intermediário}{Dado de PV \dado{1d6+2}}
\maestria{Extração}{Ao purgar qualquer aflição, você pode capturá-la num frasco em vez de dissipá-la (mantém a Profundidade original, dura 1 mês). Aplicada em arma, comida ou superfície, força teste de Vigor (CD 8 + BC) na próxima criatura exposta. Você carrega até Espírito frascos.}
\talento{Frasco Estável}{1}{Suas extrações duram um ano em vez de um mês, e você carrega o dobro de frascos.}
\talento{Purga Coletiva}{1}{O Purgar de rank Principiante passa a atingir até três criaturas adjacentes com uma conjuração.}
\talento{Leitura de Sintoma}{1}{Você sabe se alguém está afetado por algo antes dos sintomas aparecerem, incluindo maldições dormentes e venenos de efeito retardado.}
\magia{Purga Profunda}{s}{2}{4}{2 Ações / encurtada 1 / silenciosa 1}{9 metros}{Remove uma aflição de Profundidade 2 ou menor, ou reduz em 2 a Profundidade de qualquer aflição --- a redução, no máximo uma vez por aflição por Descanso Longo.}{\relax}
\magia{Muro Estéril}{\relax}{1}{3}{2 Ações / encurtada 1 / silenciosa 1}{Esfera de 9m}{Por 1 hora, gás venenoso, esporo, praga, ácido e ar contaminado não entram na área.}{\relax}
\magia{Torpor}{\relax}{1}{3}{2 Ações / encurtada 1 / silenciosa 1}{18 metros}{Teste de Vigor (CD 8 + BC) ou o alvo fica Envenenado por 1 minuto (Desvantagem em ataques e testes de atributo). Sem dano.}{\relax}
\magia{Diagnóstico de Praga}{\relax}{1}{2}{2 Ações / encurtada 1 / silenciosa 1}{18 metros}{Identifica todas as criaturas doentes, envenenadas ou amaldiçoadas num raio de 18m, a Profundidade de cada uma, e quem foi a origem.}{\relax}
\patamar{Avançado}{Dado de PV \dado{1d6+3}}
\maestria{Contra a Maré}{Você purga aflições de Profundidade 3 ou menor. A Profundidade de qualquer aflição para de subir em criaturas a até 9 metros de você, enquanto consciente. Desbloqueia Magia Combinada.}
\talento{Purificador de Guerra}{2}{Muro Estéril e Quarentena passam a cobrir o dobro da área e a durar o dobro do tempo.}
\talento{Extração Refinada}{2}{Aflições extraídas por você sobem 1 de Profundidade ao serem aplicadas em outra criatura.}
\talento{Corpo Recusado}{2}{Você é imune a veneno e doença não-mágicos, e a Profundidade de qualquer aflição sobe em você na metade da velocidade normal.}
\magia{Anular}{s}{3}{6}{3 Ações / encurtada 2 / silenciosa 1}{9 metros}{Remove uma condição de qualquer origem do alvo: Envenenado, Paralisado, Petrificado, Cego, Surdo, Amedrontado, Atordoado, Congelado, Em Chamas, Atolado, Desequilibrado ou Marcado. Também remove aflições de Profundidade 3 ou menor.}{\relax}
\magia{Quarentena}{\relax}{2}{5}{3 Ações / encurtada 2 / silenciosa 1}{Esfera de 12m}{Por 10 minutos, nada tóxico, infeccioso ou amaldiçoado atravessa a borda da área, nos dois sentidos.}{\relax}
\magia{Corrosão}{\relax}{2}{5}{3 Ações / encurtada 2 / silenciosa 1}{18 metros}{\textbf{Dano:} 4d8 de dano ácido (dobrado contra construtos e armaduras pesadas)\par
Teste de Vigor (CD 8 + BC). Falha: dano e Envenenado por 1 minuto. Metal não-mágico exposto perde 2 de CA permanentemente.}{\relax}
\magia{Sangue Trocado}{\relax}{2}{4}{3 Ações / encurtada 2 / silenciosa 1}{Toque}{Você transfere uma aflição de um alvo para você mesmo, reduzindo a Profundidade dela em 2 no processo.}{\relax}
\patamar{Santo}{Dado de PV \dado{1d8+3}}
\maestria{Estado Anulado}{Você purga aflições de Profundidade 4 ou menor. Toda criatura que você purgar fica imune àquela aflição específica por 24 horas. Você anula qualquer condição (dos outros) com um toque, sem PM, uma vez por turno.}
\talento{A Mão que Não Erra}{3}{Uma vez por Descanso Longo, você purga uma aflição de Profundidade um ponto acima do seu limite.}
\magia{Purificação}{s}{4}{11}{4 Ações / encurtada 3 / silenciosa 2}{Esfera de 30m}{\textbf{Ritual:} não pode ser encurtado.\par
Toda criatura, água, solo, alimento e estrutura na área é purgada de aflições de Profundidade 4 ou menor.}{Que a terra esqueça o que foi despejado nela. Que a carne esqueça o que entrou nela. Purificação!}
\magia{Selar a Maldição}{\relax}{3}{9}{4 Ações / encurtada 3 / silenciosa 2}{Toque}{Uma aflição que você não consegue purgar fica congelada em Profundidade e sintomas por até um ano --- não melhora, mas para de piorar.}{\relax}
\patamar{Rei}{Dado de PV \dado{1d8+3}}
\maestria{Veneno}{Você conjura aflições diretamente, sem frasco nem contato. Venenos criados por você começam na Profundidade igual à metade do seu Bônus de Rank, arredondado para cima. Seus efeitos de veneno ignoram Resistência a veneno. Você purga aflições de Profundidade 5 ou menor.}
\talento{Duas Faces}{4}{Quando purgar uma aflição de uma criatura, gaste 1 Ação para aplicá-la imediatamente em outra criatura visível a até 9m, sem frasco e sem teste.}
\magia{Sopro Podre}{s}{5}{13}{5 Ações / encurtada 4 / silenciosa 3}{Cone de 18m}{\textbf{Dano:} 10d8 de dano de veneno\par
Teste de Vigor com Desvantagem. Falha: dano, Envenenado por 10 minutos, e uma aflição de Profundidade 3 que continua subindo. Sucesso: metade. Não funciona em construtos, mortos-vivos ou quem não respira.}{\relax}
\magia{Toque do Fim}{\relax}{4}{10}{5 Ações / encurtada 4 / silenciosa 3}{Toque}{Teste de Vigor (CD 8 + BC). Falha: o alvo recebe uma aflição de Profundidade 4 à sua escolha, que sobe 1 por dia e mata ao chegar a 6. Só um mago de patamar igual ou superior consegue removê-la.}{\relax}
\patamar{Imperador}{Dado de PV \dado{1d8+4}}
\maestria{O Corpo Limpo}{Você purga qualquer aflição de Profundidade 5 ou menor sem rolagem, ritual ou tempo --- inclusive petrificação completa, maldições hereditárias e parasitas mágicos. É permanentemente imune a veneno, doença, maldição e petrificação. Uma vez por turno, conjure magia de Desintoxicação de rank Avançado ou inferior em Silenciosa sem gastar Ação.}
\talento{Nada Entra}{4}{Todos os aliados a até 18 metros de você compartilham a sua imunidade a veneno, doença e maldição enquanto permanecerem ao seu alcance.}
\magia{O Mundo Sem Praga}{s}{6}{22}{6 Ações / silenciosa 4}{Raio de 3 km}{\textbf{Ritual:} não pode ser encurtado.\par
Toda aflição de Profundidade 5 ou menor deixa de existir dentro do raio: em pessoas, água, solo, ar e paredes. Efeito narrativo permanente.}{\relax}
\magia{Nome do Veneno}{\relax}{5}{18}{6 Ações / silenciosa 4}{45 metros}{Você transfere uma aflição inteira, com Profundidade intacta, de uma criatura afetada para outra visível, sem teste (com Desvantagem contra criaturas de rank Deus).}{\relax}

\subsection{Barreira e Proteção}
\begin{arvorecard}{Barreira e Proteção}{Cura e Suporte}{Espírito \textbullet\ Recurso: PM}
Anti-magia medida em regras, não em dano. Fraqueza estrutural: barreiras distorcem mana, e aço não é mana --- quase nada contra o pilar do Corpo.
\end{arvorecard}
\patamar{Principiante}{Dado de PV \dado{1d6+2}}
\maestria{O Primeiro Círculo}{Você desenha barreiras (1 Ação + PM da magia), esferas centradas num ponto à sua escolha. Sustenta uma por vez, aplicando Selado e a condição de Fluxo Interrompido que você escolher. Você vê mana: barreiras, encantamentos, itens mágicos e invisibilidade mágica aparecem como contorno luminoso, sem custo.}
\talento{Reserva do Selador}{1}{+2 PM por patamar seu em Barreira. Aplicado sozinho na ficha, e cresce a cada patamar novo que você abrir nela.}
\talento{Círculo Portátil}{1}{Sua barreira passa a se mover com você, centrada no seu corpo, em vez de ficar fixa num ponto.}
\talento{Mão de Giz}{1}{Você desenha círculos permanentes em superfícies. Leva 1 hora e o dobro do PM, mas a barreira fica lá depois que você for embora.}
\magia{Círculo Menor}{s}{2}{3}{2 Ações / encurtada 1 / silenciosa 1}{Esfera de 6m}{Barreira de 6m de raio por 1 minuto. Criaturas dentro ficam Seladas. Você escolhe Estagnação ou Fonte.}{Aqui dentro a regra é outra. Eu escrevi a regra. Círculo Menor!}
\magia{Recusa}{\relax}{1}{2}{1 Reação}{18 metros}{1 Reação, quando uma criatura conjurar magia de rank Principiante: a magia falha e o PM se perde. Contra rank Intermediário, ainda acontece, mas com metade dos dados.}{\relax}
\magia{Selo de Objeto}{\relax}{1}{1}{2 Ações / encurtada 1 / silenciosa 1}{Toque}{Um objeto mágico, arma encantada ou item amaldiçoado fica inerte por 1 hora.}{\relax}
\magia{Anteparo}{\relax}{1}{2}{2 Ações / encurtada 1 / silenciosa 1}{9 metros}{Uma placa de mana de 3m × 3m por 3 turnos. Qualquer magia de rank Principiante que a atravesse é anulada; de rank Intermediário tem os dados reduzidos à metade.}{\relax}
\magia{Leitura de Trama}{\relax}{1}{1}{2 Ações / encurtada 1 / silenciosa 1}{18 metros}{Você identifica exatamente qual magia está ativa numa criatura/objeto/local, de qual escola e rank, e quanto tempo falta. Revela armadilhas mágicas e barreiras alheias.}{\relax}
\patamar{Intermediário}{Dado de PV \dado{1d6+3}}
\maestria{Interdição}{Magias conjuradas dentro de uma barreira sua custam +2 PM para quem não for seu aliado. Você sustenta duas barreiras simultaneamente, e pode escolher até Espírito criaturas para quem a barreira simplesmente não se aplica.}
\talento{Trama Fina}{1}{Sua Recusa passa a anular magias de rank Intermediário por completo.}
\talento{Barreira Persistente}{1}{Suas barreiras continuam de pé por 1 minuto depois de você ficar inconsciente ou sair do alcance.}
\talento{Peneira}{1}{Você declara uma escola de magia à qual a sua barreira não se aplica.}
\magia{Domo}{s}{2}{5}{2 Ações / encurtada 1 / silenciosa 1}{Esfera de 12m}{Barreira de 12m por 10 minutos. Além de Selado e Fluxo Interrompido, impede passagem de efeitos mágicos pela superfície nos dois sentidos. Criaturas e flechas atravessam normalmente.}{\relax}
\magia{Amarra}{\relax}{1}{4}{2 Ações / encurtada 1 / silenciosa 1}{18 metros}{Teste de Espírito (CD 8 + BC). Falha: por 1 minuto, o alvo não pode conjurar magia alguma. Ele ainda pode andar, correr e bater.}{\relax}
\magia{Espelho de Mana}{\relax}{1}{4}{1 Reação}{Pessoal}{1 Reação. A próxima magia de alvo único de rank Intermediário ou inferior dirigida a você é devolvida ao conjurador, com a CD original dele.}{\relax}
\magia{Silêncio de Mana}{\relax}{1}{3}{2 Ações / encurtada 1 / silenciosa 1}{Esfera de 9m}{Por 1 minuto, criaturas na área não conseguem iniciar conjuração --- quem já estava recitando pode terminar.}{\relax}
\patamar{Avançado}{Dado de PV \dado{1d8+3}}
\maestria{Selo de Conjuração}{O Selado das suas barreiras impõe Desvantagem em testes de resistência contra as suas magias para quem está dentro. Criaturas dentro que tentarem conjurar acima do limite sofrem 2d10 de dano psíquico e perdem a Ação. Desbloqueia Magia Combinada.}
\talento{Trama Densa}{2}{Todas as suas barreiras ganham o dobro de PV.}
\talento{Selo Cirúrgico}{2}{Amarra deixa de permitir teste de resistência contra criaturas de rank inferior ao seu.}
\talento{Duas Mãos, Três Círculos}{2}{Você sustenta três barreiras simultaneamente.}
\magia{Recinto}{s}{3}{7}{3 Ações / encurtada 2 / silenciosa 1}{Esfera de 18m}{Barreira de 18m por 1 hora. Impede entrada e saída física de criaturas: a superfície vira sólida para carne, mas continua atravessável por objetos inanimados. Tem 120 PV contra ataques físicos.}{\relax}
\magia{Dissipar}{\relax}{2}{6}{3 Ações / encurtada 2 / silenciosa 1}{27 metros}{Encerre um efeito mágico ativo de rank Avançado ou inferior: magia sustentada, barreira alheia, encantamento, invocação, condição de origem mágica. Sem teste, sem disputa.}{\relax}
\magia{Campo Nulo}{\relax}{2}{6}{3 Ações / encurtada 2 / silenciosa 1}{Esfera de 12m}{Por 3 turnos, nenhuma magia funciona dentro da área, incluindo as suas. Sustentadas de fora são suspensas; invocações desaparecem; itens mágicos ficam inertes.}{\relax}
\magia{Redoma}{\relax}{2}{5}{3 Ações / encurtada 2 / silenciosa 1}{9 metros}{Uma esfera de 1,5m encapsula uma criatura: totalmente isolada, não conjura. Tem 60 PV.}{\relax}
\patamar{Santo}{Dado de PV \dado{1d8+3}}
\maestria{Espaço Recusado}{Suas barreiras impedem teleporte, invocação, passagem dimensional e qualquer forma de aparecer/desaparecer na área. Você sustenta barreiras sem limite de distância, no mesmo continente. Você vê e lê qualquer barreira, encantamento ou selo do mundo apenas olhando.}
\talento{O Selo Não Cede}{3}{Suas barreiras não podem ser destruídas por Dissipar/Anulação de patamar inferior ao seu, e criaturas presas dentro não saem por meios mágicos de nenhum patamar.}
\magia{Interdito}{s}{4}{12}{4 Ações / encurtada 3 / silenciosa 2}{Esfera de 45m}{\textbf{Ritual:} não pode ser encurtado.\par
Por 24 horas: ninguém conjura acima do rank Santo, ninguém recupera PM, ninguém teleporta, nenhuma invocação existe, e item mágico de rank Rei ou inferior fica inerte. Aliados designados ficam isentos.}{\relax}
\magia{Recusar o Mundo}{\relax}{3}{10}{1 Reação}{Pessoal}{1 Reação. Anule completamente uma magia dirigida a você ou aliado a 18m, de rank Rei ou inferior, sem teste. Uma vez por combate. Contra Imperador, dano reduzido à metade sem condições.}{\relax}
\patamar{Rei}{Dado de PV \dado{1d8+3}}
\maestria{Anulação}{Gastando 1 Reação e 4 PM, você anula qualquer magia de rank Imperador ou inferior no instante em que é conjurada, dentro de 45 metros, sem teste. Um número de vezes por combate igual ao seu Espírito. O conjurador perde o PM e as Ações gastas.}
\talento{Retorno}{4}{Quando você anular uma magia com Anulação, o conjurador sofre dano psíquico igual ao PM que gastou, e não pode reconjurar até o fim do combate.}
\magia{O Círculo do Rei}{s}{5}{16}{5 Ações / encurtada 4 / silenciosa 3}{Esfera de 300m}{\textbf{Ritual:} não pode ser encurtado.\par
Por uma semana, a região obedece a três regras à sua escolha: nenhuma magia acima de um rank; ninguém recupera PM; nada teleporta/invoca/atravessa; um tipo de criatura não entra; nada morre aqui.}{\relax}
\magia{Prisão Absoluta}{\relax}{4}{14}{5 Ações / encurtada 4 / silenciosa 3}{45 metros}{Teste de Espírito com Desvantagem (CD 8 + BC). Falha: selada numa redoma de 3m por 1 hora, isolada e sem agir. Criaturas de patamar igual ou superior ao seu repetem o teste ao fim de cada um dos seus turnos. Só libertável por você ou por Dissipar de rank Imperador.}{\relax}
\patamar{Imperador}{Dado de PV \dado{1d8+4}}
\maestria{Lei Local}{O Selado das suas barreiras passa a valer contra todo rank, incluindo Deus. Ao erguer qualquer barreira, declare uma regra arbitrária proibitiva que passa a valer dentro dela. Uma vez por turno, erga ou desfaça uma barreira sem gastar Ação.}
\talento{Barreira Viva}{4}{Suas barreiras persistem mesmo depois da sua morte, até serem dissipadas por um Imperador ou por rank Deus.}
\magia{Mundo Fechado}{s}{6}{24}{6 Ações / silenciosa 4}{Esfera de 1,5 km}{\textbf{Ritual:} não pode ser encurtado.\par
Por 1 hora, dentro da esfera nenhuma magia de nenhum rank funciona para ninguém além de você e quem você designar. Invocações se desfazem, itens mágicos morrem, barreiras alheias caem, voo mágico cessa.}{Eu não apago a magia do mundo. Eu recorto um pedaço do mundo onde ela nunca foi convidada. Mundo Fechado!}
\magia{Selo do Nome}{\relax}{5}{20}{6 Ações / silenciosa 4}{Toque}{Você sela permanentemente uma única magia, técnica ou habilidade de uma criatura tocada. Ela nunca mais consegue usar aquilo, até um mago de Barreira de patamar igual ou superior desfazer.}{\relax}

\subsection{Espíritos e Feras}
\begin{arvorecard}{Espíritos e Feras}{Invocação}{Espírito \textbullet\ Recurso: PM}
A menor lista de feitiços do mundo --- o que você compra aqui não são feitiços, são relações (Pactos).
\end{arvorecard}
\patamar{Principiante}{Dado de PV \dado{1d6+2}}
\maestria{O Primeiro Círculo}{Você desenha círculos e mantém 1 Pacto ativo por vez. Um invocado permanece por 1 hora ou até 0 PV (não pode ser chamado de novo até o próximo Descanso Longo). Invocados têm PV = 10 × seu Bônus de Rank e usam o seu BC para acertar e para CD. Você entende e é entendido por qualquer criatura com Pacto, sem idioma comum.}
\talento{Pacto: Lobo Cinzento}{1}{3d8 mordida, deslocamento 12m, rastreia por cheiro com Vantagem. Se outro aliado estiver adjacente ao alvo, derruba.}
\talento{Pacto: Corvo Mensageiro}{1}{Frágil (metade dos PV), voa 18m. Você vê e ouve pelo que ele vê e ouve a qualquer distância. Não luta.}
\talento{Pacto: Salamandra}{1}{2d8 mordida + 2d6 ígneo, imune a fogo, aplica Em Chamas. Acende fogueiras e derrete fechaduras; do tamanho de um gato.}
\talento{Pacto: Espírito de Pedra}{1}{3d6 soco, Resistência a dano físico, Deslocamento 6m. Não recua nunca. Serve para segurar uma porta.}
\talento{Reserva do Invocador}{1}{+2 PM por patamar seu em Invocação. Aplicado sozinho na ficha, e cresce a cada patamar novo que você abrir nela.}
\talento{Traço Rápido}{1}{Desenhar um círculo em combate passa a custar o PM normal, sem o dobro.}
\talento{Círculo Guardado}{1}{Você carrega um círculo pré-desenhado. Usá-lo dispensa o tempo de desenho, e ele aguenta três invocações antes de se apagar.}
\magia{Chamado}{s}{2}{3}{1 Ações}{Círculo}{Invoca uma criatura com quem você tenha Pacto. Ela surge no círculo e age a partir do próximo turno. Invocar um Pacto acima do Principiante custa PM adicional: +3 PM por patamar de diferença (Intermediário 6, Avançado 9, Santo 12, Rei 15). Desenhar o círculo custa 10 minutos fora de combate, ou 1 Ação e o dobro do PM no meio da luta.}{Eu desenhei o caminho e paguei a passagem. Venha, e o que for combinado será cumprido. Chamado!}
\patamar{Intermediário}{Dado de PV \dado{1d6+2}}
\maestria{Vínculo}{2 Pactos ativos ao mesmo tempo. Você vê pelos olhos de qualquer invocado seu, gastando 1 Ação, a qualquer distância (seu corpo fica Cego enquanto isso). Ordens específicas passam a custar 1 Ação para todos os seus invocados de uma vez.}
\talento{Pacto: Urso das Cavernas}{1}{4d10 garra, Grande, empurra 3m a cada acerto. Absorve dano por você.}
\talento{Pacto: Serpente de Névoa}{1}{Ataque com veneno (Vigor ou Envenenado). Move-se por qualquer fresta; invisível em terreno enevoado.}
\talento{Pacto: Espírito do Vento}{1}{Não ataca. Concede Voo (18m) a um aliado que ele toque, enquanto durar a invocação.}
\talento{Pacto: Grifo}{1}{Voa 24m, 3d10 garra, e carrega uma pessoa.}
\talento{Sangue no Círculo}{1}{Você pode pagar o PM de uma invocação com PV, na razão de 2 PV por 1 PM.}
\talento{Vontade Firme}{1}{Invocados seus são imunes a Amedrontado e não podem ser dominados, expulsos ou dissipados por efeito de patamar inferior ao do próprio Pacto invocado.}
\talento{Companhia}{1}{Um invocado à sua escolha permanece 8 horas em vez de 1.}
\magia{Retorno}{s}{2}{2}{1 Ações}{90 metros}{Um invocado seu volta imediatamente para o círculo, ou é dispensado. Se dispensado antes de chegar a 0 PV, pode ser chamado de novo neste mesmo dia.}{\relax}
\patamar{Avançado}{Dado de PV \dado{1d6+3}}
\maestria{Círculo Rápido}{3 Pactos ativos. Invocar deixa de exigir círculo desenhado: gaste 1 Ação e trace no ar. Invocados passam a ter PV = 15 × seu Bônus de Rank e recebem seu Bônus de Rank no dano. Desbloqueia Magia Combinada.}
\talento{Pacto: Quimera}{2}{Três cabeças, três ataques por turno de 2d8 cada, um deles com sopro elemental à sua escolha ao invocar.}
\talento{Pacto: Espírito Antigo}{2}{Não luta. Responde uma pergunta por invocação sobre algo que aconteceu antes de você nascer --- verdadeira, ainda que enviesada.}
\talento{Pacto: Golem de Guerra}{2}{Enorme, 6d8 por golpe, Resistência a todo dano físico, Deslocamento 6m. Não pode ser movido nem derrubado.}
\talento{Pacto: Alcateia}{2}{Invoca cinco lobos de uma vez, cada um com um quarto dos PV normais e 1d8 de mordida em vez de 3d8, agindo na mesma Iniciativa.}
\talento{Duas Vidas}{2}{Um invocado reduzido a 0 PV pode ser chamado de novo após um Descanso Curto, em vez de Longo.}
\talento{Empatia Absoluta}{2}{Você compartilha os sentidos de todos os seus invocados ao mesmo tempo, sem gastar Ação e sem ficar cego.}
\talento{Pacto Emprestado}{2}{Um aliado pode comandar um dos seus invocados com as próprias Ações, sem que você gaste nada.}
\magia{Substituição}{s}{3}{6}{1 Reação}{18 metros}{1 Reação, quando você ou um aliado for alvo de um ataque: um invocado seu troca de lugar com o alvo instantaneamente e recebe o ataque no lugar dele.}{\relax}
\patamar{Santo}{Dado de PV \dado{1d8+3}}
\maestria{Corpos Artificiais}{4 Pactos ativos. Você constrói o corpo de um invocado à sua escolha entre sessões: escolha duas melhorias permanentes (+50\% PV, +1 ataque por turno, sentido especial, resistência a um tipo de dano, ou voo). Invocados não desaparecem a 0 PV --- o corpo quebra, a consciência volta e recebe um corpo novo após um Descanso Longo.}
\talento{Pacto: Espírito da Chama Antiga}{3}{4d8 de dano ígneo em área de 6m por turno, sem gastar suas Ações. Não obedece bem: role Espírito (CD 15) a cada turno, ou ele escolhe o próprio alvo.}
\talento{Pacto: Sentinela de Aço}{3}{Enorme, 100 PV além do normal, intercepta um ataque por rodada contra um aliado adjacente, automaticamente. Não ataca.}
\talento{Pacto: Mensageiro do Alto}{3}{Atravessa continentes numa hora, entrega qualquer coisa a qualquer pessoa que você já tenha visto, e volta.}
\talento{O Nome Verdadeiro}{3}{Você aprende o nome verdadeiro de um invocado. Ele passa a obedecer qualquer ordem sem teste, inclusive ordens suicidas --- e passa a te odiar em silêncio.}
\magia{Círculo de Convocação}{s}{4}{12}{4 Ações}{Círculo de 12m}{\textbf{Ritual:} não pode ser encurtado.\par
Invoca todos os seus Pactos de uma vez, no mesmo turno. Uma vez por Descanso Longo.}{\relax}
\patamar{Rei}{Dado de PV \dado{1d8+4}}
\maestria{Círculo Permanente}{5 Pactos ativos. Um invocado à sua escolha fica permanentemente ao seu lado, sem custo de PM e sem limite de tempo --- ele come, dorme e tem opinião. Círculos que você desenhar em pedra ou metal duram para sempre e podem ser usados por outras pessoas, se souberem o nome certo.}
\talento{Pacto: Fera Ancestral}{4}{Gigantesca. 4d10 por golpe, três golpes por turno, voa, e respira um elemento à sua escolha em cone de 18m. Ela concorda em vir; não concorda em ficar.}
\talento{Pacto: Espírito do Contrato}{4}{Não luta e não pode ser ferido. Enquanto existir, qualquer acordo verbal na presença dele é vinculante: quem quebrar sofre 10d10 de dano psíquico, onde estiver.}
\talento{Legião}{4}{Ao invocar, você pode chamar três cópias de um mesmo Pacto de patamar Avançado ou inferior, pagando o PM uma vez.}
\magia{Troca de Lugares}{s}{5}{12}{1 Ações}{Alcance ilimitado}{Você e um invocado seu trocam de posição instantaneamente, em qualquer lugar do mundo.}{\relax}
\patamar{Imperador}{Dado de PV \dado{1d10+4}}
\maestria{O Grande Círculo}{6 Pactos ativos, todos podem estar em campo ao mesmo tempo. Invocar deixa de custar Ação: você chama um invocado por turno de graça. Seus invocados usam os seus PV como reserva de emergência --- ao chegar a 0, você pode transferir dano para si mesmo e mantê-lo de pé.}
\talento{Ninguém Chega Sozinho}{4}{Todos os seus invocados em campo recebem +2 na CA, +2 no acerto e imunidade a efeitos de dissipação de patamar Rei ou inferior.}
\magia{O Chamado que Não se Recusa}{s}{6}{22}{6 Ações}{Ilimitado (em qualquer lugar do mundo)}{\textbf{Ritual:} não pode ser encurtado.\par
Você invoca uma criatura com quem não tem Pacto --- qualquer criatura de rank Rei ou inferior que já tenha visto, viva ou morta. Teste de Espírito com Desvantagem: se falhar, obedece por 1 minuto; se passar, vem mesmo assim, furiosa, e o Mestre decide o que ela faz.}{\relax}
\magia{Corpo Emprestado}{\relax}{5}{18}{1 Ações}{Pessoal}{Por 10 minutos, você transfere sua consciência para o corpo de um invocado, usando os PV, ataques e sentidos dele; seu corpo fica inconsciente e protegido pelo círculo. Se o corpo emprestado morrer, você volta com 1 PV e um nível de Exaustão.}{\relax}

\section{Árvore do Corpo}

\subsection{Estilo Deus da Espada}
\begin{arvorecard}{Estilo Deus da Espada}{Espadachim}{Força \textbullet\ Recurso: PT}
"A vitória é de quem se move primeiro." Sem defesa, sem contra-ataque, o maior dano do livro. Doutrina: sem armadura média/pesada (perde tudo se vestir), CA base -2, nenhuma técnica concede Reações defensivas --- exceto a Reversão de Luz.
\end{arvorecard}
\patamar{Principiante}{Dado de PV \dado{1d8+3} \textbullet\ +1 degrau(s) de Dado de Arma}
\maestria{Quem Se Move Primeiro}{Empates de Iniciativa sempre são seus. No primeiro turno de qualquer combate, seu primeiro ataque tem Vantagem e rola o Dado de Arma uma vez a mais. Você nunca fica Surpreso enquanto empunhar espada e estiver consciente.}
\talento{Braço de Ferro}{1}{+4 PV por patamar seu nesta árvore. Aplicado sozinho na ficha, e cresce a cada patamar novo que você abrir nela.}
\talento{Fio Perfeito}{1}{Sua arma nunca lasca, entorta ou quebra por meios mundanos, e ataques contra objetos/estruturas causam dano dobrado.}
\talento{Pavio Curto}{1}{Vantagem em Intimidação, Desvantagem em qualquer teste social que exija paciência.}
\tecnica{Corte de Braço}{s}{2}{\relax}{1 Ações}{Corpo a corpo}{Ataque. Se acertar, o alvo faz teste de Vigor (CD 8 + Força + Rank): falha larga o que segura e não usa aquele braço até o fim do próximo turno. Contra rank Santo ou superior, apenas larga a arma.}
\tecnica{Investida}{\relax}{1}{\relax}{1 Ações}{Corpo a corpo}{\textbf{Dano:} +1 Dado de Arma se percorreu 9m ou mais\par
Avance até o dobro do Deslocamento em linha reta e ataque ao final. Sem Reações até seu próximo turno.}
\tecnica{Golpe Contínuo}{\relax}{1}{1}{1 Ações}{Corpo a corpo}{Se o ataque anterior neste turno acertou, este não pode errar por menos de 5.}
\tecnica{Roupa Leve}{\relax}{1}{\relax}{Passivo}{Passivo}{Sem armadura: +3m de Deslocamento e +2 na CA nos turnos em que você se moveu 3m ou mais. Vestir qualquer armadura desliga isto.}
\tecnica{Sem Recuo}{\relax}{1}{\relax}{Passivo}{Passivo}{Você é imune a Amedrontado e a qualquer efeito que force recuo ou fuga.}
\patamar{Intermediário}{Dado de PV \dado{1d10+3} \textbullet\ +1 degrau(s) de Dado de Arma \textbullet\ +1 PT}
\maestria{Aura Precoce}{Exceção do livro: o Deus da Espada acorda o Touki no 2º patamar, não no 3º. Você recebe Pontos de Touki e pode usar Touki Concentrado e Lâmina de Touki. Ainda não recebe o Manto de Touki completo (chega no Avançado).}
\talento{Punho Duplo}{1}{Você empunha arma de duas mãos com uma só, sem penalidade, e ganha +1 grau de Dado quando usa as duas mãos.}
\talento{Aço Rápido}{1}{+2 PT Máximos.}
\talento{Cavaleiro Vencido}{1}{Contra soldados, guardas e cavaleiros comuns, seus ataques acertam automaticamente com resultado 10 ou mais no dado.}
\tecnica{Passo Encurtado}{s}{2}{1}{1 Ações}{9 metros}{Desloque-se até 9m em linha reta, atravessando espaço de criaturas sem provocar oportunidade, e ataque imediatamente. Contra alvo que ainda não agiu, Vantagem.}
\tecnica{Dois Cortes}{\relax}{1}{1}{1 Ações}{Corpo a corpo}{\textbf{Dano:} Segundo ataque: um degrau abaixo\par
Dois ataques contra o mesmo alvo. Se ambos acertarem, o alvo fica Caído.}
\tecnica{Quebra-Armadura}{\relax}{1}{1}{1 Ações}{Corpo a corpo}{\textbf{Dano:} +2 Dados de Arma contra armadura completa\par
Ataque que ignora metade da CA de armadura não-mágica e de escudos.}
\tecnica{Leitura de Abertura}{\relax}{1}{\relax}{1 Ações}{Corpo a corpo}{Observe uma criatura por um turno inteiro sem atacar. No próximo turno, seu primeiro ataque contra ela tem Vantagem e crítico em 19-20.}
\patamar{Avançado}{Dado de PV \dado{1d10+4} \textbullet\ +2 degrau(s) de Dado de Arma \textbullet\ +1 PT}
\maestria{Velocidade Encarnada}{Você recebe o Manto de Touki completo. Dois degraus de Dado de Arma neste patamar. Se você não se mover no turno, recebe 1 Ação adicional só para atacar.}
\talento{Fôlego de Aço}{2}{+3 PT Máximos.}
\talento{Mira no Osso}{2}{Seus acertos críticos rolam o Dado de Arma três vezes em vez de duas.}
\talento{Espada Emprestada}{2}{Você recebe uma espada de qualidade superior: +1 degrau de Dado de Arma permanente, e ela conta como mágica.}
\tecnica{Espada do Silêncio}{s}{3}{2}{1 Ações}{Corpo a corpo}{\textbf{Dano:} +2 Dados de Arma\par
Ataque que não pode ser alvo de Reação alguma. Não faz som algum --- quem não estiver olhando não percebe que houve um ataque.}
\tecnica{Corte Ascendente}{\relax}{2}{1}{1 Ações}{Corpo a corpo}{Ataque que lança o alvo 4,5m para cima. Ele cai no fim do turno dele com dano de queda e Caído, a menos que voe. Vantagem contra ele enquanto no ar.}
\tecnica{Sede}{\relax}{2}{\relax}{Passivo}{Passivo}{Sempre que reduzir uma criatura a 0 PV, recupere imediatamente 1 PT e gaste 1 Ação extra ainda neste turno. Uma vez por turno.}
\patamar{Santo}{Dado de PV \dado{1d12+4} \textbullet\ +1 degrau(s) de Dado de Arma \textbullet\ +1 PT}
\maestria{Controle Absoluto}{Seus acertos críticos ocorrem em 19-20. Touki Concentrado passa a ser gratuito uma vez por turno. Quando for alvo de efeito de área com teste para metade do dano, um sucesso resulta em nenhum dano.}
\tecnica{Espada de Luz}{s}{4}{3}{2 Ações}{Corpo a corpo}{\textbf{Dano:} Dado de arma rolado três vezes (dobrado contra armadura pesada/estrutura)\par
Ataque com Vantagem que ignora todo bônus de CA de armadura, escudo e Cobertura. Não pode ser alvo de Reação, exceto a Reversão de Luz. Preço: sua CA cai em 5 até o início do próximo turno. Pré-requisito para o rank Rei da Espada.}
\patamar{Rei}{Dado de PV \dado{1d12+5} \textbullet\ +2 degrau(s) de Dado de Arma \textbullet\ +1 PT}
\maestria{A Lâmina do Deus}{Dois degraus de Dado de Arma. Você recebe uma espada da coleção do Deus da Espada: arma mágica, indestrutível, ignora Resistência a cortante. Sua Espada de Luz passa a custar 2 PT, e a penalidade de CA cai para -3.}
\tecnica{Reversão de Luz}{s}{5}{3}{1 Reação}{3 metros}{A única Reação da árvore inteira: 1 Reação quando alguém usar Espada de Luz contra você ou aliado a 3m. Anula por completo e ataca o conjurador --- se acertar, ele perde permanentemente o uso da mão até receber Cura de rank Santo ou superior.}
\tecnica{Corte do Horizonte}{\relax}{4}{3}{1 Ações}{Linha de 27m}{Corte em linha reta de 27m por 1,5m. Cada criatura sofre dano completo, teste de Agilidade para metade. Corta paredes, portões e barreiras mágicas com PV ao meio.}
\tecnica{Um Só Movimento}{\relax}{4}{2}{1 Ações}{4,5 metros}{Uma vez por combate: ataque todas as criaturas hostis dentro de 4,5m com uma única rolagem, comparada contra a CA de cada. Dano completo em todas.}
\patamar{Imperador}{Dado de PV \dado{2d6+5} \textbullet\ +2 degrau(s) de Dado de Arma \textbullet\ +1 PT}
\maestria{A Vitória de Quem Se Move Primeiro}{Dois degraus de Dado de Arma. Críticos ocorrem em 18-20. A primeira vez em cada combate que acertar uma criatura com PV cheios, role o Dado de Arma duas vezes o normal --- isto NÃO se aplica a técnicas que já multiplicam o Dado de Arma por três ou mais (Espada de Luz, Espada de Luz Verdadeira): elas já são o pico do estilo e não empilham com esta Maestria. Você não pode ser Atrasado, Lentificado nem ter a Iniciativa reduzida.}
\tecnica{Espada de Luz Verdadeira}{s}{6}{5}{2 Ações}{Corpo a corpo}{\textbf{Dano:} Dado de arma rolado cinco vezes\par
Uma vez por combate: acerta automaticamente, sem rolagem. Ignora CA, Cobertura, Manto de Touki, armadura mágica e barreira física. Só responde à Reversão de Luz de rank Rei ou superior (com Desvantagem). Preço: sua CA cai para 10 até o início do próximo turno, sem Reações.}
\tecnica{Sem Segunda Vez}{\relax}{5}{3}{1 Reação}{9 metros}{1 Reação, quando uma criatura a até 9m tentar fugir, teleportar, curar acima de metade dos PV, ou conjurar magia Santo ou superior. Ataque automático; se o dano igualar/superar metade dos PV atuais dela, a ação é cancelada e o recurso gasto se perde.}
\tecnica{Herança de Gal}{\relax}{5}{\relax}{Passivo}{Passivo}{Escolha uma técnica de rank Santo ou inferior de qualquer estilo marcial. Você a possui, com o seu próprio Bônus de Rank do Deus da Espada.}

\subsection{Estilo Deus da Água}
\begin{arvorecard}{Estilo Deus da Água}{Espadachim}{Vigor \textbullet\ Recurso: PT}
Aparar e devolver. Considerado o mais fraco dos três estilos porque um Deus da Água que enfrenta quem não ataca não faz nada --- mas é o único que mantém outras pessoas vivas. Sem relação com a Magia de Água --- chamado Suishin-ryū nas referências cruzadas.
\end{arvorecard}
\patamar{Principiante}{Dado de PV \dado{1d8+3} \textbullet\ +1 degrau(s) de Dado de Arma}
\maestria{Fluxo}{Reação: quando uma criatura adjacente erra um ataque corpo a corpo contra você, você contra-ataca imediatamente com dano de arma normal, sem custo de PT. No Principiante, uma vez por rodada.}
\talento{Casco de Tartaruga}{1}{+4 PV por patamar seu nesta árvore. Aplicado sozinho na ficha, e cresce a cada patamar novo que você abrir nela.}
\talento{Olho na Mão}{1}{Vantagem em Intuição para prever a próxima ação de uma criatura; o Mestre é obrigado a dar uma dica honesta.}
\talento{Paciência de Pedra}{1}{Você é imune a Amedrontado e a efeitos que forcem você a agir contra a vontade. Provocação não funciona em você.}
\tecnica{Aparar}{s}{2}{\relax}{1 Reação}{Corpo a corpo}{1 Reação, quando alvo de ataque corpo a corpo: some seu Bônus de Rank à CA contra aquele ataque, resolvido depois de ver a rolagem. Se com isso o ataque errar, o Fluxo dispara e não consome sua Reação daquele turno.}
\tecnica{Guarda do Corpo}{\relax}{1}{\relax}{1 Reação}{1,5 metros}{1 Reação: quando aliado a até 1,5m for alvo de ataque, você vira o alvo. Aplique sua CA --- se o ataque errar, o Fluxo dispara.}
\tecnica{Peso Não Atrapalha}{\relax}{1}{\relax}{Passivo}{Passivo}{Você usa armadura pesada sem penalidade de Furtividade, deslocamento ou fadiga, e dorme com ela vestida.}
\tecnica{Provocar}{\relax}{1}{\relax}{1 Ações}{Visão e audição}{Teste de Espírito (CD 8 + Espírito + Rank). Se falhar, no próximo turno o alvo tem que atacar você se conseguir alcançá-lo. Uma vez por criatura por combate.}
\tecnica{Base Firme}{\relax}{1}{\relax}{Passivo}{Passivo}{Você não pode ser movido contra a vontade, derrubado ou empurrado por nada que permita teste de Força ou Agilidade.}
\patamar{Intermediário}{Dado de PV \dado{1d8+3} \textbullet\ +1 degrau(s) de Dado de Arma}
\maestria{A Armadura Não Pesa}{Proficiência com toda armadura e escudo, ignorando penalidades. Usando armadura pesada, +1 na CA além do valor normal. O seu Fluxo passa a disparar duas vezes por rodada.}
\talento{Guarda Longa}{1}{Seu alcance de Reação corpo a corpo aumenta para 3 metros.}
\talento{Aço Calmo}{1}{+2 PT Máximos.}
\talento{Nome de Reidar}{1}{Seu cônjuge abandonou a própria casa pra estar com você --- costume herdado da princesa. Você tem uma pessoa de lealdade absoluta no mundo.}
\tecnica{Devolver}{s}{2}{1}{1 Reação}{Corpo a corpo}{\textbf{Dano:} Dano de arma + metade do dano que o ataque causaria\par
1 Reação, quando um ataque corpo a corpo contra você errar. Você devolve o golpe com a força dele somada à sua.}
\tecnica{Trava de Lâmina}{\relax}{1}{1}{1 Reação}{Corpo a corpo}{1 Reação, ao aparar um ataque: disputa de Força. Se vencer, o inimigo não ataca no próximo turno, e você pode desarmá-lo.}
\tecnica{Muralha de Um Homem}{\relax}{1}{\relax}{Passivo}{Passivo}{Enquanto você não se mover no seu turno, aliados adjacentes recebem Cobertura Superior contra ataques à distância vindos da sua direção.}
\tecnica{Contra-Investida}{\relax}{1}{1}{1 Reação}{Corpo a corpo}{1 Reação, quando um inimigo se mover para dentro do seu alcance corpo a corpo: você ataca antes que ele chegue, interrompendo a ação dele se acertar.}
\patamar{Avançado}{Dado de PV \dado{1d10+4} \textbullet\ +1 degrau(s) de Dado de Arma \textbullet\ +1 PT}
\maestria{A Postura}{Você recebe o Manto de Touki completo e a Postura de Água (Kamae, 1 Ação pra entrar): Deslocamento 0, não ataca, soma Bônus de Rank à CA, ganha Reações extras (Bônus de Rank ÷ 2, arredondado pra cima), e o Fluxo passa a disparar também quando você apara com sucesso. Sair é livre ao se mover.}
\talento{Postura Móvel}{2}{Em Postura, você pode se mover até 3 metros por turno sem sair dela.}
\talento{Segunda Guarda}{2}{+1 Reação por turno, mesmo fora da Postura.}
\talento{Escudo Vivo}{2}{Guarda do Corpo passa a alcançar 3 metros e pode ser usada uma vez por turno sem gastar Reação.}
\tecnica{Fluxo Verdadeiro}{s}{3}{2}{1 Reação}{9 metros}{1 Reação: apara qualquer flecha, virote, projétil mágico, ou magia de ataque de alvo único de rank Avançado ou inferior. O efeito é anulado, ou redirecionado pra outra criatura a 9m com a CD original. Não funciona contra área nem Espada de Luz.}
\tecnica{Correnteza}{\relax}{2}{1}{Passivo}{Contra-ataque de Fluxo}{Gaste 1 PT no início da sua rodada para ativar. Cada contra-ataque de Fluxo que acertar naquela rodada causa +1 Dado de Arma cumulativo, até o teto de +2 Dados de Arma num mesmo contra-ataque: o terceiro causa dois dados extras, e do quarto em diante o bônus permanece em dois.}
\tecnica{Peso da Água}{\relax}{2}{1}{1 Reação}{Corpo a corpo}{1 Reação, ao aparar: o atacante faz teste de Força ou fica Caído e perde o restante das Ações dele no turno.}
\patamar{Santo}{Dado de PV \dado{1d10+4} \textbullet\ +1 degrau(s) de Dado de Arma \textbullet\ +1 PT}
\maestria{Ler o Fluxo}{O Mestre é obrigado a te informar o que um inimigo pretende fazer antes de resolver a ação. Fluxo Verdadeiro passa a funcionar contra magias de qualquer rank, exceto área e rank Imperador. Você é imune a acertos críticos em Postura.}
\tecnica{O Primeiro Segredo}{s}{4}{3}{1 Reação}{Corpo a corpo}{Um dos Cinco Segredos --- criado com o Mestre. Escolha um gatilho (inimigo ataca você/aliado, conjura, se move, foge) e duas opções de orçamento (anula o efeito; devolve o dano; aplica uma condição --- Caído, Preso, Desarmado ou Atordoado; atinge todos a 3m; +5 CA a você e aliado), mais uma Amarra (condição de uso).}
\tecnica{Espelho}{\relax}{3}{2}{1 Reação}{Corpo a corpo}{1 Reação: quando um inimigo usar contra você técnica marcial de rank Santo ou inferior, você a executa de volta imediatamente, mesmo sem possuí-la.}
\patamar{Rei}{Dado de PV \dado{1d12+5} \textbullet\ +1 degrau(s) de Dado de Arma \textbullet\ +1 PT}
\maestria{A Arte da Provocação}{Provocar vira gratuito, uma vez por turno. Quem falhar ataca você com Desvantagem. Em Postura, inimigos que optarem por não atacar você e estiverem ao alcance sofrem Desvantagem em tudo naquele turno.}
\tecnica{O Segundo Segredo}{s}{5}{3}{1 Reação}{Corpo a corpo}{Invente seu segundo Segredo com o mesmo molde, escolhendo três opções de orçamento em vez de duas.}
\tecnica{Maré de Retorno}{\relax}{4}{3}{1 Ações}{Postura}{Requer Postura. Até o início do próximo turno, todo ataque corpo a corpo que errar você dispara Fluxo automaticamente, até um número de Reações igual ao seu Bônus de Rank por rodada (continua sendo a maior contagem de Reações do livro).}
\tecnica{Nada Passa}{\relax}{4}{2}{1 Reação}{Pessoal}{1 Reação: você apara um efeito de área --- ele acontece, mas não atinge você nem aliados adjacentes a você.}
\patamar{Imperador}{Dado de PV \dado{1d12+5} \textbullet\ +1 degrau(s) de Dado de Arma \textbullet\ +1 PT}
\maestria{Domínio Absoluto}{Suas Reações em Postura deixam de ter limite numérico (uma vez por criatura por turno). Seu alcance de Reação corpo a corpo aumenta para 4,5 metros. Você pode entrar em Postura como Reação no instante em que o combate começar.}
\tecnica{Reino da Espada da Privação}{s}{6}{5}{2 Ações}{Esfera de 12m}{Requer Postura (não pode sair). Estabelece um domínio esférico de 12m: qualquer hostil que se mova, ataque, conjure ou canalize mana dentro é atingido automaticamente, uma vez por turno dela. Dura enquanto mantiver Postura e tiver PT (1 PT/turno). Um alvo completamente imóvel não é atingido, e não responde à Espada de Luz Verdadeira a tempo.}
\tecnica{O Terceiro Segredo}{\relax}{5}{4}{1 Ações}{Corpo a corpo}{Invente seu terceiro Segredo, com quatro opções de orçamento.}

\subsection{Estilo Deus do Norte}
\begin{arvorecard}{Estilo Deus do Norte}{Espadachim}{Força ou Agilidade \textbullet\ Recurso: PT}
Não é esgrima refinada --- é um método para continuar vivo. Truques, improviso e terreno; a maioria dos aventureiros aprende Norte.
\end{arvorecard}
\patamar{Principiante}{Dado de PV \dado{1d8+2} \textbullet\ +1 degrau(s) de Dado de Arma}
\maestria{Sobreviver é Vencer}{+1 degrau no Dado de Arma. Proficiência universal: qualquer arma ou objeto improvisado, sem penalidade (objetos usam d6 base). [Improviso] O Improviso: uma vez por combate, descreva uma manobra usando o cenário e role o atributo indicado pelo Mestre com Vantagem.}
\talento{Três Bainhas}{1}{Você carrega armas escondidas. Sacar uma arma nova é livre, e você nunca fica realmente desarmado.}
\talento{Osso Duro}{1}{+4 PV por patamar seu nesta árvore. Aplicado sozinho na ficha, e cresce a cada patamar novo que você abrir nela.}
\talento{Pés no Chão}{1}{Você não sofre penalidade de deslocamento em terreno difícil, escombros ou gelo.}
\tecnica{Forma Quadrúpede}{s}{2}{\relax}{1 Ações}{Corpo a corpo}{\textbf{Dano:} +1d6 no ataque ao final do avanço\par
1 Ação: avance até o dobro do seu Deslocamento em linha reta (ataques à distância contra você têm Desvantagem durante o avanço). Ao final, gaste a próxima Ação para atacar com +1d6.}
\tecnica{Arremesso de Espada}{\relax}{1}{\relax}{1 Ações}{18 metros}{\textbf{Dano:} Dano de arma normal\par
1 Ação: ataque à distância (Agilidade + Rank) até 18m. Alternativamente, use como Reação pra interceptar um projétil ou magia de rank Principiante (teste de Agilidade CD 12 pra anular). Você fica desarmado.}
\tecnica{Bala de Lágrimas}{\relax}{1}{\relax}{1 Ações}{Cone de 3m}{Requer um saco de temperos preparado (2 PO na cidade, ou 1 hora colhendo ervas). Teste de Vigor (CD 8 + Agilidade + Rank) ou o alvo fica Cego até o fim do próximo turno.}
\tecnica{Primeiros Socorros de Campo}{\relax}{1}{\relax}{1 Ações}{Toque (adjacente)}{\textbf{Dano:} Vigor + Bônus de Rank em PV\par
Um aliado adjacente (ou você) recupera PV e remove todas as Marcas da Morte. Uma vez por Descanso Curto por criatura. Concede a perícia Medicina.}
\tecnica{Leitura de Rastro}{\relax}{1}{\relax}{Passivo}{Passivo}{Você ganha as perícias Sobrevivência e Percepção, e rola com Vantagem para rastrear alvos por terreno pisado, sangue ou hábitos.}
\tecnica{Golpe Baixo}{\relax}{1}{\relax}{1 Ações}{Corpo a corpo}{\textbf{Dano:} Metade do dado de arma\par
Teste de Força ou Agilidade (CD 8 + Força + Rank) ou o alvo fica Caído. Contra alvos Caídos, seus ataques corpo a corpo têm Vantagem.}
\patamar{Intermediário}{Dado de PV \dado{1d8+3} \textbullet\ +1 degrau(s) de Dado de Arma}
\maestria{O Corpo Ferido}{Você não sofre Desvantagem por lutar Caído, com uma mão só, em terreno difícil, no escuro, ou sob 1 nível de Exaustão. Com metade ou menos dos PV Máximos, +2 de dano em ataques corpo a corpo. Testes do Fio da Vida com Vantagem.}
\talento{Segunda Chance [Improviso]}{1}{Uma vez por combate, ao errar um ataque, reposicione usando o cenário e repita a rolagem.}
\talento{Facção Errante}{1}{Escolha uma arma exótica: ganha dado base d8 e um efeito de condição (empurrar, prender ou derrubar) uma vez por combate.}
\talento{Estômago de Mercenário}{1}{Você ignora fome, sede, clima extremo e uma noite sem dormir. Vantagem contra veneno.}
\tecnica{Finta do Norte}{s}{2}{\relax}{1 Ações}{Corpo a corpo}{Teste de Enganação (Espírito) contra a Intuição do alvo. Se vencer, seu próximo ataque neste ou no próximo turno acerta automaticamente, e é crítico se o alvo estiver desprevenido.}
\tecnica{Empunhadura Dupla}{\relax}{1}{\relax}{1 Ações}{Corpo a corpo}{\textbf{Dano:} Arma secundária: um degrau abaixo do normal\par
Você empunha duas armas de uma mão. Enquanto duas armas empunhadas, recebe +1 na CA.}
\tecnica{Desarme}{\relax}{1}{\relax}{1 Ações}{Corpo a corpo}{Disputa de Força ou Agilidade. Se vencer, a arma do alvo voa 3m. Se vencer por 10+, você a pega no ar e passa a usá-la com todos os seus degraus de Dado de Arma.}
\tecnica{Passo do Terreno [Improviso]}{\relax}{1}{\relax}{1 Ações}{Deslocamento}{Mova-se até o Deslocamento ignorando terreno difícil, escalando livremente. Ao terminar em posição elevada, +2 CA e Vantagem no próximo ataque, até sair de lá.}
\tecnica{Ferro Frio}{\relax}{1}{\relax}{Passivo}{Passivo}{Com um quarto ou menos dos PV Máximos, você tem Resistência a todo dano físico e é imune a Amedrontado e Atordoado.}
\patamar{Avançado}{Dado de PV \dado{1d10+3} \textbullet\ +1 degrau(s) de Dado de Arma \textbullet\ +1 PT}
\maestria{O Despertar do Touki}{Você passa a vestir Touki conscientemente: recebe o Manto de Touki e a reserva de Pontos de Touki (Cap. 3). Exclusivo do Norte: você pode aplicar Touki a objetos improvisados --- uma cadeira revestida de aura parte um escudo de aço.}
\talento{Fôlego Longo}{2}{+3 PT Máximos.}
\talento{Mão Trocada}{2}{Você usa Empunhadura Dupla sem redução de degrau na arma secundária.}
\talento{Instinto de Presa}{2}{Você não pode ser surpreendido enquanto consciente, e rola Iniciativa com Vantagem.}
\tecnica{Corte Reverso}{s}{3}{1}{1 Ações}{Corpo a corpo}{\textbf{Dano:} +1d10 se ignorar uma Reação defensiva\par
Se o ataque for aparado/bloqueado/anulado por Reação, ele acontece mesmo assim, ignorando a Reação. Contra Deus da Água, não gasta PT.}
\tecnica{Túmulo de Aço [Improviso]}{\relax}{2}{1}{1 Ações}{Corpo a corpo}{\textbf{Dano:} 3d10 (contundente)\par
Requer cenário utilizável. Teste de Agilidade (CD 8 + Força + Rank). Falha: dano e Preso até gastar 1 Ação para se soltar. Sucesso: metade do dano e não fica Preso.}
\tecnica{Quebra-Guarda}{\relax}{2}{1}{1 Ações}{Corpo a corpo}{Ataque que ignora Cobertura, bônus de escudo e metade da CA de armaduras não-mágicas. Contra barreiras mágicas com PV, dano dobrado.}
\tecnica{Pele de Ferro}{\relax}{2}{\relax}{Passivo}{Passivo}{O bônus de CA do seu Manto de Touki passa a ser o Bônus de Rank completo, não a metade.}
\patamar{Santo}{Dado de PV \dado{1d10+4} \textbullet\ +1 degrau(s) de Dado de Arma \textbullet\ +1 PT}
\maestria{Estilo Pessoal}{Crie, com o Mestre, uma técnica assinatura própria: custo 1 Ação e 1 PT; escolha duas opções (dano extra 2d10; aplica condição com teste de resistência; atinge dois alvos; alcance 9m; +4 CA até o próximo turno); e defina uma Amarra --- condição que a técnica exige pra funcionar.}
\tecnica{Cruz Nebulosa}{s}{4}{2}{2 Ações}{Corpo a corpo}{\textbf{Dano:} +4d8 no segundo corte, ignora Manto de Touki\par
Requer Empunhadura Dupla e uma terceira arma embainhada. Ataque com Vantagem; se o alvo bloquear/aparar, o segundo corte acerta automaticamente.}
\tecnica{Contra-Água}{\relax}{3}{1}{1 Reação}{Corpo a corpo}{1 Reação, quando um inimigo errar um ataque contra você: ataque imediatamente com Vantagem. Contra praticante do Deus da Água, não gasta PT.}
\tecnica{Golpe do Desespero}{\relax}{3}{2}{2 Ações}{Corpo a corpo}{\textbf{Dano:} Dado de arma rolado três vezes\par
Só usável com metade ou menos dos PV. Depois de usar, 1 nível de Exaustão até o próximo Descanso Longo.}
\patamar{Rei}{Dado de PV \dado{1d12+4} \textbullet\ +1 degrau(s) de Dado de Arma \textbullet\ +1 PT \textbullet\ desbloqueio 2 PA (exceção)}
\maestria{Leitura de Batalha}{Ao final do primeiro turno de combate, identifique estilo, rank aproximado e maior fraqueza de um inimigo --- contra ele, +2 em acertos e testes de resistência pelo resto do combate. Vantagem de Estilo contra o Deus da Água não é mais anulada por diferença de rank. O Improviso pode ser usado três vezes por combate.}
\tecnica{Mil Facções}{s}{5}{2}{1 Ações}{Corpo a corpo}{Replique, neste turno, qualquer técnica de rank Avançado ou inferior de qualquer estilo marcial que já tenha visto, mesmo sem possuí-la. Uma vez por combate.}
\tecnica{Dança de Aço}{\relax}{4}{2}{1 Ações}{3 metros}{Ataque cada criatura à sua escolha dentro de 3m, com rolagem separada para cada, sem Desvantagem por aliados na área.}
\tecnica{Aura Cortante}{\relax}{4}{3}{1 Ações}{Linha de 27 metros}{\textbf{Dano:} +2d10, somados ao dano de arma normal\par
Libera o Touki como lâmina em linha. Teste de Agilidade (CD 8 + Força + Rank) para metade. Corta estruturas de pedra e madeira no caminho.}
\patamar{Imperador}{Dado de PV \dado{1d12+5} \textbullet\ +1 degrau(s) de Dado de Arma \textbullet\ +1 PT}
\maestria{Nada é Regra}{Você recebe uma Ação adicional em cada turno (total 4). Técnicas [Improviso] funcionam mesmo em ambientes vazios. Gaste 1 PT para transformar uma falha crítica sua em rolagem normal.}
\tecnica{Golpe do Fim da Linha}{s}{6}{4}{3 Ações}{Corpo a corpo}{\textbf{Dano:} Dado de arma rolado quatro vezes + 4d12\par
Uma vez por combate. Acerta automaticamente, ignora Touki, Cobertura e armadura. Se o alvo estiver com metade ou menos dos PV, dano dobrado.}
\tecnica{Aura de Comando}{\relax}{5}{3}{1 Ações}{18 metros}{1 minuto: aliados em 18m recebem +2 em acertos, imunidade a Amedrontado, e podem repetir um teste de resistência falho por turno.}
\tecnica{Corpo Impossível}{\relax}{5}{\relax}{Passivo}{Passivo}{Enquanto tiver ao menos 1 PT, você é imune a Paralisia, Petrificação, Preso e Atordoado. Membros quebrados não reduzem seus atributos.}

\subsection{Lutador}
\begin{arvorecard}{Lutador}{Guerreiro / Parrudice}{Força \textbullet\ Recurso: PT}
Peso e impacto --- a única árvore onde lutar de mãos vazias é opção legítima até o fim. Degrada o inimigo (Quebrantado) em vez de matar rápido.
\end{arvorecard}
\patamar{Briguento}{Dado de PV \dado{1d10+3} \textbullet\ +1 degrau(s) de Dado de Arma}
\maestria{O Corpo é a Arma}{Seus ataques desarmados usam Dado Base d6 (com seus degraus normais). Proficiência com toda arma pesada, de duas mãos e improvisada; empunhar arma de duas mãos com uma só custa apenas um degrau a menos. Você aplica a condição Quebrantado e usa o combo de Momento (+1 Dado de Arma correndo 6m+ em linha reta antes de atacar).}
\talento{Couro Grosso}{1}{+4 PV por patamar seu nesta árvore. Aplicado sozinho na ficha, e cresce a cada patamar novo que você abrir nela.}
\talento{Punho de Mineiro}{1}{Seus ataques desarmados causam dano letal e contam como arma pesada para qualquer efeito.}
\talento{Sem Vergonha}{1}{Você pega qualquer objeto do cenário e o usa como arma pesada (d6, todos os seus degraus); ele quebra depois de três acertos.}
\tecnica{Investida Devastadora [Peso]}{s}{2}{\relax}{1 Ações}{Corpo a corpo}{\textbf{Dano:} +2 Dados de Arma\par
Requer 6m de corrida em linha reta. Se acertar, teste de Força (CD 8 + Força + Rank) ou o alvo fica Caído e ganha 2 acúmulos de Quebrantado.}
\tecnica{Agarrão [Impacto]}{\relax}{1}{\relax}{1 Ações}{Corpo a corpo}{Disputa de Força. Vencendo, o alvo fica Agarrado (Deslocamento 0, Desvantagem contra outros) e você o arrasta 3m por turno. Enquanto agarrado, seus ataques desarmados contra ele acertam automaticamente.}
\tecnica{Golpe Circular [Peso]}{\relax}{1}{\relax}{1 Ações}{Corpo a corpo}{\textbf{Dano:} Dado de arma um degrau abaixo\par
Ataque contra todas as criaturas adjacentes, com rolagem separada para cada.}
\tecnica{Cabeçada [Impacto]}{\relax}{1}{\relax}{1 Ações}{Corpo a corpo}{\textbf{Dano:} Metade do dado (você sofre 1d4)\par
Teste de Vigor (CD 8 + Força + Rank) ou o alvo fica Atordoado até o fim do próximo turno dele.}
\tecnica{Quebrar Equipamento [Peso]}{\relax}{1}{\relax}{1 Ações}{Corpo a corpo}{Teste de Força (CD 8 + Força + Rank): falha quebra o escudo ou arma mundana do alvo, ou a inutiliza por 3 turnos se for mágica.}
\patamar{Combatente}{Dado de PV \dado{1d10+3} \textbullet\ +1 degrau(s) de Dado de Arma}
\maestria{Não Para de Vir}{O bônus de Momento sobe para +2 Dados de Arma. Sem penalidade lutando agarrado, em espaço apertado, no chão ou preso --- na verdade, +2 nessas situações. Criaturas Agarradas por você ganham 1 acúmulo de Quebrantado no início do turno delas, automaticamente.}
\talento{Mãos Grandes}{1}{Você agarra criaturas de qualquer tamanho e carrega o dobro de peso.}
\talento{Fôlego de Fosso}{1}{+2 PT Máximos.}
\talento{Cicatriz Velha}{1}{Enquanto estiver com metade ou menos dos PV, seus ataques ganham +1 Dado de Arma.}
\tecnica{Arremesso [Impacto]}{s}{2}{\relax}{1 Ações}{9 metros}{\textbf{Dano:} 4d8 + Força (contundente)\par
Requer alvo Agarrado, de até duas categorias de tamanho acima da sua. Arremesse até 9m: fica Caído, e quem estiver no ponto de queda faz teste de Agilidade ou também sofre metade e cai.}
\tecnica{Golpe Ascendente [Peso]}{\relax}{1}{1}{1 Ações}{Corpo a corpo}{Ataque que lança o alvo 4,5m para cima. Ele cai no fim do turno com dano de queda e Caído.}
\tecnica{Trava [Impacto]}{\relax}{1}{1}{1 Ações}{Corpo a corpo}{Requer alvo Agarrado. Teste de Vigor do alvo (CD 8 + Força + Rank). Falha: Preso e Incapacitado enquanto você mantiver (você só pode manter).}
\tecnica{Peso Contra Peso [Peso]}{\relax}{1}{1}{1 Reação}{Corpo a corpo}{1 Reação, quando tentarem te empurrar, derrubar ou atravessar seu espaço: falha automática, e ganham 2 acúmulos de Quebrantado.}
\patamar{Veterano}{Dado de PV \dado{1d12+4} \textbullet\ +1 degrau(s) de Dado de Arma \textbullet\ +1 PT}
\maestria{Nada Segura}{Você recebe o Manto de Touki completo e a reserva de PT. Seus ataques atravessam: ao reduzir uma criatura a 0 PV, ou acertar alvo com 4+ acúmulos de Quebrantado, o ataque atinge outra criatura adjacente com dano completo. O limite de acúmulos passa a ser o dobro do seu Bônus de Rank.}
\talento{Braço de Bigorna}{2}{Seus ataques desarmados sobem um degrau adicional de Dado de Arma, permanentemente.}
\talento{Duas Mãos, Uma Arma}{2}{Empunhar arma de duas mãos com uma só deixa de custar degrau, e você pode segurar um escudo na outra.}
\talento{Cai Quem Toca}{2}{Criaturas que te acertarem corpo a corpo ganham 1 acúmulo de Quebrantado.}
\tecnica{Esmagar [Peso]}{s}{3}{2}{1 Ações}{Corpo a corpo}{\textbf{Dano:} Dado de arma rolado três vezes\par
Contra alvo Caído, Agarrado, Preso ou Atolado, acerta automaticamente e aplica 3 acúmulos de Quebrantado.}
\tecnica{Onda de Choque [Peso]}{\relax}{2}{2}{1 Ações}{Esfera de 6m}{\textbf{Dano:} 5d8 (contundente)\par
Teste de Agilidade (CD 8 + Força + Rank). Falha: dano, Caído e 1 acúmulo de Quebrantado. Estruturas e barreiras com PV sofrem dano dobrado.}
\tecnica{Estrangular [Impacto]}{\relax}{2}{1}{1 Ações}{Corpo a corpo}{\textbf{Dano:} 3d6 por turno\par
Requer alvo Agarrado. Ele não consegue falar, gritar nem recitar cântico algum enquanto você mantiver.}
\patamar{Campeão}{Dado de PV \dado{1d12+4} \textbullet\ +1 degrau(s) de Dado de Arma \textbullet\ +1 PT}
\maestria{Ainda Estou Aqui}{Você não pode ser derrubado, empurrado, agarrado, arremessado nem movido contra a vontade por nada de patamar igual ou inferior. Ao ser reduzido a 0 PV, termina o turno atual antes de cair, com todas as Ações restantes (1x/combate). Imune a Amedrontado e Atordoado.}
\talento{Colheita}{3}{Cada criatura reduzida a 0 PV por você devolve 1 PT e permite mover-se 3m sem gastar Ação; uma vez por turno, permite também um ataque ou agarrão extra.}
\talento{Peso Absoluto}{3}{O Arremesso alcança 18m e o dano dobra.}
\tecnica{Ruína [Peso]}{s}{4}{3}{2 Ações}{Corpo a corpo}{\textbf{Dano:} Dado de arma rolado quatro vezes (triplicado contra estruturas)\par
Ignora todo bônus de CA de escudo e Cobertura, e aplica acúmulos de Quebrantado iguais ao seu Bônus de Rank de uma vez.}
\tecnica{Mão na Garganta [Impacto]}{\relax}{3}{2}{1 Ações}{9 metros}{Atravesse até 9m ignorando ataques de oportunidade e agarre automaticamente uma criatura visível, sem disputa. Ela ganha 2 acúmulos de Quebrantado.}
\patamar{Mestre de Guerra}{Dado de PV \dado{2d6+5} \textbullet\ +1 degrau(s) de Dado de Arma \textbullet\ +1 PT}
\maestria{Arma Viva}{Seu alcance corpo a corpo aumenta para 3 metros, com qualquer arma ou mãos vazias. Seus ataques contam como armas de cerco contra estruturas. Cada acúmulo de Quebrantado no alvo concede a você +1 no acerto contra ele.}
\tecnica{Não Sobra Formação [Peso]}{s}{5}{3}{1 Ações}{Linha de 18m × 3m}{Ataque contra todas as criaturas na linha, dano completo com rolagem comparada à CA de cada. Quem for atingido é arremessado 4,5m e fica Caído.}
\tecnica{Prensa [Impacto]}{\relax}{4}{3}{1 Ações}{Corpo a corpo}{\textbf{Dano:} 8d10 + Força (contundente, automático)\par
Requer alvo Agarrado. O alvo fica Quebrantado ao máximo instantaneamente. Se isso o reduzir a 0 PV, não pode ser estabilizado por meios mundanos.}
\patamar{Lenda Viva}{Dado de PV \dado{2d6+5} \textbullet\ +1 degrau(s) de Dado de Arma \textbullet\ +1 PT}
\maestria{Nada Fica de Pé}{Você recebe 1 Ação adicional por turno, usável só para atacar ou agarrar. Toda criatura que começar o turno adjacente a você ganha 1 acúmulo de Quebrantado. Criaturas com acúmulos iguais ao dobro do seu Bônus de Rank ficam Incapacitadas.}
\talento{Sem Arma Nenhuma}{4}{Seus ataques desarmados sobem para o Dado de Arma de um montante (d10 base) e contam como mágicos, de cerco e adamantinos.}
\tecnica{O Golpe que Fecha a Conta [Peso]}{s}{6}{5}{2 Ações}{Corpo a corpo}{\textbf{Dano:} Dado de arma rolado cinco vezes + 1d12 por acúmulo de Quebrantado do alvo, até o seu Bônus de Rank em acúmulos\par
Uma vez por combate. Acerta automaticamente e ignora CA, Cobertura, escudo e Manto de Touki.}
\tecnica{Terremoto Pessoal [Peso]}{\relax}{5}{4}{2 Ações}{Esfera de 18m}{\textbf{Dano:} 14d10 (contundente)\par
Teste de Agilidade para metade. Toda estrutura na área desaba; terreno vira difícil permanentemente; voadores a menos de 9m do chão caem.}

\subsection{Cavalaria e Escudos}
\begin{arvorecard}{Cavalaria e Escudos}{Tank / Defensor}{Vigor \textbullet\ Recurso: PT}
Aparar e não devolver nada --- diferente do Suishin-ryū, o Defensor não tem Fluxo nem contragolpe. A pergunta única é quem está atrás de mim. Maior PV do livro; gasta PT mais rápido que qualquer outra árvore.
\end{arvorecard}
\patamar{Escudeiro}{Dado de PV \dado{1d10+4} \textbullet\ +1 degrau(s) de Dado de Arma}
\maestria{Interpor}{Você desbloqueia 'Sob Minha Guarda': designe aliados (até seu Bônus de Rank) como protegidos; a até 3m, gaste 1 Reação para que todo o dano de um ataque contra ele venha para você (não reduzível por Resistência, mas sim por PT). Se ele sofrer dano que você não interceptou, recupere 1 PT. Você é proficiente com toda armadura e escudo, e com escudo recebe +2 na CA além do normal.}
\talento{Ombro de Pedra}{1}{+4 PV por patamar seu nesta árvore. Aplicado sozinho na ficha, e cresce a cada patamar novo que você abrir nela.}
\talento{Montaria}{1}{Você monta, treina e acalma qualquer besta de carga. Sobre ela, você não cai por efeito que permita teste, e sua montaria também está Sob Sua Guarda.}
\talento{Sono de Ferro}{1}{Você dorme de armadura completa sem penalidade e acorda pronto. Vantagem contra Exaustão por marcha ou vigília.}
\tecnica{Muralha de Um}{s}{2}{1}{1 Ações}{3 metros}{Até o início do próximo turno, você não pode se mover, e aliados a até 3m recebem Cobertura Superior (+5 CA) e Resistência a dano de área. Você não recebe nenhum benefício.}
\tecnica{Golpe de Escudo}{\relax}{1}{\relax}{1 Ações}{Corpo a corpo}{\textbf{Dano:} 1d8 + Força (contundente)\par
Teste de Força do alvo (CD 8 + Vigor + Rank) ou é empurrado 3m e fica Caído.}
\tecnica{Puxar}{\relax}{1}{\relax}{1 Ações}{9 metros}{Um aliado a até 9m é puxado para adjacente a você e fica Sob Sua Guarda imediatamente, mesmo excedendo seu limite.}
\tecnica{Provocar Ódio}{\relax}{1}{\relax}{1 Ações}{Visão}{Teste de Espírito (CD 8 + Vigor + Rank). Falha: no próximo turno da criatura, ataques contra qualquer um que não seja você têm Desvantagem. (Não é a Provocação do Suishin-ryū, que força o ataque pra alimentar contragolpe --- aqui você só torna caro ignorar você.)}
\patamar{Guarda}{Dado de PV \dado{1d12+4}}
\maestria{Peso do Aço}{O alcance de Sob Minha Guarda sobe para 4,5 metros. Você não pode ser empurrado, derrubado, agarrado nem movido contra a vontade com os pés no chão e escudo na mão. Com armadura pesada, Resistência a dano de área.}
\talento{Dois Escudos}{1}{Você empunha um escudo em cada mão: +2 na CA adicional, e não pode atacar.}
\talento{Fôlego de Sentinela}{1}{+2 PT Máximos.}
\talento{A Porta Sou Eu}{1}{Enquanto bloquear uma passagem de até 3m, criaturas Médias ou menores não atravessam sem antes te derrubar.}
\tecnica{Aguentar o Baque}{s}{2}{1}{1 Reação}{Pessoal}{\textbf{Dano:} Reduz 1d10 + Vigor + Bônus de Rank\par
1 Reação, ao interceptar dano por Sob Minha Guarda: reduza aquele dano antes de aplicá-lo em você. Empilha com quantas Reações tiver.}
\tecnica{Escudo Erguido}{\relax}{1}{1}{1 Ações}{Pessoal}{1 minuto: ataques à distância mundanos contra você e aliados a 3m erram automaticamente. Você não pode correr enquanto sustentar.}
\tecnica{Formação}{\relax}{1}{\relax}{Passivo}{Passivo}{Aliados adjacentes a você somam seu Bônus de Rank aos testes de resistência contra efeitos de área e contra ser movido.}
\tecnica{Cavalgada}{\relax}{1}{\relax}{1 Ações}{Deslocamento da montaria}{Requer montaria. Avance até o dobro do deslocamento dela atravessando linhas inimigas --- no caminho, teste de Força ou Caídas. Sem ataques de oportunidade.}
\patamar{Protetor}{Dado de PV \dado{1d12+5} \textbullet\ +1 degrau(s) de Dado de Arma \textbullet\ +2 PT}
\maestria{Escudo Estendido}{Você recebe o Manto de Touki completo e a reserva de PT --- recebe 2 PT por patamar em vez de 1. O Touki é projetado para fora do corpo: Sob Minha Guarda alcança 9 metros e exige apenas linha de visão. Interceptar dano deixa de custar Reação uma vez por rodada.}
\talento{Casco}{2}{Você tem Resistência a dano físico de armas mundanas enquanto empunhar escudo.}
\talento{Guarda Ampla}{2}{O número de aliados Sob Sua Guarda passa a ser o dobro do seu Bônus de Rank.}
\talento{Aço Paciente}{2}{+4 PT Máximos.}
\tecnica{Não Ele}{s}{3}{2}{1 Reação}{9 metros}{1 Reação: intercepte um ataque, magia ou efeito de área inteiro dirigido a um aliado Sob Sua Guarda, mesmo de alvo único e mesmo sem alcance físico. Você sofre o efeito completo no lugar dele, inclusive condições.}
\tecnica{Redirecionar}{\relax}{2}{1}{1 Reação}{9 metros}{Ao interceptar um ataque à distância ou projétil mágico, desvie-o para uma criatura hostil à sua escolha a até 9m, usando a rolagem original.}
\tecnica{Fôlego Emprestado}{\relax}{2}{2}{1 Ações}{9 metros}{\textbf{Dano:} PV Temporários = Vigor + Bônus de Rank\par
Um aliado Sob Sua Guarda recebe PV Temporários e remove uma condição de Amedrontado, Atordoado ou Caído.}
\patamar{Guardião}{Dado de PV \dado{2d6+5} \textbullet\ +2 PT}
\maestria{Aegis}{Aliados Sob Sua Guarda recebem, passivamente e sem custo: redução de dano igual ao seu Bônus de Rank contra todo dano recebido, imunidade a acertos críticos, e o direito de repetir um teste de resistência falho por turno.}
\talento{Aço Vivo}{3}{Sua armadura e escudo se reparam sozinhos após cada Descanso Curto e não podem ser destruídos por efeito algum de patamar inferior ao seu.}
\talento{Contagem de Corpos}{3}{Para cada aliado Sob Sua Guarda que não tenha sofrido dano neste combate, você recebe +1 na CA, cumulativo.}
\tecnica{Custe o Que Custar}{s}{4}{3}{1 Reação}{Pessoal}{1 Reação, quando um aliado Sob Sua Guarda chegaria a 0 PV: ele fica com 1 PV e é movido 9m para fora do perigo. Você sofre todo o dano excedente, sem receber Marcas da Morte neste turno.}
\tecnica{Bastião Menor}{\relax}{3}{2}{1 Ações}{Esfera de 6m}{1 minuto: criaturas hostis que entrarem gastam o dobro do deslocamento, e nenhum efeito de área de fora atinge quem está dentro sem antes te atingir.}
\patamar{Muralha}{Dado de PV \dado{2d6+6} \textbullet\ +1 degrau(s) de Dado de Arma \textbullet\ +2 PT}
\maestria{Ninguém Passa}{Criaturas hostis não podem se mover para além de você (3m para cada lado da sua linha) sem antes vencer uma disputa de Força ou Vigor. Você intercepta dano por Sob Minha Guarda sem gastar Reação, quantas vezes quiser por rodada. Imune a Paralisia, Petrificação, Preso e efeitos que impeçam agir, com 1+ PT.}
\tecnica{A Linha}{s}{5}{4}{1 Ações}{18 metros}{1 minuto: todo dano dirigido a qualquer aliado a até 18 metros vem para você automaticamente, reduzido pelo seu Bônus de Rank. Você não pode se mover, atacar nem ser curado enquanto sustentar.}
\tecnica{Ordem de Recuo}{\relax}{4}{2}{1 Ações}{18 metros}{Todos os aliados a até 18m movem-se imediatamente até o próprio Deslocamento em direção a um ponto indicado, sem ataques de oportunidade e sem gastar as Ações deles.}
\patamar{Bastião}{Dado de PV \dado{2d8+6} \textbullet\ +2 PT}
\maestria{Enquanto Eu Estiver de Pé}{Nenhum aliado Sob Sua Guarda pode ser reduzido a menos de 1 PV enquanto você estiver consciente e a até 18m dele --- o excedente vem para você, sempre, sem custo. Você recebe 1 Ação adicional (mover-se, interpor-se, proteger). Ao chegar a 0 PV, gaste todos os PT e volte a 1 PV, uma vez por combate.}
\talento{Nome na Porta}{4}{Aliados Sob Sua Guarda ficam imunes a Amedrontado, e inimigos que falharem ao tentar atravessar sua linha ficam Abalados (Desvantagem até o fim do próximo turno).}
\tecnica{O Muro Final}{s}{6}{6}{2 Ações}{Todo o campo de batalha}{Uma vez por Descanso Longo. Por 1 minuto, nenhum aliado seu pode morrer --- todo dano letal é transferido para você, e você não cai abaixo de 1 PV durante a duração. Quando acaba, todo o dano acumulado é aplicado em você de uma vez. Você provavelmente morre.}
\tecnica{Aço Inquebrável}{\relax}{5}{4}{1 Ações}{Pessoal}{Por 3 turnos, você é imune a todo dano de patamar Rei ou inferior, e criaturas hostis a até 9m não conseguem se afastar de você.}

\subsection{Estilo Vendaval}
\begin{arvorecard}{Estilo Vendaval}{Estilo Híbrido}{Força ou Agilidade \textbullet\ Recurso: PT}
O que sobra quando a sobrevivência do Norte para de precisar de distância --- nascido de quem já domina o aço e o vento ao mesmo tempo.
\par\medskip\textbf{Pré-requisito:} Pré-requisito narrativo: Rank Avançado ou superior em Estilo Deus do Norte E em Magia de Vento. Não é uma árvore que se escolhe na criação --- é o que emerge de já ter dominado as outras duas. O jogo não impede a compra de outro jeito; o Mestre nega até os dois pré-requisitos estarem cumpridos, do mesmo jeito que já faz com a Raça Dragão e o Rank Deus.
\end{arvorecard}
\patamar{Principiante}{Dado de PV \dado{1d8+2} \textbullet\ +1 degrau(s) de Dado de Arma}
\maestria{Passo sem Peso}{+3 metros de Deslocamento. Uma vez em cada um dos seus turnos, afastar-se de um inimigo não provoca ataque de oportunidade. Uma vez por combate, ao errar um ataque corpo a corpo, o vento absorve o impacto: reposicione-se até 3m sem gastar Ação.}
\talento{Fôlego do Vendaval}{1}{+2 PT Máximos.}
\talento{Ouvido do Vendaval}{1}{Vantagem em Percepção auditiva, e você percebe magia de Vento sendo conjurada a até 60 metros.}
\talento{Pés que o Vento Segura}{1}{Você nunca sofre penalidade de Deslocamento em terreno difícil, gelo ou vento contrário, e nunca escorrega.}
\tecnica{Corte que Chega Antes}{s}{2}{\relax}{1 Ações}{Corpo a corpo}{\textbf{Dano:} Dado de arma normal\par
Seu golpe viaja com o vento: ataque corpo a corpo que também empurra o alvo 3m na direção do golpe (teste de Força CD 8 + Força + Rank pra resistir só ao empurrão --- o dano acerta de qualquer forma).}
\tecnica{Arremesso Cortante}{\relax}{1}{\relax}{1 Ações}{18 metros}{\textbf{Dano:} Dado de arma normal\par
Ataque à distância até 18m arremessando sua arma como um golpe de vento cortante, usando seu Dado de Arma e Bônus de Rank normalmente. A arma retorna à sua mão ao fim do turno, puxada pelo vento.}
\tecnica{Sopro sob a Lâmina}{\relax}{1}{\relax}{1 Ações}{Cone de 4,5 metros}{Teste de Vigor (CD 8 + Força + Rank) ou o alvo fica Desequilibrado e Cego por poeira levantada até o fim do próximo turno dele. Sem dano.}
\tecnica{Passo Vazio de Combate}{\relax}{1}{\relax}{1 Ações}{Deslocamento}{Desloque-se até seu Deslocamento sem provocar ataques de oportunidade e ignorando terreno difícil. Se terminar adjacente a um inimigo que não te viu chegar, seu próximo ataque tem Vantagem.}
\patamar{Intermediário}{Dado de PV \dado{1d8+3} \textbullet\ +1 degrau(s) de Dado de Arma}
\maestria{Lâmina no Vácuo}{Seus ataques corpo a corpo ignoram metade da Cobertura do alvo --- o vento carrega seu golpe ao redor de obstáculos parciais. Uma vez por turno, gaste 1 PT pra repetir um teste de resistência de Agilidade que tenha falhado.}
\talento{Corrente de Apoio Marcial}{1}{Quando um aliado a até 9m for empurrado, arremessado ou puxado por qualquer efeito, gaste 1 PT como Reação pra dobrar a distância do movimento dele, na direção que ele quiser.}
\talento{Fuga no Vendaval}{1}{Ao ser reduzido a 0 PV, o vento te arremessa 6m para longe da fonte do dano antes de cair.}
\talento{Instinto de Corrente de Ar}{1}{Você não pode ser surpreendido enquanto consciente, e sente a aproximação de qualquer voo hostil a até 30 metros.}
\tecnica{Redemoinho de Aço}{s}{2}{1}{1 Ações}{3 metros}{\textbf{Dano:} Dado de arma normal em cada alvo\par
Gire com sua arma, atacando cada criatura à escolha dentro de 3m, com rolagem separada para cada, sem Desvantagem por aliados na área --- o vento abre espaço pra sua lâmina.}
\tecnica{Golpe que Corta o Ar}{\relax}{1}{1}{1 Ações}{18 metros}{\textbf{Dano:} Dado de arma normal (cortante)\par
Ataque à distância com sua arma corpo a corpo, como se o fio dela se esticasse no vento. Se acertar, o alvo fica Desequilibrado.}
\tecnica{Passo Entre Rajadas}{\relax}{1}{1}{1 Reação}{9 metros}{1 Reação, quando um inimigo errar um ataque contra você: desloque-se até 9m antes que ele perceba. Uma vez por turno.}
\patamar{Avançado}{Dado de PV \dado{1d10+3} \textbullet\ +1 degrau(s) de Dado de Arma \textbullet\ +1 PT}
\maestria{Fio do Vendaval}{Enquanto vestir o Manto de Touki, sua arma ganha +3 metros de alcance em ataques corpo a corpo --- o fio do golpe se estende pelo vento --- e você ignora metade da CA de armaduras não-mágicas.}
\talento{Segunda Rajada}{2}{+3 PT Máximos.}
\talento{Vórtice Pessoal}{2}{Ataques à distância mundanos contra você sofrem Desvantagem enquanto estiver se movendo, e você nunca sofre dano de queda.}
\talento{Leitura do Vento}{2}{Vantagem em Iniciativa, e você identifica de onde veio qualquer ataque à distância que te acertou, mesmo escondido.}
\tecnica{Corte do Horizonte Curto}{s}{3}{2}{1 Ações}{Linha de 12 metros}{\textbf{Dano:} Dado de arma normal + 3d8 (cortante)\par
Libere seu golpe como uma lâmina de vento em linha. Teste de Agilidade (CD 8 + Força + Rank) pra metade do dano extra --- o dano de arma sempre acerta cheio.}
\tecnica{Prisão de Ar e Aço}{\relax}{2}{1}{1 Ações}{9 metros}{Teste de Força (CD 8 + Força + Rank). Falha: o alvo é erguido do chão, Preso e Desequilibrado por 1 minuto, cortado suavemente pelo vento que o segura no lugar (repete teste no fim de cada turno).}
\tecnica{Golpe Sem Peso}{\relax}{2}{\relax}{Passivo}{Passivo}{Seus ataques corpo a corpo não sofrem Desvantagem por empunhar arma pesada com uma mão só, e você nunca fica Desequilibrado pelo próprio ataque.}
\patamar{Santo}{Dado de PV \dado{1d10+4} \textbullet\ +1 degrau(s) de Dado de Arma \textbullet\ +1 PT}
\maestria{Olho da Tempestade}{Em combate, uma bolha de calmaria de 3m ao seu redor nega Cobertura e vento contrário a você e a aliados adjacentes; fora dela, o ar reage aos seus movimentos, e seu Deslocamento aumenta em metade.}
\tecnica{Corte que o Vento Termina}{s}{4}{2}{2 Ações}{Corpo a corpo, depois Linha de 18 metros}{\textbf{Dano:} Dado de arma normal + 5d10 (cortante)\par
Um golpe corpo a corpo que, se acertar, libera uma segunda lâmina de vento em linha atrás do alvo. Teste de Agilidade (CD 8 + Força + Rank) pra metade do dano da segunda lâmina.}
\tecnica{Rajada que Segue a Lâmina}{\relax}{3}{2}{1 Ações}{Esfera de 6 metros}{\textbf{Dano:} 3d10 (contundente)\par
Teste de Força (CD 8 + Força + Rank). Falha: dano, arremessadas 6m e Caídas. Vento forte na área por 1 minuto --- projéteis mundanos erram.}
\tecnica{Fôlego Infinito do Vendaval}{\relax}{3}{\relax}{Passivo}{Passivo}{Você ignora fome, sede, clima extremo e a Exaustão por falta de sono. Enquanto tiver ao menos 1 PT, é imune a Cego e Surdo.}
\patamar{Rei}{Dado de PV \dado{1d12+4} \textbullet\ +1 degrau(s) de Dado de Arma \textbullet\ +1 PT}
\maestria{Lâmina Sem Fronteira}{Seus ataques corpo a corpo passam a ter alcance de 6 metros, sempre --- o vento carrega o fio da sua arma até o alvo. A Vantagem de Estilo (Cap. 3) se aplica mesmo à distância.}
\tecnica{Mil Cortes no Vendaval}{s}{5}{3}{1 Ações}{Esfera de 9 metros}{\textbf{Dano:} Dado de arma rolado duas vezes + Bônus de Rank (cortante)\par
Ataque cada criatura à sua escolha dentro de 9m, com rolagem separada para cada, sem Desvantagem por aliados na área --- o vento não escolhe lado, você escolhe por ele.}
\tecnica{Golpe que Não Tem Origem}{\relax}{4}{2}{1 Ações}{27 metros}{\textbf{Dano:} Dado de arma normal + 4d8\par
Ataque à distância com sua arma, sem precisar de linha de visão direta --- o vento contorna obstáculos parciais. Ignora Cobertura parcial.}
\tecnica{Pele que o Vento Protege}{\relax}{4}{\relax}{Passivo}{Passivo}{Enquanto vestir o Manto de Touki, você tem Resistência a dano de projétil mundano, e ataques à distância contra você têm Desvantagem.}
\patamar{Imperador}{Dado de PV \dado{1d12+5} \textbullet\ +1 degrau(s) de Dado de Arma \textbullet\ +1 PT}
\maestria{O Golpe que Já Passou}{Uma vez por combate, declare que um golpe seu já aconteceu antes da cena começar: ele acerta automaticamente e ignora Touki, Cobertura e armadura. Enquanto tiver ao menos 1 PT, você não pode ser Agarrado, Preso nem Desequilibrado por vento ou empurrão de qualquer fonte.}
\tecnica{Corte do Horizonte Infinito}{s}{6}{4}{3 Ações}{Linha de 300 metros}{\textbf{Dano:} Dado de arma rolado quatro vezes + Força + Bônus de Rank + 4d10 (cortante)\par
Uma vez por combate. Acerta automaticamente; o dano de arma e o bônus sempre acertam cheio. Muralhas, torres e florestas na trajetória são cortadas ao meio. Teste de Agilidade com Desvantagem pra metade do dano extra do vento (4d10).}
\tecnica{Vendaval de Aço}{\relax}{5}{3}{1 Ações}{Esfera de 18 metros}{\textbf{Dano:} 6d10 (cortante)\par
Teste de Força com Desvantagem (CD 8 + Força + Rank). Falha: dano cheio, arremessadas 9m, Caídas e Desequilibradas. Sucesso: metade do dano, sem os outros efeitos.}
\tecnica{Presença que Corta o Ar}{\relax}{5}{\relax}{1 Ações}{18 metros}{1 minuto: aliados em 18m recebem +2 em acertos e Deslocamento dobrado. Inimigos em 18m sofrem Desvantagem em ataques à distância contra qualquer um do seu grupo.}

\subsection{Arquearia}
\begin{arvorecard}{Arquearia}{Arqueiro}{Agilidade \textbullet\ Recurso: PT}
Contra quem não veste Touki, o maior dano sustentado do jogo, a 90 metros e sem gastar recurso. Contra rank Santo+, o Manto de Touki reduz cada tiro pelo dobro do Bônus de Rank do alvo --- a árvore concede só uma forma cara de furar isso.
\end{arvorecard}
\patamar{Atirador}{Dado de PV \dado{1d8+2} \textbullet\ +1 degrau(s) de Dado de Arma}
\maestria{Olho do Caçador}{Você nunca sofre Desvantagem por distância longa. Ignora Cobertura Leve e não sofre penalidade por atirar em meio a aliados engajados. Identifica a distância exata até qualquer coisa visível.}
\talento{Aljava Cheia}{1}{Você nunca fica sem flechas em terreno com madeira, e fabrica munição durante um Descanso Curto.}
\talento{Passo e Tiro}{1}{Disparar não provoca ataques de oportunidade.}
\talento{Braço Firme}{1}{+4 PV por patamar seu nesta árvore. Aplicado sozinho na ficha, e cresce a cada patamar novo que você abrir nela.}
\tecnica{Disparo Duplo}{s}{2}{\relax}{1 Ações}{Alcance da arma}{\textbf{Dano:} Segundo disparo: um degrau abaixo\par
Dois disparos contra o mesmo alvo ou alvos diferentes. Se ambos acertarem o mesmo alvo, ele fica Marcado; e faz um teste de Vigor (CD 8 + Agilidade + Bônus de Rank) ou também fica Caído.}
\tecnica{Tiro de Contenção}{\relax}{1}{\relax}{1 Ações}{Alcance da arma}{\textbf{Dano:} Metade do dado de arma\par
Teste de Vigor (CD 8 + Agilidade + Rank) ou o alvo tem o Deslocamento reduzido à metade por 2 turnos.}
\tecnica{Tiro de Objeto}{\relax}{1}{\relax}{1 Ações}{Qualquer distância visível}{Você acerta um objeto específico e pequeno a qualquer distância visível, sem rolagem.}
\tecnica{Flecha de Sinal}{\relax}{1}{\relax}{Passivo}{Passivo}{Você carrega flechas de assobio, de fumaça colorida e incendiárias.}
\patamar{Caçador}{Dado de PV \dado{1d8+2} \textbullet\ +1 degrau(s) de Dado de Arma}
\maestria{A Marca Fica}{A condição Marcado passa a durar o combate inteiro. Você pode Marcar sem atacar (1 Ação de observação). Contra criaturas Marcadas, sabe automaticamente PV aproximado, resistências e se ela veste Touki.}
\talento{Leitura de Presa}{1}{Vantagem em Sobrevivência e Percepção para rastrear e emboscar. (Se você já tem isto pelo Tático ou pelo Norte, escolha outro talento --- o mesmo bônus não empilha.)}
\talento{Corda Rápida}{1}{Recarregar besta deixa de custar Ação.}
\talento{Distância É Segurança}{1}{Contra criaturas a mais de 18 metros, você recebe +2 na CA.}
\tecnica{Tiro Certeiro}{s}{2}{1}{1 Ações}{Alcance da arma}{Um disparo contra alvo Marcado que acerta automaticamente e crita em 19-20.}
\tecnica{Chuva de Flechas}{\relax}{1}{1}{1 Ações}{Esfera de 6m ao alcance máximo}{Teste de Agilidade (CD 8 + Agilidade + Rank) ou sofrem o dano de arma completo. Ignora Cobertura que não seja teto.}
\tecnica{Armadilha de Caça}{\relax}{1}{\relax}{Passivo}{Local (10 minutos para montar)}{\textbf{Dano:} 3d6\par
Uma criatura que entrar sofre o dano e fica Presa até passar num teste de Força. Você pode ter até Agilidade armadilhas ativas.}
\tecnica{Tiro Perfurante}{\relax}{1}{1}{1 Ações}{Linha}{O disparo atravessa o alvo e atinge a próxima criatura na linha com o mesmo dano. Contra alvo Atolado, Preso ou Congelado, atravessa até três.}
\patamar{Franco-Atirador}{Dado de PV \dado{1d8+3} \textbullet\ +1 degrau(s) de Dado de Arma \textbullet\ +1 PT}
\maestria{Aura na Corda}{Você recebe o Manto de Touki completo e a reserva de PT. Seus disparos contam como mágicos e ignoram Resistência a perfurante --- mas ainda não furam o Manto de Touki de ninguém. O alcance de todas as suas armas de disparo dobra.}
\talento{Fôlego Estável}{2}{+3 PT Máximos.}
\talento{Três na Corda}{2}{Disparo Duplo passa a ser triplo, com o terceiro disparo dois degraus abaixo.}
\talento{Nunca Aqui}{2}{Depois de atirar, gaste 1 PT para se mover 9m sem provocar oportunidade e refazer Furtividade imediatamente.}
\tecnica{Tiro do Céu}{s}{3}{2}{2 Ações}{Dobro do alcance máximo}{\textbf{Dano:} Dado de arma rolado três vezes\par
Se o alvo estiver Marcado e não souber que você existe, o acerto é crítico automático.}
\tecnica{Tiro Interrompido}{\relax}{2}{1}{1 Reação}{Alcance da arma}{1 Reação, quando uma criatura visível começar a conjurar, beber uma poção ou usar um item. Disparo automático; se causar dano, teste de Espírito ou perde a ação e o recurso gasto.}
\tecnica{Ninho}{\relax}{2}{\relax}{Passivo}{Posição preparada (10 minutos)}{Enquanto permanecer na posição: Vantagem em todos os disparos, +2 no dano, e ninguém determina sua posição sem Percepção com Desvantagem contra sua Furtividade.}
\patamar{Olho de Águia}{Dado de PV \dado{1d10+3} \textbullet\ +1 degrau(s) de Dado de Arma \textbullet\ +1 PT}
\maestria{A Distância Não Existe}{Você acerta qualquer alvo que consiga ver, sem limite de alcance. Enxerga com precisão perfeita a até 1 km, atravessando neblina, escuridão e chuva. Contra alvos Marcados, seus disparos ignoram Cobertura Total, desde que exista qualquer trajetória física.}
\talento{Marca Perene}{3}{A condição Marcado passa a durar até o próximo Descanso Longo, mesmo se a criatura fugir, se esconder ou atravessar o continente. Você a encontra de novo.}
\talento{Contra-Bateria}{3}{Quando um inimigo te atacar à distância, você sabe exatamente onde ele está e o Marca automaticamente.}
\talento{Peso da Aljava}{3}{Suas flechas causam +1d8 contra criaturas Grandes ou maiores.}
\tecnica{Um Alvo, Um Tiro}{s}{4}{3}{2 Ações}{Alcance ilimitado (visão)}{Uma vez por combate. Contra alvo Marcado com metade ou menos dos PV: o disparo reduz o alvo a 0 PV automaticamente, sem rolagem. Não funciona contra quem veste Touki, rank Rei ou superior, nem mais de 200 PV máximos.}
\tecnica{Tempestade de Setas}{\relax}{3}{2}{1 Ações}{30 metros}{Você dispara contra todas as criaturas hostis que consiga ver dentro de 30m, com rolagem separada para cada. Dano completo em todas.}
\tecnica{Flecha Amarrada}{\relax}{3}{1}{1 Ações}{Alcance da arma}{Teste de Força (CD 8 + Agilidade + Rank) ou o alvo fica Preso, com Deslocamento 0, até gastar 1 Ação e passar no teste.}
\patamar{Predador}{Dado de PV \dado{1d10+4} \textbullet\ +1 degrau(s) de Dado de Arma \textbullet\ +1 PT}
\maestria{A Flecha que Fura}{O patamar que o arqueiro esperou a campanha inteira: gastando 3 PT (Flecha de Touki), um único disparo ignora completamente o Manto de Touki do alvo e toda redução de dano contra projéteis. É caro de propósito, e não existe segunda forma.}
\talento{Aljava Divina}{4}{A Flecha de Touki passa a custar 2 PT em vez de 3.}
\tecnica{Caçada}{s}{5}{2}{1 Ações}{Alcance da arma}{Escolha um alvo Marcado. Até o fim do combate, seus disparos contra ele têm Vantagem e critam em 18-20; e o primeiro disparo de cada turno que acertar causa +1d10 cumulativo, até o máximo de +3d10. Trocar de alvo encerra o efeito.}
\tecnica{Flecha do Fim da Estrada}{\relax}{4}{4}{1 Reação}{Visão}{1 Reação, quando uma criatura visível tentar fugir, teleportar, voar para fora de alcance ou entrar por uma porta. Disparo automático com Flecha de Touki embutida; teste de Vigor ou a fuga falha.}
\patamar{Lenda da Flecha}{Dado de PV \dado{1d12+4} \textbullet\ +1 degrau(s) de Dado de Arma \textbullet\ +1 PT}
\maestria{O Tiro Que Já Aconteceu}{Você recebe 1 Ação adicional por turno, usável apenas para um único disparo --- nunca para Disparo Duplo, Três na Corda ou qualquer técnica nomeada. Seus disparos não podem ser aparados, bloqueados nem interceptados por efeito de rank Rei ou inferior, incluindo o Fluxo Verdadeiro. Criaturas Marcadas por você não conseguem se esconder de você, nem por magia.}
\talento{Nenhum Deles Chegou Perto}{4}{Enquanto tiver ao menos 1 PT e linha de visão, criaturas hostis não conseguem se aproximar a menos de 9 metros sem passar num teste de Espírito (CD 8 + Agilidade + Bônus de Rank).}
\tecnica{A Flecha que Não Erra}{s}{6}{5}{2 Ações}{Qualquer lugar já visto}{\textbf{Dano:} Dado de arma rolado cinco vezes + Agilidade + Bônus de Rank\par
Uma vez por combate. Acerta automaticamente, ignora CA, Cobertura, Manto de Touki, armadura mágica e barreiras físicas. Se o alvo estiver Marcado, não precisa vê-lo agora.}
\tecnica{Céu Cheio}{\relax}{5}{4}{2 Ações}{Esfera de 30m a qualquer distância}{\textbf{Dano:} 12d10 (perfurante)\par
Teste de Agilidade para metade. Criaturas Marcadas na área não têm direito ao teste. A área vira terreno difícil eriçado de hastes.}

\section{Árvore de Utilidade}

\subsection{Furtividade e Armadilhas}
\begin{arvorecard}{Furtividade e Armadilhas}{Batedor e Ladrão}{Agilidade \textbullet\ Recurso: PP}
Domínio da Preparação: coisas e lugares. Faixa exclusiva: Dano Furtivo --- só o Ladino causa dano acima do trivial nesta árvore. A pergunta dele: "como eu entro?"
\end{arvorecard}
\patamar{Gatuno}{Dado de PV \dado{1d6+1}}
\maestria{Olho Treinado e Dano Furtivo}{Escopo: um objeto, um cômodo, uma pessoa comum. Armadilhas e gatilhos não pedem teste para serem percebidos --- você simplesmente os vê. Uma vez por turno, ao acertar um alvo desprevenido, cego, imobilizado ou contra o qual tenha Vantagem, some +1d6 de Dano Furtivo por patamar que possua nesta árvore.}
\talento{Mãos Rápidas}{1}{Você tira e coloca objetos em bolsos alheios com teste de Agilidade contra a Percepção do alvo. Em combate, 1 Ação para roubar item não empunhado.}
\talento{Engenhoca}{1}{Você desarma e constrói armadilhas. Montar leva 10 minutos; dano modesto (2d6), CD pra perceber 8 + Agilidade + Bônus de Rank.}
\talento{Duelista de Rua}{1}{Com arma leve, +2 na CA contra um oponente por turno, e não provoca oportunidade ao se afastar dele.}
\talento{Nunca Preso}{1}{Você escapa de qualquer contenção não-mágica com 1 Ação e Vantagem Absoluta.}
\talento{Boticário}{1}{Venenos e sedativos de campo. O alvo faz teste de Vigor (CD 8 + Intelecto + Bônus de Rank) ou fica Envenenado por 1 minuto.}
\talento{Falsário}{1}{Selos, brasões, cartas de crédito. Detectar a falsificação exige CD 8 + Intelecto + Bônus de Rank.}
\tecnica{Primeiro Golpe}{s}{2}{\relax}{1 Ações}{Corpo a corpo}{\textbf{Dano:} Dano Furtivo triplicado\par
Uma vez por combate. Requer que o alvo ainda não tenha agido, ou não saiba onde você está. Depois de usar, você é só alguém com uma faca.}
\patamar{Sombra}{Dado de PV \dado{1d6+2}}
\maestria{Duas Saídas}{Escopo: um edifício inteiro, uma rotina de trabalho, um pequeno grupo. Você identifica automaticamente todas as saídas de qualquer ambiente, incluindo improvisadas. Uma vez por cena, gastando 1 PP, declare que existe uma saída onde o Mestre não tinha planejado.}
\talento{Mapa dos Ratos}{1}{Em cidade onde já passou um dia, conhece a geografia oculta: esgotos, becos, telhados, casas seguras --- sem teste.}
\talento{Leitura de Cena}{1}{Você reconstrói o que aconteceu num ambiente fechado só olhando: móveis arrastados, o que foi levado, quantas pessoas estiveram.}
\talento{Contrabandista}{1}{Você vende qualquer item incriminador e compra itens que não estão à venda --- por preço alto e favores.}
\talento{Dedos de Mana}{1}{Requer 1 patamar em escola de magia. Conjure magias de rank Principiante sem cântico nem gesto visível, ao custo da versão Encurtada.}
\talento{Sombra Longa}{1}{Você se esconde mesmo observado, com qualquer distração ou obstáculo parcial. Escuridão total dá Vantagem Absoluta em Furtividade.}
\talento{Passo de Gato}{1}{Você se move em velocidade normal sem ruído algum, e escalar custa deslocamento normal.}
\tecnica{Passo Vazio}{s}{2}{\relax}{1 Ações}{Pessoal}{Uma vez por combate: você sai do combate --- não é alvo válido, não é atingido por área, ninguém determina sua posição, até você agir, tocar em alguém ou o combate acabar. Ao reaparecer, sua primeira ação tem Vantagem e reativa o Primeiro Golpe.}
\patamar{Especialista}{Dado de PV \dado{1d8+2} \textbullet\ +1 PP}
\maestria{O Nome Certo}{Escopo: um quarteirão, uma guilda local, um cofre de banco. Para qualquer informação dentro do seu Escopo, você não investiga: nomeia a pessoa que a possui, e o Mestre confirma que ela existe.}
\talento{Quem Puxa as Cordas}{2}{Teste de Enganação estendido ao longo de dias que planta uma decisão na cabeça do alvo, que jura tê-la pensado sozinho.}
\talento{Mestre-Chave}{2}{Fechaduras e cofres mundanos viram questão de tempo, não teste. Barreiras mágicas ainda te barram.}
\talento{Veneno Refinado}{2}{Requer Boticário. Seus venenos impõem Desvantagem no teste e podem ser calibrados: sono, paralisia parcial, mudez, febre retardada.}
\talento{Marca da Casa}{2}{Estude uma organização por uma semana: hierarquia, senhas, uniformes. Vantagem em testes sociais e de Furtividade contra membros dela.}
\tecnica{Ponto Cego}{s}{3}{1}{1 Ações}{Ambiente do combate}{Declare uma sabotagem feita neste ambiente antes do combate: o lustre cai (4d6, Caído, 3m); a porta dos reforços está pregada (perdem 3 turnos); o chão está encharcado de óleo; a arma de um inimigo foi limada (quebra no crítico ou falha crítica).}
\patamar{Mestre Espião}{Dado de PV \dado{1d8+3} \textbullet\ +1 PP}
\maestria{Segunda Face}{Escopo: uma cidade inteira, uma casa nobre menor, uma rota comercial. Você mantém uma identidade falsa completa e documentada; magia de detecção de mentiras não te denuncia, porque no momento em que fala, você é aquela pessoa. Mantém identidades = seu Intelecto.}
\talento{Homem Dentro}{2}{Escolha uma organização urbana: você tem um agente permanente lá dentro, que age a seu favor uma vez por sessão.}
\talento{A Faca do Amigo}{2}{1 Ação e 1 PP: um inimigo secundário na cena trabalha pra você --- age uma vez, depois foge ou morre.}
\talento{Nada Escrito}{2}{Você memoriza documentos com uma leitura e escreve em cifras que exigem Intelecto igual ou superior pra quebrar.}
\talento{Faca no Escuro}{2}{Seu Dano Furtivo funciona contra alvos adjacentes a um aliado que tenha atacado o alvo desde o seu último turno.}
\tecnica{A Faca Certa}{s}{3}{1}{1 Ações}{Visão}{Declare que já sabotou o equipamento de um inimigo visível que você já viu antes desta cena. Escolha um: perde todo bônus de armadura por 1 minuto; a arma quebra no próximo ataque; não usa itens por 3 turnos.}
\patamar{Fantasma}{Dado de PV \dado{1d10+3} \textbullet\ +1 PP}
\maestria{Não Estive Aqui}{Escopo: uma capital, um palácio, uma organização continental. Nenhuma evidência física da sua presença persiste mais de uma hora. Magia de rastreamento, adivinhação e visão do passado não te encontra. Testemunhas descrevem você de forma contraditória.}
\talento{Mão na Coroa}{3}{Seu Escopo inclui uma corte real ou organização continental fixa; dentro dela, seus fatos custam 1 PP a menos (mínimo 1).}
\talento{O Dossiê}{3}{Sobre qualquer indivíduo nomeado que já tenha encontrado, você possui um segredo comprometedor. Gaste 2 PP para revelá-lo.}
\talento{Saída Limpa}{3}{Uma vez por Descanso Longo, você e aliados a 9m simplesmente não estão mais lá --- reaparecem num lugar seguro já conhecido.}
\tecnica{Você Não Vai Chegar Lá}{s}{4}{2}{1 Reação}{Visão}{1 Reação, quando um inimigo declara que vai alcançar um objetivo. Ele não chega --- declare o motivo dentro do seu Escopo (ponte serrada, chave sumida, corredor murado). Não causa dano nenhum e frequentemente ganha o combate.}
\patamar{Lenda Oculta}{Dado de PV \dado{1d10+4} \textbullet\ +1 PP}
\maestria{O Fato Consumado}{Escopo: um reino, um continente, o registro histórico. Seus fatos podem ter sido executados por outra pessoa, anos atrás, a seu mando. Uma vez por Descanso Longo, gastando 4 PP, declare que a situação atual foi arranjada por você (o Mestre escolhe um detalhe que saiu do controle). Recupere 2 PP em Descanso Curto.}
\talento{Nome que Não Existe}{3}{Você apaga uma pessoa dos registros do mundo, ou insere uma. Leva meses e é irreversível.}
\talento{A Mão Longa}{3}{Seu Escopo passa a cobrir um continente inteiro.}
\talento{Herança}{3}{Escolha um talento de qualquer patamar desta árvore que não possua --- alguém do seu séquito sabe fazer aquilo. Troque a cada Descanso Longo; nunca mais de um por vez.}
\tecnica{O Homem Que Nunca Esteve Lá}{s}{4}{4}{3 Ações}{Visão}{Uma vez por Descanso Longo. Declare que o inimigo à sua frente já perdeu: os aliados dele nunca foram dele; o que veio buscar não está mais aqui; a autoridade dele foi revogada; ou ele está sozinho e a saída está fechada. Não funciona contra feras, mortos-vivos, elementais ou rank Deus.}

\subsection{Bardo e Interação}
\begin{arvorecard}{Bardo e Interação}{Bardo}{Espírito \textbullet\ Recurso: PP}
Domínio da Preparação: pessoas e reputação. Faixa exclusiva: estado emocional --- só o Bardo altera o que um inimigo sente. A pergunta dele: "quem eu convenço?"
\end{arvorecard}
\patamar{Aprendiz}{Dado de PV \dado{1d6+1}}
\maestria{A Plateia}{Escopo: uma pessoa que já te ouviu tocar ou falar. Enquanto estiver tocando/cantando/falando (sem custo de Ação fora de combate), aliados que te ouvem somam seu Bônus de Rank em um teste de perícia por cena, à escolha deles. Você nunca dorme na rua --- uma apresentação garante cama e comida.}
\talento{Ouvido Absoluto}{1}{Você imita qualquer voz já ouvida e reproduz sotaques. Aprende idiomas em dias.}
\talento{Cantiga de Marcha}{1}{O grupo viaja mais rápido e ignora o primeiro nível de Exaustão por marcha, enquanto você tocar.}
\talento{Insulto Afiado}{1}{1 Ação: teste de Espírito contra Intuição do alvo. Se vencer, o próximo ataque dele tem Desvantagem.}
\talento{Colecionador de Histórias}{1}{Sobre pessoa/família/cidade/artefato conhecidos, você sabe uma coisa verdadeira e uma exagerada, e distingue qual é qual.}
\talento{Contrato de Bardo}{1}{Uma vez por mês de jogo, um patrono cobre as despesas do grupo em troca de você registrar os feitos deles.}
\tecnica{Inspiração}{s}{2}{\relax}{1 Ações}{Voz}{\textbf{Dano:} Dado de Inspiração: 1d6 (1º-2º patamar), 2d6 (3º-4º), 3d6 (5º-6º)\par
Número de vezes por Descanso Longo igual ao seu Espírito. Um aliado que te ouça recebe um Dado de Inspiração, somável a qualquer teste até o fim da cena, mesmo após ver o resultado. O Dado de Inspiração e a Maestria A Plateia não contam para o Teto de Auxílio +5 (Cap. 4, §4).}
\patamar{Artista}{Dado de PV \dado{1d6+2}}
\maestria{Ler a Sala}{Escopo: uma taverna, uma tropa, um público, uma família. Ao entrar em qualquer cena com pessoas, o Mestre deve dizer, sobre cada indivíduo relevante: o que ele quer, do que tem medo, e que existe algo que ele esconde (não qual).}
\talento{Cantiga de Ninar}{1}{Criaturas não-hostis que te ouvirem 10 minutos adormecem, salvo teste de Vigor. Não funciona em combate.}
\talento{Mestre de Cerimônias}{1}{Você controla uma multidão: acalmar, iniciar tumulto ou direcionar a atenção de todos. Teste de Espírito com Vantagem contra a CD do Mestre; num sucesso a multidão age como você quer por 1 minuto, e você escolhe uma pessoa nela que não é afetada.}
\talento{A Máscara}{1}{Enquanto sustentar uma persona, testes pra detectar mentira em você têm Desvantagem. Dura uma noite.}
\talento{Nome nas Bocas}{1}{Gastando 1 PP, planta um rumor numa comunidade. Em três dias, todo mundo acredita --- não precisa ser verdade.}
\talento{Duelo de Canções}{1}{Você desafia alguém pra disputa artística/verbal; recusar é desonra pública. Vença uma disputa de Atuação e o perdedor fica Amedrontado por você até o fim da cena, e ninguém que assistiu testemunha contra você naquela comunidade.}
\tecnica{Insulto que Fica}{s}{2}{\relax}{1 Ações}{Voz}{Teste de Espírito contra Espírito do alvo. Se vencer, por 3 turnos ele só consegue pensar em você: Desvantagem em ataques que não sejam contra você, e não pode usar habilidades que exijam concentração ou cálculo. O risco: ele vai te atacar.}
\patamar{Trovador}{Dado de PV \dado{1d6+2} \textbullet\ +1 PP}
\maestria{A Canção Não Para}{Escopo: um vilarejo, uma companhia mercenária, uma corte pequena. Você sustenta um efeito de Bardo indefinidamente sem gastar Ação nem concentração. Segurar dois ao mesmo tempo exige 1 Ação por turno.}
\talento{Réquiem}{2}{Aliados que te ouvem ficam imunes a Amedrontado e têm Vantagem contra efeitos que manipulem emoção ou mente.}
\talento{Diplomata de Guerra}{2}{Você negocia trégua no meio de um combate: teste de Espírito (CD 8 + Espírito + Bônus de Rank), quem falhar para de lutar por 1 minuto e escuta.}
\talento{Voz que Alcança}{2}{Sua voz é ouvida claramente a até 300 metros, atravessa tempestade e ruído de batalha.}
\talento{A Balada Instrutiva}{2}{Você transforma informação complexa em canção memorizável permanentemente pelo grupo em 10 minutos.}
\tecnica{Canção de Guerra}{s}{3}{\relax}{1 Ações}{Voz}{Sustentada de graça pela sua Maestria. Aliados que te ouvem recebem +2 em acertos, imunidade a Amedrontado, e ignoram a penalidade do primeiro nível de Exaustão. Acaba se você for silenciado, nocauteado ou morto.}
\patamar{Virtuoso}{Dado de PV \dado{1d8+3} \textbullet\ +1 PP}
\maestria{Precede Você}{Escopo: uma cidade inteira. Ao entrar em qualquer cidade dentro do seu Escopo pela primeira vez, escolha como é conhecido lá (herói, artista, excêntrico, agente estrangeiro) --- a escolha é verdade. Muda uma vez por cidade, levando semanas.}
\talento{O Favor Antigo}{2}{Gastando 2 PP, uma pessoa importante na cena te deve algo --- dá uma informação, passagem, ou benefício da dúvida.}
\talento{Elegia}{2}{Uma vez por combate, cante para um inimigo que perdeu aliados nesta luta: teste de Espírito com Desvantagem ou ele deixa o combate, sem morrer.}
\talento{A Corte na Palma}{2}{Uma vez por cena, gastando 1 PP: em ambiente formal, você não rola o teste social --- declara o resultado desejado, desde que não contrarie interesses vitais de alguém presente.}
\tecnica{A Verdade que Dói}{s}{3}{1}{1 Ações}{Voz}{Diga, na frente de todos, algo verdadeiro sobre um inimigo (dentro do seu Escopo e Domínio). Aliados dele fazem teste de Espírito (CD 8 + Espírito + Bônus de Rank): quem falhar age com Desvantagem enquanto o alvo estiver na cena.}
\patamar{Maestro}{Dado de PV \dado{1d8+3} \textbullet\ +1 PP}
\maestria{Voz que Comanda}{Escopo: um reino. Qualquer criatura capaz de ouvir e sentir emoção é afetada pelas suas habilidades, mesmo sem idioma comum. Seus efeitos de área social atingem todos que te ouvem, sem limite de número. Você é imune a manipulação da própria emoção.}
\talento{A Canção que Todos Sabem}{3}{Uma composição sua se espalhou. Gastando 2 PP, declare que ela contém um sinal que a pessoa certa reconhece ao ouvir.}
\talento{Silêncio Absoluto}{3}{1 Ação: ninguém em 18m fala/canta/recita, incluindo você. Enquanto sustentar, você não pode usar outra habilidade de Bardo.}
\talento{Herdeiro de Todas as Bocas}{3}{Você fala todos os idiomas do Mundo de Seis Faces, incluindo o Divino e dialetos demoníacos antigos.}
\tecnica{Coro}{s}{4}{2}{1 Ações}{Esfera de 18m}{Uma vez por combate. Escolha pavor, fúria ou devoção. Hostis na área que te ouçam fazem teste de Espírito: quem falhar foge (pavor), ataca o mais próximo (fúria), ou não ataca você/aliados (devoção) por 2 turnos. Não funciona em criaturas sem emoção.}
\patamar{Voz do Mundo}{Dado de PV \dado{1d10+4} \textbullet\ +1 PP}
\maestria{A História Oficial}{Escopo: um continente. Uma vez por Descanso Longo, gastando 4 PP, escolha um evento que testemunhou ou do qual participou: sua versão dele vira a verdade aceita. Desmentir exige provas materiais e testemunha de reputação equivalente. Recupere 2 PP em Descanso Curto.}
\talento{Nome Imortal}{3}{Escolha uma pessoa: ela entra pra história como herói ou monstro, permanentemente.}
\talento{A Marcha}{3}{Uma canção sua vira hino de um movimento. Gastando 3 PP, declare que ele age agora a seu favor.}
\talento{Público Universal}{3}{Suas Preparações sociais alcançam qualquer continente, mesmo os que você nunca visitou.}
\tecnica{O Fim da Canção}{s}{4}{4}{3 Ações}{Voz}{Uma vez por Descanso Longo. Encerre a batalha declarando publicamente por que ela não faz mais sentido. Hostis capazes de ouvir e raciocinar fazem teste de Espírito com Desvantagem; quem falhar encerra as hostilidades. Exige que a razão seja real e acessível --- não funciona em quem luta por prazer, fome ou ordem divina.}

\subsection{Navegação e Liderança}
\begin{arvorecard}{Navegação e Liderança}{Sobrevivência e Táticas}{Intelecto \textbullet\ Recurso: PP}
Domínio da Preparação: tempo e logística. Faixa exclusiva: economia de ação --- só o Tático concede Ações, mexe na Iniciativa e reposiciona aliados. A pergunta dele: "onde e quando isso acontece?"
\end{arvorecard}
\patamar{Explorador}{Dado de PV \dado{1d8+1}}
\maestria{Onde Pisar}{Escopo: a próxima hora, o trecho de estrada à frente. Enquanto liderar a marcha, o grupo nunca se perde e ignora terreno difícil natural. O grupo nunca é surpreendido --- emboscadas ainda acontecem, mas vocês agem no primeiro turno. Sempre encontram água, abrigo e um lugar defensável.}
\talento{Mapa Vivo}{1}{Você desenha e lê mapas; regiões que já atravessou ficam registradas e podem ser vendidas.}
\talento{Suprimento}{1}{O grupo consome metade de ração/água/forragem, e você sempre sabe quantos dias faltam para o problema começar.}
\talento{Sinais}{1}{Código de gestos e assobios com o grupo: comunicação a 200m sem falar.}
\talento{Conhecimento de Bestas}{1}{Sobre qualquer monstro visto, identifica espécie, comportamento de caça e uma fraqueza real.}
\talento{Olho de Cerco}{1}{Olhando pra uma fortificação/acampamento, estima defensores, suprimento, tempo de resistência e o ponto fraco.}
\tecnica{Primeiro a Ver}{s}{2}{\relax}{Passivo}{Passivo}{Uma vez por combate: se o grupo entrar em combate vindo de uma marcha conduzida por você, todos os aliados somam seu Bônus de Rank na Iniciativa, e você escolhe quem age primeiro.}
\patamar{Rastreador}{Dado de PV \dado{1d8+2}}
\maestria{Quem Passou Por Aqui}{Escopo: o dia de hoje, uma região de um dia de viagem. Rastros deixam de ser teste: você sabe quantos eram, o que carregavam, há quanto tempo, e se estavam com pressa. Prevê o destino provável de uma trilha seguida por uma hora.}
\talento{Marcha Forçada}{1}{O grupo viaja o dobro da distância por dia; seu teste de Vigor final tem Vantagem, e você isenta um aliado por dia.}
\talento{Terreno Conhecido}{1}{Escolha um tipo de terreno: velocidade total nele, e Vantagem em navegação e ocultação.}
\talento{Cavalaria}{1}{Você treina e monta qualquer besta de carga; não cai por efeito que permita teste de Agilidade.}
\talento{Retirada Ordenada}{1}{1 Ação: até o fim do próximo turno, aliados que se afastarem de inimigos não provocam oportunidade.}
\tecnica{Antecipação}{s}{2}{1}{1 Reação}{Voz}{1 Reação, quando um inimigo declara uma ação. Um aliado à sua escolha usa imediatamente 1 Ação para reagir. Uma vez por combate.}
\patamar{Guia}{Dado de PV \dado{1d8+2} \textbullet\ +1 PP}
\maestria{O Terreno Escolhido}{Escopo: a semana, uma província, uma rota comercial. Com 10 minutos antes de um combate previsto, você escolhe uma característica definidora do campo de batalha (elevação, gargalo, chão instável, luz favorável, cobertura pesada) --- e o Mestre não pode discordar.}
\talento{Engenharia de Campo}{2}{Com uma hora e ajuda do grupo, ergue paliçada, ponte improvisada, fosso ou trincheira: 40 PV, Cobertura Superior.}
\talento{Voz de Sargento}{2}{1 Ação: um aliado remove Atordoado ou Caído e repete um teste de resistência falho.}
\talento{Logística de Guerra}{2}{Você abastece até cinquenta pessoas indefinidamente em território hostil.}
\talento{Emboscada Planejada}{2}{Requer 1 PP ao usar. Com uma hora de preparo, todos os aliados agem antes de qualquer inimigo no primeiro turno, e os inimigos ficam Surpresos.}
\tecnica{Ponto de Estrangulamento}{s}{3}{1}{1 Ações}{Campo de batalha}{Declare que existe um gargalo de 3m de largura no campo de batalha. Inimigos só atravessam por ali; aliados posicionados nele têm +3 CA e Vantagem em ataques de oportunidade; efeitos de área inimigos atingem no máximo dois personagens.}
\patamar{Estrategista}{Dado de PV \dado{1d10+3} \textbullet\ +1 PP}
\maestria{A Ordem de Batalha}{Escopo: o mês, um campo de batalha inteiro, o suprimento de uma tropa. Aliados podem usar sua rolagem de Iniciativa. Uma vez por turno, sem gastar Ação, troque a posição de dois aliados na ordem de Iniciativa. Você sempre sabe quem ainda não agiu e o que cada inimigo fez.}
\talento{Foco de Fogo}{2}{1 Ação: até o fim do turno, aliados que atacarem o alvo apontado somam seu Bônus de Rank ao dano.}
\talento{Prever o Golpe}{2}{1 Reação: quando aliado a 18m for atingido, ele recebe +4 na CA contra aquele ataque, resolvido retroativamente.}
\talento{A Guerra Antes da Guerra}{2}{Uma vez por Descanso Longo, gastando 2 PP: a força inimiga chegou em pior estado --- reduza o número de inimigos em um terço, ou dê a todos 1 nível de Exaustão.}
\talento{Doutrina}{2}{Escolha um estilo/escola: contra praticantes dele, aliados que ouvirem suas instruções recebem +2 na CA.}
\tecnica{Manobra}{s}{3}{1}{1 Ações}{Voz}{Até três aliados que te ouçam movem-se imediatamente até o próprio Deslocamento, sem provocar oportunidade e sem gastar as Ações deles.}
\patamar{Comandante}{Dado de PV \dado{1d10+3} \textbullet\ +1 PP}
\maestria{Comando}{Escopo: a estação, uma campanha militar. Uma vez por turno, gastando 1 Ação sua, conceda 1 Ação adicional a um aliado que a use imediatamente. Aliados sob seu comando não são movidos contra a vontade e nunca ficam Surpresos. Você lidera até 500 pessoas sem testes.}
\talento{Segundo Escalão}{3}{Gastando 2 PP, uma força aliada que estava a caminho chega agora --- patrulha, mercenários, guarda da cidade.}
\talento{Sem Baixas}{3}{Uma vez por Descanso Longo, quando um aliado chegaria a 0 PV, ele fica com 1 PV e se move 9m para fora do perigo.}
\talento{Ordem de Marcha}{3}{Aliados que te ouvem ganham +3m de Deslocamento e podem atravessar espaços ocupados por aliados livremente.}
\tecnica{Avante}{s}{4}{3}{1 Ações}{Voz}{Uma vez por combate: todos os aliados que te ouçam recebem 1 Ação adicional neste turno, que não pode conjurar magia de rank Santo ou superior.}
\patamar{Senhor da Guerra}{Dado de PV \dado{1d12+4} \textbullet\ +1 PP}
\maestria{A Guerra Já Acabou}{Escopo: o ano, uma guerra inteira. Uma vez por Descanso Longo, gastando 4 PP, declare uma condição estratégica já em vigor (estrada cortada, porto bloqueado, exército sem pagamento). Você pode declarar que um confronto planejado não vai acontecer. Recupere 2 PP em Descanso Curto.}
\talento{O Mapa É Meu}{3}{Seu Escopo passa a cobrir um continente inteiro: movimentação de tropas, rotas, colheitas, estações.}
\talento{Reputação de Aço}{3}{Exércitos que saibam que você comanda o outro lado sofrem penalidade de moral; alguns recusam engajamento.}
\talento{Herança de Comando}{3}{Escolha um talento de qualquer patamar desta árvore que não possua. Troque a cada Descanso Longo; nunca mais de um por vez.}
\tecnica{A Batalha Que Você Escolheu}{s}{4}{4}{3 Ações}{Campo de batalha}{Uma vez por Descanso Longo. Declare que este confronto foi montado por você e escolha duas: todos os aliados recebem 1 Ação extra nos próximos 3 turnos; os reforços inimigos não vêm; o terreno muda a seu favor; o comandante inimigo desconfia do segundo em comando. Exige inimigos organizados --- contra um monstro solitário, não faz nada.}

